% macro conflict, first definition kept: \DeclareMathOperator{\card}{card}
% macro conflict, first definition kept: \DeclareMathOperator{\Ob}{Ob}
% macro conflict, first definition kept: \newcommand{\Ind}{\operatorname{Ind}}
% macro conflict, first definition kept: \newcommand{\Ind}{\operatorname{Ind}}
% macro conflict, first definition kept: \newcommand{\cercle}[1]{\tikz[baseline=(c.base)]\node[draw,circle,inner sep=1pt](c){#1};}
% macro conflict, first definition kept: \newcommand{\Ind}{\operatorname{Ind}}
% macro conflict, first definition kept: \newcommand{\cardred}{\mathop{\textup{card\,réd}}}
% macro conflict, first definition kept: \newcommand{\defeq}{\overset{\textup{déf}}{=}}
% macro conflict, first definition kept: \newcommand{\cardred}{\mathop{\textup{card\,réd}}}
% macro conflict, first definition kept: \newcommand{\Ind}{\operatorname{Ind}}
% macro conflict, first definition kept: \newcommand{\circled}[1]{\tikz[baseline=(c.base),commutative diagrams/math mode=false]\node[circle,draw,inner sep=1pt](c){\upshape #1};}
% macro conflict, first definition kept: \newcommand{\Ob}{\mathrm{Ob}\,}
% macro conflict, first definition kept: \newcommand{\Fl}{\mathrm{Fl}\,}
% macro conflict, first definition kept: \newcommand{\cardred}{\operatorname{card\,r\acute{e}d}}
% macro conflict, first definition kept: \newcommand{\circled}[1]{\tikz[baseline=(c.base),commutative diagrams/math mode=false]\node[draw,circle,inner sep=0.8pt](c){\upshape #1};}
% macro conflict, first definition kept: \newcommand{\Ob}{\mathrm{Ob}\,}
% macro conflict, first definition kept: \newcommand{\Fl}{\mathrm{Fl}}
% macro conflict, first definition kept: \newcommand{\Hom}{\mathrm{Hom}}
% macro conflict, first definition kept: \newcommand{\Ind}{\operatorname{Ind}}
% macro conflict, first definition kept: \newcommand{\cardred}{\operatorname{card}\,\text{\upshape réd}}
% macro conflict, first definition kept: \newcommand{\defeq}{\overset{\text{\upshape déf}}{=}}
\documentclass[11pt]{article}
\usepackage{fontspec}
\usepackage[english,french]{babel}
\usepackage{amsmath}
\usepackage{amssymb}
\usepackage{amsthm}
\usepackage{mathrsfs}
\usepackage{tikz-cd}
\usepackage{enumitem}

\newcommand{\printedpage}[1]{\marginpar{\footnotesize p.~#1}}
\newcommand{\msp}[1]{\marginpar{\footnotesize [#1]}}
\newcommand{\circled}[1]{\tikz[baseline=(c.base),commutative diagrams/math mode=false]\node[draw,circle,inner sep=1pt](c){#1};}
\newcommand{\ovalbox}[1]{\tikz[baseline=(c.base)]\node[draw,rounded corners=6pt,inner sep=2pt](c){#1};}
\newcommand{\Ob}{\operatorname{Ob}}
\newcommand{\card}{\operatorname{card}}
\DeclareMathOperator{\Ind}{Ind}
\newcommand{\Ord}{\operatorname{Ord}}
\newcommand{\cardred}{\mathop{\textup{card\,réd}}\nolimits}
\newcommand{\cercle}[1]{\tikz[baseline=(c.base)]\node[draw,circle,inner sep=1pt](c){\textup{#1}};}
\newcommand{\Fl}{\operatorname{Fl}}
\newcommand{\Hom}{\operatorname{Hom}}
\newcommand{\id}{\mathrm{id}}
\newcommand{\limind}{\varinjlim}
\newcommand{\defeq}{\overset{\text{déf}}{=}}
\newcommand{\uHom}{\underline{\mathrm{Hom}}}
\newcommand{\End}{\operatorname{End}}
\newcommand{\edbr}[1]{\textup{\texttt{#1}}}
\newcommand{\uFl}{\underline{\mathrm{Fl}}}
\newcommand{\uIs}{\underline{\mathrm{Is}}}
\newcommand{\uMon}{\underline{\mathrm{Mon}}}
\newcommand{\uEp}{\underline{\textup{Ép}}}
\newcommand{\uGamma}{\underline{\Gamma}}
\newcommand{\Ens}{\mathrm{Ens}}
\newcommand{\pr}{\mathrm{pr}}
\newcommand{\diag}{\mathrm{diag}\,}
\newcommand{\cocart}{\mathrm{cocart}}
\newcommand{\cart}{\mathrm{cart}}
\newcommand{\cc}{\mathrm{c}}
\newcommand{\bDelta}{\boldsymbol{\Delta}}
\newcommand{\cE}{\mathcal{E}}
\newcommand{\cC}{\mathcal{C}}
\newcommand{\cM}{\mathcal{M}}

% Reader layout only; mathematical readings of the edition are retained.
\usepackage[a4paper,left=24mm,right=32mm,top=22mm,bottom=25mm,
  marginparwidth=25mm,marginparsep=4mm]{geometry}
\usepackage{microtype}
\setlength{\emergencystretch}{3em}
\setlength{\parskip}{3pt}
\hypersetup{pdftitle={Les Derivateurs: Chapter XVIII, pp. 316--429, with commentary},
 pdfauthor={Alexander Grothendieck},
 colorlinks=true,linkcolor=black,citecolor=black,urlcolor=blue}
\AtBeginDocument{%
 \begin{otherlanguage*}{english}
 \begin{titlepage}
 \vspace*{24mm}
 \noindent{\Large Alexander Grothendieck\par}
 \vspace{7mm}
 \noindent{\huge\bfseries Les D\'erivateurs\par}
 \vspace{5mm}
 \noindent{\Large Chapter XVIII, part d\par}
 \vspace{12mm}
 \noindent{\large Accessible conditions, categories of structures, and gluing\par}
 \vspace{12mm}
 \noindent Text with commentary\par
 \vspace{8mm}
 \noindent Manuscript pages 316--429\par
 \noindent Edition printed pages 119--163\par
 \vspace{10mm}
 \noindent\textbf{Reading the apparatus.}\quad The French text retains the edition's
 readings, including its slips. English commentary is set off by a grey rule.
 The separate claim and correction registers distinguish proof status,
 mathematical repairs, and documentary readings.\par
 \vfill
 \noindent Commentary and audit AI-produced by GPT-6 Astra (OpenAI, ChatGPT Pro mode),\par
 \noindent checked and integrated by Claude (Anthropic)\par
 \vspace{4mm}
 \noindent{\footnotesize Source: the electronic edition deciphered and transcribed
 by M.~K\"unzer, J.~Malgoire and G.~Maltsiniotis, version 17 December 2025.
 \par}
 \end{titlepage}
 \end{otherlanguage*}
}

\begin{document}
% ---- from der-ch18_p121-125.tex
\printedpage{119}
% head: A. Grothendieck : Dérivateurs, Ch. XVIII, Édition des Universités de Montpellier II et Paris VII
satisfaisant aux contraintes $(**)$, $(***)$, $(****)$. On va faire la construction en trois étapes.

\circled{1} Départ :
\[
i_0 \;=\; i\ .
\]
Comme $J'$ est cofinal dans $J$, existe $j_0\in\Ob J'$ et
\[
\varphi(i_0)\xrightarrow{u_0} j_0\ .
\]

\msp{page 316}

\circled{2} Supposons que les $i_\beta$, $j_\beta$ soient construits pour $\beta<\alpha$, ainsi que les $\tau_{\beta',\beta}:i_\beta\longrightarrow i_{\beta'}$, $\tau'_{\beta',\beta}:j_\beta\longrightarrow j_{\beta'}$ pour $\beta,\beta'<\alpha$, enfin les $u_\beta:\varphi(i_\beta)\longrightarrow j_\beta$, $v_\beta:j_\beta\longrightarrow\varphi(i_{\beta+1})$ pour $\beta<\alpha$ resp. pour $\beta+1<\alpha$, de façon à satisfaire aux contraintes $(**)$, $(****)$, i.e. commutativités dans
\[
\left\{
\begin{array}{l}
\left.
\begin{array}{ll}
\circled{1} &
\begin{tikzcd}[ampersand replacement=\&]
\varphi(i_\beta) \arrow[r, "u_\beta"] \arrow[rr, bend right=15, "\varphi(\tau_{\beta+1,\beta})"'] \& j_\beta \arrow[r, "v_\beta"] \& \varphi(i_{\beta+1})
\end{tikzcd}
\\[5ex]
\circled{2} &
\begin{tikzcd}[ampersand replacement=\&]
j_\beta \arrow[r, "v_\beta"] \arrow[rr, dotted, bend right=15, "\tau'_{\beta+1,\beta}"'] \& \varphi(i_{\beta+1}) \arrow[r, "u_{\beta+1}"] \& j_{\beta+1}
\end{tikzcd}
\end{array}
\right\}\quad \textit{si}\ \beta+1<\alpha
\\[8ex]
\begin{array}{ll}
\circled{3} &
\begin{tikzcd}[ampersand replacement=\&]
\varphi(i_\beta) \arrow[r, "u_\beta"] \arrow[d, "\varphi(\tau_{\beta',\beta})"'] \& j_\beta \arrow[d, "\tau'_{\beta',\beta}"] \\
\varphi(i_{\beta'}) \arrow[r, "u_{\beta'}"] \& j_{\beta'}
\end{tikzcd}
\qquad\qquad \text{si } \beta<\beta'<\alpha
\end{array}
\\[10ex]
\begin{array}{ll}
\circled{4} &
\begin{tikzcd}[ampersand replacement=\&]
j_\beta \arrow[r, "v_\beta"] \arrow[d, "\tau'_{\beta',\beta}"'] \& \varphi(i_{\beta+1}) \arrow[d, "\varphi(\tau_{\beta'+1,\beta+1})"] \\
j_{\beta'} \arrow[r, "v_{\beta'}"] \& \varphi(i_{\beta'+1})
\end{tikzcd}
\qquad\qquad \text{si } \beta<\beta' \text{ et } \beta'+1<\alpha\ ,
\end{array}
\end{array}
\right.
\]

\msp{page 317}

il faut construire
% foot: Version 17 décembre 2025

\printedpage{120}
% head: A. Grothendieck : Dérivateurs, Ch. XVIII, Édition des Universités de Montpellier II et Paris VII
\begin{itemize}[leftmargin=*]
\item[1)] $i_\alpha\in\Ob I$,\quad et un homomorphisme $(i_\beta)_{\beta\in I_\alpha}\longrightarrow i_\alpha$, i.e. des $\tau_{\alpha,\beta}$ : $i_\beta\longrightarrow i_\alpha$ $(\beta\in I_\alpha)$, compatibles aux translations $\tau_{\beta',\beta}$ pour $\beta<\beta'<\alpha$.
\item[2)] $j_\alpha\in\Ob J'$,\quad et un homomorphisme $(j_\beta)_{\beta\in I_\alpha}\longrightarrow j_\alpha$, i.e. des $\tau'_{\alpha,\beta}$ : $j_\beta\longrightarrow j_\alpha$ $(\beta\in I_\alpha)$, compatibles aux translations $\tau'_{\beta',\beta}$ pour $\beta<\beta'<\alpha$.
\item[3)] Une flèche
\[
u_\alpha:\varphi(i_\alpha)\longrightarrow j_\alpha\ .
\]
\item[4)] Si $\alpha$ n'est pas ordinal limite, i.e. $\alpha=\alpha_0+1$, une flèche
\[
v_{\alpha_0}:j_{\alpha_0}\longrightarrow\varphi(i_{\alpha_0+1})=\varphi(i_\alpha)\ ;
\]
\end{itemize}

et ceci de façon à satisfaire aux contraintes \circled{1}\,\circled{2}\,\circled{3}\,\circled{4} quand on y suppose $\beta+1\leqslant\alpha$ (au lieu de $\beta'<\alpha$) dans \circled{1} et \circled{2}, et $\beta'\leqslant\alpha$ (au lieu de $\beta'<\alpha$) dans \circled{3}, et $\beta'+1\leqslant\alpha$ (au lieu de $\beta'+1<\alpha$) dans \circled{4}. Donc pour les contraintes $1^\circ$, $2^\circ$, $4^\circ$, il n'y a en tenir compte \textit{que} si $\alpha$ n'est pas ordinal limite, donc $\alpha=\alpha_0+1$, et en y faisant % [sic] il n'y a en tenir -> il n'y a à en tenir, the preposition « à » appears to be missing
\[
\beta' \;=\; \alpha_0\ .
\]
Quant à la contrainte $3^\circ$, il faut en
\msp{page 318}
tenir compte pour $\beta'=\alpha$. Donc au total, il faut exiger
\[
(3_\alpha)(\beta)\qquad\qquad
\begin{tikzcd}
\varphi(i_\beta) \arrow[r, "u_\beta"] \arrow[d, "\varphi(\circled{$\tau_{\alpha,\beta}$})"'] & j_\beta \arrow[d, "\tau'_{\alpha,\beta}" {draw, circle, inner sep=1pt}] \\
\varphi(i_\alpha) \arrow[r, "u_\alpha" {draw, circle, inner sep=1pt}] & j_\beta % [sic] j_\beta -> j_\alpha, the target of \tau'_{\alpha,\beta} and u_\alpha is j_\alpha
\end{tikzcd}
\qquad \text{commutatif pour } \beta<\alpha\ (\footnotemark[140]),
\]
\footnotetext[140]{\textbf{N.B.} Il suffit de l'exiger pour $\beta=\alpha_0$, dans le cas $\alpha=\alpha_0+1$, cf. ci-dessous, p. 323.}
(par restriction sur $\alpha$), et si $\alpha=\alpha_0+1$ est ordinal non limite, il faut exiger de plus
\[
(4_\alpha)(\beta)\qquad\qquad
\begin{tikzcd}
j_\beta \arrow[r, "v_\beta"] \arrow[d, "\tau'_{\alpha_0,\beta}"'] & \varphi(i_{\beta+1}) \arrow[d, "\varphi(\circled{$\tau_{\alpha,\beta+1}$})"] \\
j_{\alpha_0} \arrow[r, "v_{\alpha_0}" {draw, circle, inner sep=1pt}] & \varphi(i_{\alpha_0+1})=\varphi(i_\alpha)
\end{tikzcd}
\qquad \beta<\alpha_0
\]
\[
(1_\alpha)\qquad\qquad
\begin{tikzcd}
\varphi(i_{\alpha_0}) \arrow[r, "u_{\alpha_0}"] \arrow[rr, bend right=15, "\varphi(\circled{$\tau_{\alpha,\alpha_0}$})"'] & j_{\alpha_0} \arrow[r, "v_{\alpha_0}" {draw, circle, inner sep=1pt}] & \varphi(i_{\alpha_0+1})=\varphi(i_\alpha)
\end{tikzcd}
\]
% foot: Version 17 décembre 2025

\printedpage{121}
% head: A. Grothendieck : Dérivateurs, Ch. XVIII, Édition des Universités de Montpellier II et Paris VII
\[
(2_\alpha)\qquad\qquad
\begin{tikzcd}
j_{\alpha_0} \arrow[r, "v_{\alpha_0}" {draw, circle, inner sep=1pt}] \arrow[rr, bend right=15, "\tau'_{\alpha,\alpha_0}"' {draw, circle, inner sep=1pt}] & \varphi(i_{\alpha_0+1})=\varphi(i_\alpha) \arrow[r, "u_\alpha" {draw, circle, inner sep=1pt}] & j_\alpha\ .
\end{tikzcd}
\]

Dans ces compatibilités, j'ai cerclé les flèches $\tau_{\alpha,\beta}$, $\tau'_{\alpha,\beta}$ $(\beta<\alpha)$, $\Bigl(\ \tau_{\alpha,\beta+1}\ (\beta+1<\alpha)\ \Bigr)$, $v_{\alpha_0}$, $\Bigl(\ \tau_{\alpha,\alpha_0}\ \Bigr)$, $u_\alpha$, $\Bigl(\ \tau'_{\alpha,\alpha_0}\ \Bigr)$, qui ont été introduites dans les donnés 1) 2) 3) 4), page 317. % [sic] donnés -> données, the noun « données » is meant

\msp{page 319}

On va distinguer donc les deux cas.

\textit{2 A)} $\alpha=\alpha_0+1$.

Pour construire $i_\alpha$ et les $\tau_{\alpha,\beta}:i_\beta\longrightarrow i_\alpha$, il suffit de construire $i_\alpha$ et
\[
\circled{$\tau_{\alpha,\alpha_0}$}\ :\ i_{\alpha_0}\ \longrightarrow\ i_\alpha=i_{\alpha_0+1}\ .
\]
De même pour $j_\alpha$ et les $\tau'_{\alpha,\beta}$, il suffit de construire
\[
\circled{$\tau'_{\alpha,\alpha_0}$}\ :\ j_{\alpha_0}\ \longrightarrow\ j_\alpha\ .
\]
Il faut de plus se donner
\[
j_{\alpha_0}\ \xrightarrow{\circled{$v_{\alpha_0}$}}\ \varphi(\underbrace{\,i_\alpha\,}_{=i_{\alpha_0+1}})\ \xrightarrow{\circled{$u_\alpha$}} j_\alpha\ ,
\]
donc en tout deux objets $i_\alpha$, $j_\alpha$ (dans $I$ et dans $J'$ respectivement) et quatre flèches $\tau_{\alpha,\alpha_0}$, $\tau'_{\alpha,\alpha_0}$, $v_{\alpha_0}$, $u_\alpha$. Il faut quatre compatibilités (page 318), où les deux dernières $3_\alpha(\beta)$, $4_\alpha(\beta)$ dépendent d'un paramètre $\beta<\alpha$, i.e. $\beta\leqslant\alpha_0$, ou de $\beta<\alpha_0$.

\msp{page 320 biffé}

\msp{page 321}

On va d'abord s'occuper de deux premières compatibilités $1_\alpha$, $2_\alpha$, en termes de la donnée de
\[
i_{\alpha_0}\ \in\ \Ob I\ ,\qquad j_{\alpha_0}\ \in\ \Ob J'\ ,\qquad u_{\alpha_0}\ :\ \varphi(i_{\alpha_0})\ \longrightarrow\ j_{\alpha_0}\ .
\]
Les deux compatibilités voulues pour les deux objets $i_\alpha$, $j_\alpha$ et les quatre flèches $\tau_{\alpha,\alpha_0}$, $\tau'_{\alpha,\alpha_0}$, $v_{\alpha_0}$, $u_\alpha$ sont indiquées dans le diagramme
\[
(*)\qquad\qquad
\begin{tikzcd}
\varphi(i_{\alpha_0}) \arrow[r, "u_{\alpha_0}"] \arrow[rr, bend right=15, "\varphi(\tau_{\alpha,\alpha_0})"'] & j_{\alpha_0} \arrow[r, "v_{\alpha_0}\ ?"] \arrow[rr, bend left=20, "\tau'_{\alpha,\alpha_0}"] & \varphi(\underbrace{\,i_\alpha\,}_{?}) \arrow[r, "u_\alpha\ ?"] & j_\alpha\ .
\end{tikzcd}
\]

On note que comme $\varphi$ est cofinal, $\exists\ i_\alpha=i\in\Ob I$ et $j_{\alpha_0}\xrightarrow{v_{\alpha_0}}\varphi(i)$. Mais on veut que le composé
\[
w\ :\ \varphi(i_{\alpha_0})\ \xrightarrow{u_{\alpha_0}}\ j_{\alpha_0}\ \xrightarrow{v_{\alpha_0}}\ \varphi(i)
\]
soit de la forme $\varphi(\tau)$, où $\tau=\tau_{\alpha,\alpha_0}:i_{\alpha_0}\longrightarrow i_\alpha$ est une flèche dans $I$. Comme $\varphi$ est un foncteur cofinal de catégories filtrantes, on sait
% foot: Version 17 décembre 2025
\printedpage{122}
% head: A. Grothendieck : Dérivateurs, Ch. XVIII, Édition des Universités de Montpellier II et Paris VII
\msp{page 322}
qu'on peut trouver $i\xrightarrow{\rho}i'$ dans $I$, tel que la flèche composée $\varphi(i_{\alpha_0})\xrightarrow{w}\varphi(i)\xrightarrow{\varphi(\rho)}\varphi(i')$ soit de la forme $\varphi(\tau)$. Donc quitte à remplacer $i$ par $i'$, et $v_{\alpha_0}$ par le composé $v'_{\alpha_0}$ : $j_{\alpha_0}\xrightarrow{v_{\alpha_0}}\varphi(i)\xrightarrow{\varphi(\rho)}\varphi(i')$, on peut supposer que $v_{\alpha_0}u_{\alpha_0}$ est de la forme $\varphi(\tau)$. Cela nous donne donc en principe l'objet $i=i_\alpha$, et les flèches $v_{\alpha_0}:j_{\alpha_0}\longrightarrow\varphi(i)$ et $\tau_{\alpha,\alpha_0}:i_{\alpha_0}\longrightarrow i_\alpha$, de façon à satisfaire à la première compatibilité dans $(*)$. Il reste à construire $j_\alpha$, $\tau'_{\alpha,\alpha_0}$, $u_\alpha$ de façon à satisfaire à la deuxième compatibilité $(*)$. Or $J'$ étant cofinal dans $J$, il existe $j_\alpha$ dans $J'$ et une flèche $u_\alpha:\varphi(i_\alpha)\longrightarrow j_\alpha$. D'autre part, comme $J'\longrightarrow J$ est pleinement fidèle, la flèche composée $j_{\alpha_0}\xrightarrow{v_{\alpha_0}}\varphi(i_\alpha)\xrightarrow{u_\alpha}j_\alpha$ dans $J$, entre objets de $J'$, provient d'une unique flèche $\tau'_{\alpha,\alpha_0}:j_{\alpha_0}\longrightarrow j_\alpha$ dans $J'$.

\msp{page 323}

Ainsi, on a construit les deux objets et les quatre flèches voulues, de façon à satisfaire aux contraintes $1_\alpha$, $2_\alpha$ de la page 318. Les contraintes $3_\alpha(\beta)$, $4_\alpha(\beta)$ sont-elles alors satisfaites automatiquement ?
\[
\begin{tikzcd}
\varphi(i_\beta) \arrow[r, "u_\beta"] \arrow[d] & j_\beta \arrow[d] \\
\varphi(i_{\alpha_0}) \arrow[r, "u_{\alpha_0}"] \arrow[d, "\varphi(\tau_{\alpha,\alpha_0})"'] & j_{\alpha_0} \arrow[d, "\tau'_{\alpha,\alpha_0}"] \\
\varphi(i_\alpha) \arrow[r, "u_\alpha"] & j_\alpha
\end{tikzcd}
\qquad\qquad \beta<\alpha, \text{ i.e. } \beta\leqslant\alpha_0\ .
\]
Le carré supérieur est commutatif par hypothèse de récurrence, pour que le carré total (composé) soit commutatif, il suffit que le carré inférieur le soit. On a donc une double-flèche $(\tau'_{\alpha,\alpha_0}u_{\alpha_0},u_\alpha\varphi(\tau_{\alpha,\alpha_0})):\varphi(i_\alpha)\rightrightarrows j_\alpha$, elle n'est peut-être pas triviale, mais on sait, $J'$ étant filtrante, qu'on peut l'égaliser par une flèche $j_\alpha\longrightarrow\bar{j}_\alpha$ dans $J'$. Donc quitte à remplacer $j_\alpha$ par % [sic] \varphi(i_\alpha) -> \varphi(i_{\alpha_0}), both composites of the double arrow start at \varphi(i_{\alpha_0})
\msp{page 324}
$\bar{j}_\alpha$, $u_\alpha$ par le composé $\bar{u}_\alpha:\varphi(i_0)\longrightarrow j_\alpha\longrightarrow\bar{j}_\alpha$ et $\tau'_{\alpha,\alpha_0}$ par le composé $\bar{\tau}'_{\alpha,\alpha_0}:j_{\alpha_0}\longrightarrow j_\alpha\longrightarrow\bar{j}_\alpha$ (sans toucher à $i_\alpha$, $v_{\alpha_0}$, $\tau_{\alpha_0,\alpha}$), on peut obtenir la commutativité, d'où $3_\alpha(\beta)$. % [sic] \varphi(i_0) -> \varphi(i_\alpha), u_\alpha has source \varphi(i_\alpha); [sic] \tau_{\alpha_0,\alpha} -> \tau_{\alpha,\alpha_0}, the arrow is written \tau_{\alpha,\alpha_0} everywhere else

Reste la contrainte $4_\alpha(\beta)$. Pour ceci, considérons le diagramme composite
\[
\begin{tikzcd}
j_\beta \arrow[r, "v_\beta"] \arrow[d, "\tau'_{\alpha_0,\beta}"'] & \varphi(i_{\beta+1}) \arrow[d, "\varphi(\tau_{\alpha,\beta+1})"] \\
j_{\alpha_0} \arrow[r, "v_{\alpha_0}"] & \varphi(i_{\alpha_0+1})=\varphi(i_\alpha)\ .
\end{tikzcd}
\]

Si $\alpha_0$ n'est pas un ordinal limite, mais de la forme $\alpha_1+1$, on peut encore se ramener au cas où $\beta=\alpha_1$, en procédant comme ci-dessus. Mais on peut procéder dans tous les cas, sans un tel artifice (\,!), en considérant les double-flèches
\[
(*)\qquad\qquad
\begin{tikzcd}
j_\beta \arrow[r, shift left, "\varphi(\tau_{\alpha,\beta+1})v_\beta"] \arrow[r, shift right, "v_{\alpha_0}\tau'_{\alpha_0,\beta}"'] & \varphi(i_\alpha)
\end{tikzcd}
\]
% foot: Version 17 décembre 2025
\printedpage{123}
% head: A. Grothendieck : Dérivateurs, Ch. XVIII, Édition des Universités de Montpellier II et Paris VII
déduites du carré précédent. Comme
\msp{page 325}
l'ensemble de ces double-flèches est de cardinal $\card\alpha<\pi$, et que $J$ est grande devant $\pi$ (il suffirait qu'il le soit pour les cardinaux $\pi'<\pi$), on peut trouver une flèche $\varphi(i_\alpha)\xrightarrow{u}j$ dans $J$ qui égalise toutes ces double-flèches. Comme $\varphi$ est cofinal, on peut supposer que $j$ est de la forme $\varphi(i)$. Comme $\varphi$ est un foncteur cofinal de catégories filtrantes, il existe alors $i\xrightarrow{\rho}i'$ dans $I$, tel que le composé $\varphi(i_\alpha)\xrightarrow{u}\varphi(i)\xrightarrow{\varphi(\rho)}\varphi(i')$ soit de la forme $\varphi(\sigma)$, pour $\sigma:i_\alpha\longrightarrow i'$. Au total, on a trouvé $\sigma:i_\alpha\longrightarrow i'$, tel que $\varphi(\sigma)$ $(=\varphi(\rho)u):\varphi(i_\alpha)\longrightarrow\varphi(i')$ égalise toutes les double-flèches envisagées dans $(*)$, p. 324. Donc il s'impose de remplacer $i_\alpha$ par $\bar{i}_\alpha=i'$, $v_{\alpha_0}$ par le composé $j_{\alpha_0}\longrightarrow\varphi(i_\alpha)\xrightarrow{\varphi(\sigma)}\varphi(\bar{i}_\alpha)$, $\tau_{\alpha,\alpha_0}$
\msp{page 326}
par le composé $\bar{\tau}_{\alpha,\alpha_0}:i_{\alpha_0}\xrightarrow{\tau_{\alpha,\alpha_0}}i_\alpha\xrightarrow{\sigma}\bar{i}_\alpha$.

Pour résumer, voici la recette pour construire, d'abord \ovalbox{$i_\alpha$, $v_{\alpha_0}$, $\tau_{\alpha,\alpha_0}$,} et \textit{ensuite} \ovalbox{$j_\alpha$, $u_\alpha$, $\tau'_{\alpha,\alpha_0}$}.

On procède d'abord comme page 321, 322 pour construire le premier triple, de façon à satisfaire à la première compatibilité $(*)$ (p. 321), i.e. à $1_\alpha$. On “corrige” ensuite ce choix de $i_\alpha$, $v_{\alpha_0}$, $\tau_{\alpha,\alpha_0}$, en remplaçant $i_\alpha$ par $\bar{i}_\alpha$ pour $\sigma:i_\alpha\longrightarrow\bar{i}_\alpha$, $v_{\alpha_0}$ par $\bar{v}_{\alpha_0}=\varphi(\sigma)v_{\alpha_0}$, et $\tau_{\alpha,\alpha_0}$ par $\bar{\tau}_{\alpha,\alpha_0}=\sigma\tau_{\alpha,\alpha_0}$, de sorte à obtenir les compatibilités $4_\alpha(\beta)$ pour $\beta<\alpha_0$, comme on vient de l'indiquer. On ne touchera plus à ces choix de $i_\alpha$, $v_{\alpha_0}$, $\tau_{\alpha,\alpha_0}$, dans la suite de la construction. On construit $(j_\alpha,u_\alpha,\tau'_{\alpha,\alpha_0})$ de façon à satisfaire à la deuxième compatibilité $(*)$ (p. 321), comme
\msp{page 327}
indiqué page 322. Ensuite, on corrige ce choix, en remplaçant $j_\alpha$ par $\bar{j}_\alpha$, avec $\rho:j_\alpha\longrightarrow\bar{j}_\alpha$, et $u_\alpha$ par $\bar{u}_\alpha=\rho u_\alpha$, $\tau'_{\alpha,\alpha_0}$ par $\bar{\tau}'_{\alpha,\alpha_0}=\rho\tau'_{\alpha,\alpha_0}$, de façon à satisfaire aux contraintes $3_\alpha(\beta)$ de la page 318.

Cela achève la construction dans le cas où $\alpha=\alpha_0+1$.

\textit{Cas 2 B : $\alpha$ est ordinal limite.}

Il faut donc construire les données 1) 2) 3) de la page 317 (à l'exclusion de la donnée 4), de façon à satisfaire à la seule contrainte $3_\alpha(\beta)$.

On commence à construire $i_\alpha$ et les $\tau_{\beta,\alpha}$ dans $I$. C'est possible, car $I_\alpha=\{\beta\mid\beta<\alpha\}$ est de cardinal $\pi_\alpha<\pi$ – il suffit donc de supposer $I$ grande devant les cardinaux $\pi'<\pi$.

On construit ensuite de même $J_\alpha$ dans $J'$ et les $\tau'_{\beta,\alpha}$ $(\beta<\alpha)$, en utilisant le % [sic] J_\alpha -> j_\alpha, the object of J' is meant
\msp{page 328}
fait que $J'$ est grande devant $\card\alpha=\card I_\alpha$. On doit de plus construire une flèche
\[
u_\alpha\ :\ \varphi(i_\alpha)\ \longrightarrow\ j_\alpha\ .
\]
C'est possible, car $J'$ étant cofinale dans $J$, il existe $j'$ dans $J'$ et $\varphi(i_\alpha)\longrightarrow J'$, et $J'$ étant filtrante, il existe $\bar{j}_\alpha$ et des flèches % [sic] \varphi(i_\alpha)\longrightarrow J' -> \varphi(i_\alpha)\longrightarrow j', the arrow goes to the object j'
$\begin{tikzcd}[ampersand replacement=\&, baseline=(\tikzcdmatrixname-1-1.base)]
j_\alpha \arrow[rr] \& \& j_\alpha \\
\& j' \arrow[ur] \&
\end{tikzcd}$, donc quitte à remplacer $j_\alpha$ par $\bar{j}_\alpha$, les $\tau'_{\alpha,\beta}$ par des $\bar{\tau}'_{\alpha,\beta}$, on trouve bien une flèche $u_\alpha:\varphi(i_\alpha)\longrightarrow j_\alpha$. Il faut ensuite s'occuper des % [sic] target j_\alpha in the small diagram -> \bar{j}_\alpha, the two arrows go to the object \bar{j}_\alpha just introduced
% foot: Version 17 décembre 2025

% ---- from der-ch18_p126-130.tex
\printedpage{124}
% head: A. Grothendieck : Dérivateurs, Ch. XVIII, Édition des Universités de Montpellier II et Paris VII

compatibilités $3_\alpha(\beta)$, donc des double-flèches
\[
\begin{tikzcd}
\varphi(i_\beta) \arrow[r, shift left, "\tau'_{\alpha,\beta}u_\beta"] \arrow[r, shift right, "u_\alpha\varphi(\tau_{\alpha,\beta})"'] & j_\alpha \; .
\end{tikzcd}
\]
Comme l'ensemble de ces double-flèches est de cardinal $\card\alpha < \pi$, et que $J'$ est grande devant ledit cardinal, on
\msp{page 329}
trouve donc une flèche
\[
\rho : j_\alpha \longrightarrow j
\]
dans $J$ qui les égalise toutes. Comme $J'$ est cofinal, on peut supposer $j$ dans $J'$, soit $\bar{j}_\alpha$. Donc quitte à remplacer $j_\alpha$ par $\bar{j}_\alpha$, et les $\tau'_{\alpha,\beta}$ par $\bar{\tau}'_{\alpha,\beta} = \rho\tau'_{\alpha,\beta}$, enfin $u_\alpha : \varphi(i_\alpha) \longrightarrow j_\alpha$ par $\bar{u}_\alpha = \rho u_\alpha : \varphi(i_\alpha) \longrightarrow \bar{j}_\alpha$, on satisfait à toutes les contraintes voulues.

Je résume le résultat (``$\pm$ tautologique''!) prouvé dans les pages 313-329.

% a vertical change bar is printed in the left margin alongside the statement from "Lemme-clef" to "suivantes :"
\textbf{Lemme-clef 4.13.17.} \textit{Considérons un carré cartésien dans} Cat
\[
\begin{tikzcd}
I \arrow[r, "\varphi"] & J \\
I' \arrow[u, hook] \arrow[r] & J' \arrow[u, hook, "\psi"'] \; .
\end{tikzcd}
\]

\textit{On fait les hypothèses suivantes :}

\msp{page 330}

\begin{enumerate}[label={}]
\item[\textit{1)}] \textit{$\varphi$ et $\psi$ sont cofinaux, $\psi$ est l'inclusion d'une sous-catégorie strictement pleine $J'$ de $J$, $I$ est filtrant (donc $J$ et $J'$ aussi).}
\end{enumerate}

\textit{On cherche des conditions pour que l'inclusion de la sous-catégorie strictement pleine $I' \longrightarrow I$ soit cofinale (auquel cas on voit tout de suite que $I'$ est filtrante, et $I' \longrightarrow J'$ est aussi cofinal). Pour ceci, on suppose qu'il existe un \textup{ordinal limite} $\alpha_0$ ayant les propriétés supplémentaires suivantes :}

\begin{enumerate}[label={}]
\item[\textit{2)}] \textit{$I$ et $J'$ sont stables par $\varinjlim$ de type $I_{\alpha_0}$ $(= \{\beta \mid \beta \ \textit{ordinal}\ < \alpha_0\})$, et les foncteurs $\varphi$, $\psi$ commutent à ces limites.}
\item[\textit{3)}] \textit{Pour tout $\alpha \in I_{\alpha_0}$, $I$ est grande devant $\card\alpha$. (Donc, $\varphi$ étant cofinal, $J$ est également grande devant $\card\alpha$, et de même, $\psi$ étant pleinement fidèle, et cofinal, $J'$ est grande devant $\card\alpha$.)}
\end{enumerate}

\textit{Sous ces conditions, $I'$ est cofinal dans $I$ (donc aussi grande devant $\card\alpha$, pour tout $\alpha \in I_{\alpha_0}$).}

La démonstration est celle donnée ci-dessus. On n'a pas eu à utiliser que l'ordinal $\alpha_0 = \alpha_\pi$ était de cette espèce particulière, seulement que c'était un cardinal limite.

\printedpage{125}
% head: A. Grothendieck : Dérivateurs, Ch. XVIII, Édition des Universités de Montpellier II et Paris VII

\msp{page 331}

Ce lemme implique, comme on l'a vu à la page 312, l'implication b $\Longrightarrow$ d dans le diagramme d'implications $(*)$ page 310, lequel est donc à présent prouvé.

Pour achever de prouver le théorème 4.13.15, il reste à prouver les implications
\[
\begin{tikzcd}[row sep=small]
\mathrm{c}' \arrow[dr, Rightarrow] & & \\
& \mathrm{a}' \arrow[r, Rightarrow] & \mathrm{a} \; . \\
\mathrm{b} \arrow[ur, Rightarrow] & &
\end{tikzcd}
\]

Démontrons donc (b ou c$'$) $\Longrightarrow$ a$'$. Il reste à prouver que $\underline{S}$ est $\pi$-accessible dans $\mathcal{M}$, i.e. que le foncteur pleinement fidèle
\[
\Ind_\pi(\underline{S}(\pi)) \longrightarrow \mathcal{M} \tag*{$(*)$}
\]
a pour image essentielle $\underline{S}$. Donc il faut prouver que si $x \in \Ob\mathcal{M}$, alors
\[
x \in \Ob\underline{S} \quad \overset{(?)}{\Longleftrightarrow} \quad
\begin{minipage}[c]{0.55\textwidth}
$x$ est dans l'image essentielle de $\Ind_\pi(\underline{S}(\pi))$, i.e. $x$ est isomorphe à $\varinjlim\limits_{i\in I} x_i$, les $x_i$ dans $\Ob\underline{S}(\pi)$, $I$ grande devant $\pi$.
\end{minipage}
\tag*{$(*)$}
\]

\msp{page 332}

$\Bigg($ Supposons d'abord que $\underline{S}$ satisfasse b) (\footnote[141]{Maladroit, on peut traiter les deux cas b) et c) en même temps, cf. bas de page 333.}). Donc on sait que
\[
x \in \Ob\underline{S} \qquad \Longleftrightarrow \qquad \underline{S}/x \text{ est cofinale dans } \mathcal{M}_\pi/x \; . \tag*{$(**)$}
\]

Nous voulons déduire l'équivalence $(*)$ de l'équivalence $(**)$.

Prouvons $\Longrightarrow$ dans $(*)$, i.e. supposons $x \in \Ob\underline{S}$, et prouvons que $x$ est dans l'image essentielle de $\Ind_\pi\underline{S}(\pi)$. En vertu de $(**)$, $\underline{S}(\pi)/x$ est une sous-catégorie (strictement pleine) cofinale dans $\mathcal{M}_\pi/x$, donc comme $\mathcal{M}_\pi/x$ est grande devant $\pi$, $\underline{S}(\pi)/x$ l'est aussi. Comme on sait que
\[
x \ = \ \varinjlim_{i\in\mathcal{M}_\pi/x} x_i
\]
($\pi$ étant cardinal utile pour $\mathcal{M}$), il s'ensuit par cofinalité qu'on a aussi
\[
x \ = \ \varinjlim_{i\in\underline{S}(\pi)/x} x_i \; ,
\]
et comme $\underline{S}(\pi)/x$ est grande devant $\pi$, cela implique que $x$ appartient à l'image essentielle de $\Ind_\pi(\underline{S}(\pi))$.

Prouvons $\Longleftarrow$ dans $(*)$, i.e. supposons que $x$ soit dans l'image essentielle de $\Ind_\pi(\underline{S}(\pi))$,
\printedpage{126}
% head: A. Grothendieck : Dérivateurs, Ch. XVIII, Édition des Universités de Montpellier II et Paris VII
\msp{page 333}
\[
x \ = \ \varinjlim_{i\in I} x_i \; , \qquad\quad \text{les } x_i \text{ dans } \underline{S}(\pi),\ I \text{ grande devant } \pi.
\]

Alors le foncteur composé $i \longmapsto x_i$,
\[
I \ \longrightarrow \ \underline{S}(\pi)/x \ \hookrightarrow \ \mathcal{M}_\pi/x
\]
est cofinal (cf. 4.4.2, page 79), donc $\underline{S}(\pi)/x \longrightarrow \mathcal{M}_\pi/x$ satisfait la condition F 1 (puisque le composé le satisfait), et comme $\underline{S}(\pi)/x \longrightarrow \mathcal{M}_\pi/x$ est pleinement fidèle, cela implique déjà que ce foncteur est cofinal. Donc par $(**)$ on en conclut $x \in \Ob\underline{S}$, q.e.d.

Supposons maintenant que $S$ satisfasse c$'$ (au lieu de b), et prouvons encore l'équivalence $(*)$. Pour ceci, écrivons
\[
\square \qquad x \ = \ \varinjlim_{i\in I} x_i \qquad\quad
\left\{
\begin{minipage}[c]{0.5\textwidth}
les $x_i \in \Ob\mathcal{M}_\pi$, $I$ \textit{ensemble ordonné}, $(x_i)_{i\in I}$ est $\pi_0$-$\pi$-adaptée (\footnotemark[142]).
\end{minipage}
\right.
\]
\footnotetext[142]{\textbf{N.B.} Ça existe, grâce à 4.13.14.}

ce qui est possible puisque $\pi$ est utile pour $\mathcal{M}$ $\Bigg)$. Définissons $I_{\underline{S}}$ comme dans l'énoncé de c), donc comme image inverse de $\underline{S}(\pi)/x \hookrightarrow \mathcal{M}_\pi/x$ par $I \longrightarrow \mathcal{M}_\pi$. Par hypothèse d$'$) \texttt{[plutôt c$'$)]}, on a
\[
x \in \Ob\underline{S} \qquad \Longleftrightarrow \qquad I_{\underline{S}} \text{ est cofinal dans } I \; . \tag*{$(**')$}
\]

En utilisant $(**')$, prouvons encore $(*)$ (page 331). En fait, on utilisera seulement qu'on a
\msp{page 334}
l'isomorphisme $\square$, avec $I$ grande devant $\pi$, et qu'on a l'équivalence $(**')$, sans avoir à utiliser que $I$ est ordonnée. Donc la démonstration vaut aussi dans le cas c).

Prouvons $\Longrightarrow$ dans $(*)$, i.e. supposons $x \in \Ob\underline{S}$, prouvons que $x$ est dans l'image essentielle de $\Ind_\pi(\underline{S}(\pi))$. En vertu de $(**')$, $I_{\underline{S}}$ est cofinale dans $I$, et comme $I_{\underline{S}} \longrightarrow I$ est une inclusion pleinement fidèle, il s'ensuit que $I_{\underline{S}}$ est également grande devant $\pi$, d'autre part on aura par cofinalité
\[
x \ = \ \varinjlim_{i\in I_{\underline{S}}} x_i \; ,
\]
où cette fois les $x_i$ sont dans $\underline{S}(\pi) = \underline{S}\cap\mathcal{M}_\pi$, d'où résulte que $x$ est dans l'image essentielle de $\Ind_\pi(\underline{S}(\pi))$.

Prouvons $\Longleftarrow$ dans $(*)$, i.e. supposons que l'on ait
\[
x \ = \ \varinjlim_{j\in J} x_j \; , \qquad\quad \text{les } x_j \text{ dans } \underline{S}(\pi),\ J \text{ grande devant } \pi \; .
\]

\printedpage{127}
% head: A. Grothendieck : Dérivateurs, Ch. XVIII, Édition des Universités de Montpellier II et Paris VII

\msp{page 335}

En invoquant à nouveau 4.4.2, on trouve que $J \longrightarrow \mathcal{M}_\pi/x$ est cofinal, et comme ce foncteur se factorise par le foncteur d'inclusion pleinement fidèle $\underline{S}(\pi)/x \hookrightarrow \mathcal{M}_\pi/x$, il s'ensuit comme tantôt que ce dernier est cofinal. Nous allons admettre pour un moment que ceci implique que l'image inverse $I_{\underline{S}}$ de cette sous-catégorie strictement pleine, par le foncteur également cofinal $I \longrightarrow \mathcal{M}_\pi/x$, est cofinal dans $I$. Ceci étant, l'équivalence $(**')$ (p. 333) implique $x \in \Ob\underline{S}$, ce qu'on voulait.

Donc il reste à prouver que si
\[
J' \ = \ \underline{S}(\pi)/x \hookrightarrow \mathcal{M}_\pi/x \ = \ J
\]
est cofinal, alors $I' = I_{\underline{S}} \longrightarrow I$ l'est aussi. Il semble qu'il faille bel et bien à nouveau invoquer les conditions de stabilité $\mathrm{L}(\pi_0,\pi)$ faites dans d$'$) \texttt{[plutôt c$'$)]} ce qui n'avait pas été nécessaire pour le cas c) et
\msp{page 336}
le lemme-clef (un peu) vache 4.13.17, où on prendra (comme dans la démonstration écrite
avant l'énoncé de cet lemme) % [sic] cet lemme -> ce lemme, demonstrative adjective "cet" before a consonant
$\alpha = \alpha_\pi \overset{\text{déf}}{=}$ plus petit ordinal $\alpha$ tel que $\card\alpha = \pi$. Donc $I_\alpha$ est de cardinal $\pi$, et grande devant tout cardinal $\pi' < \pi$, et en particulier devant $\pi'_0$. L'hypothèse $2^\circ$) de 4.13.17 est vérifiée pour $I$, par définition d'un système inductif $\pi_0$-$\pi$-adapté, et pour $J'$, grâce à l'hypothèse $\mathrm{L}(\pi_0,\pi)$ sur $\underline{S}(\pi)$. De même, l'hypothèse $3^\circ$) ($I$ grande devant tout cardinal $< \pi$) est satisfaite par définition des systèmes $\pi_0$-$\pi$-adaptés -- $I$ est même grande devant $\pi$. L'hypothèse $1^\circ$) a déjà été vue.

On a donc prouvé (par voie circulaire)
\[
\begin{tikzcd}[row sep=small]
& & & \mathrm{c}' & \\
\mathrm{a} \arrow[r, Rightarrow] & \Bigg(\mathrm{a}' \arrow[r, Leftrightarrow] & \mathrm{c} \arrow[ur, Leftrightarrow] \arrow[dr, Rightarrow] & & \Bigg)\, , \\ % Verified against PDF leaf 9: the oblique equivalence has both arrowheads.
& & & \mathrm{b} &
\end{tikzcd}
\]

\texttt{[et aussi b $\Longrightarrow$ c,]} en ayant à invoquer le lemme-clef vache pour les \textit{deux} implications
\[
\mathrm{a}' \ \Longrightarrow \ \mathrm{c} \qquad\quad \text{et} \qquad\quad \mathrm{c}' \ \Longrightarrow \ \mathrm{a}'
\]
(donc pour deux implications \textit{en sens inverses}, chose étrange \dots).

\msp{page 337}

Il reste seulement encore à prouver l'implication
\[
\mathrm{a}' \ \Longrightarrow \ \mathrm{a} \; ,
\]
ce qui est essentiellement sorital.

Cela achève la démonstration du critère d'accessibilité 4.13.15 d'une sous-catégorie pleine $\underline{S}$ de $\mathcal{M}$.

On va donner une variante un peu décantée de 4.13.15, qui (avec des hypothèses un peu plus faibles) capte l'essentiel de sa substance. L'énoncé 4.13.15 en est une conséquence quasi immédiate :

\printedpage{128}
% head: A. Grothendieck : Dérivateurs, Ch. XVIII, Édition des Universités de Montpellier II et Paris VII

\msp{page 338}

\textbf{Corollaire 4.13.18} (de 4.13.15). \textit{Soient $\mathcal{M}$ une catégorie accessible, $\pi$ un cardinal utile pour $\mathcal{M}$, $\alpha_0$ un \textup{ordinal limite} tel que l'on ait $\card\alpha_0 \leqslant \pi$. Soit $I_{\alpha_0} = \{\alpha \mid \alpha < \alpha_0\}$, et supposons que $\mathcal{M}$ soit stable par $\varinjlim$ de type $I_{\alpha_0}$.}

\textit{Soit alors $\underline{S}$ une sous-catégorie strictement pleine de $\mathcal{M}$, $\underline{S}(\pi) = \underline{S}\cap\mathcal{M}_\pi$. Désignons par $L(\alpha_0)$ la condition : $\underline{S}(\pi)$ est stable dans $\mathcal{M}_\pi$ (ou ce qui revient au même, dans $\mathcal{M}$) par $\varinjlim$ de type $I_{\alpha_0}$. Alors les conditions suivantes sont équivalentes.}

\begin{enumerate}[label={}]
\item[\textit{a)}] \textit{$\underline{S}$ est $\pi$-accessible dans $\mathcal{M}$, et est stable dans $\mathcal{M}$ par $\varinjlim$ de type $I_{\alpha_0}$.}
\item[\textit{a$'$)}] \textit{$\underline{S}$ est $\pi$-accessible et $\underline{S}(\pi)$ satisfait $L(\alpha_0)$.}
\item[\textit{b)}] \textit{$\underline{S}(\pi)$ satisfait $L(\alpha_0)$, et pour tout $x \in \Ob\mathcal{M}$, on a ceci : $x \in \Ob\underline{S}$ si et seulement si la sous-catégorie strictement pleine $\underline{S}(\pi)/x$ de $\mathcal{M}(\pi)/x$ est cofinale dans $\mathcal{M}(\pi)/x$} \texttt{[lire $\mathcal{M}_\pi/x$]}.
\item[\textit{c)}] \textit{$\underline{S}(\pi)$ satisfait $L(\alpha_0)$, de plus pour tout $x \in \Ob\mathcal{M}$, et toute représentation}
\[
x \ \simeq \ \varinjlim_{i\in I} x_i
\]
\msp{page 339}
\textit{qui est ``$\alpha_0$-$\pi$-adaptée'', i.e. telle que}
\begin{enumerate}[label={}]
\item[\textit{$1^\circ$)}] \textit{$I$ équivalente à une petite catégorie, et grande devant $\pi$,}
\item[\textit{$2^\circ$)}] \textit{les $\varinjlim$ de type $I_{\alpha_0}$ existent dans $I$,}
\item[\textit{$3^\circ$)}] \textit{le foncteur $I \longrightarrow \mathcal{M}$, $i \longmapsto x_i$, commute aux $\varinjlim$ de type $I_{\alpha_0}$, et enfin}
\item[\textit{$4^\circ$)}] \textit{les $x_i \in \Ob\mathcal{M}_\pi$}
\end{enumerate}
\textit{(comparer 4.13.11), on a ceci (où on désigne par $I_{\underline{S}}$ la sous-catégorie strictement pleine de $I$ formée des $i$ tels que $x_i \in \Ob\underline{S}$) : on a $x \in \Ob\underline{S}$ si et seulement si $I_{\underline{S}}$ est cofinale dans $I$.}
\item[\textit{c$'$)}] \textit{(\textup{Si} $I_{\alpha_0}$ est grande devant $\pi_0$, et $\pi^{\pi_0} = \pi$, pour pouvoir invoquer 4.13.14, cf. p. 335.) Comme c), mais en supposant $I$ ordonnée.}
\end{enumerate}
(\footnote[143]{\textbf{N.B.} Si on ne suppose par $I_{\alpha_0}$ grande devant $\pi_0$ et % [sic] suppose par -> suppose pas, negation "ne ... pas"
$\pi^{\pi_0} = \pi$, alors toutes les conditions de a) à c) (a, a$'$, b, c) sont équivalentes et \textit{impliquent} c, mais ne lui sont peut-être pas équivalentes.}).

\textsc{Démonstration}. Elle se fait suivant le même diagramme logique que celle de 4.13.15,
\[
\begin{tikzcd}[row sep=small]
& & & \mathrm{c}' \arrow[dr, Rightarrow, "\textbf{!}"] & \nearrow\ \begin{tabular}[b]{@{}r@{}}\footnotesize ici, il faut supposer\\ \footnotesize $I_{\alpha_0}$ grande devant $\pi_0$\\ \footnotesize et $\pi^{\pi_0} = \pi$.\end{tabular} & \\
\mathrm{a} \arrow[r, Rightarrow] & \mathrm{a}' \arrow[r, Rightarrow, "\textbf{!}"] & \mathrm{c} \arrow[ur, Rightarrow] \arrow[dr, Rightarrow] & & \mathrm{a}' \arrow[r, Rightarrow] & \mathrm{a} \; , \\
& & & \mathrm{b} \arrow[ur, Rightarrow] & &
\end{tikzcd}
\]
où on a marqué du signe\,! les deux implications qui sont ``vaches'', en ce sens qu'elles font appel au lemme-clef 4.13.17 (pages 329-330), qui se prouve par une construction transfinie un peu vissée. La démonstration des sept implications du diagramme est essentiellement
celle donné dans les pages précédentes 311-337, % [sic] celle donné -> celle donnée, agreement with feminine "celle"

% ---- from der-ch18_p131-135.tex
\setcounter{footnote}{143}

\printedpage{129}
% head: A. Grothendieck : Dérivateurs, Ch. XVIII, Édition des Universités de Montpellier II et Paris VII
% foot: Version 17 décembre 2025

\msp{page 340}
à l'exception de a$'$ $\Longrightarrow$ a, qui est un peu plus simple dans le cas actuel. On va la dégager sous forme de lemme un peu plus général. (\textbf{N.B.} La démonstration de a$'$ $\Longrightarrow$ a dans 4.13.15 avait été omise comme “soritale”\,!, cf. page 337\,\dots)

\textbf{Lemme 4.13.19.} \textit{Soient $\mathcal{M}$ une catégorie accessible dans $\mathcal{M}$, $\pi$ un cardinal utile pour $\mathcal{M}$ et tel que $\underline{S}$ soit $\pi$-accessible, i.e. $\underline{S}$ est l'image essentielle de $\Ind_\pi(\underline{S}(\pi))$ par le foncteur pleinement fidèle canonique $\Ind_\pi(\underline{S}(\pi))\longrightarrow\mathcal{M}$, où $\underline{S}(\pi)=\underline{S}\cap\mathcal{M}_\pi$. On suppose qu'il existe un ordinal limite $\alpha$, tel que $\card\alpha\leqslant\pi$ et tel que $\mathcal{M}_\pi$ et $\underline{S}(\pi)$ $(=\underline{S}_\pi)$ soient stables par $\varinjlim$ de type $I_\alpha$ soient stables par $\varinjlim$ de type $I_\alpha$, et que le foncteur d'inclusion $\underline{S}(\pi)\longrightarrow\mathcal{M}_\pi$ y commute. Soit $\Phi$ un ensemble de catégories telles que $\forall\ I\in\Phi$, on ait $\cardred I\leqslant\pi$, et supposons que $\mathcal{M}_\pi$ soit stable par $\varinjlim$ de type $\Phi$, et que $\underline{S}(\pi)$ soit stable dans $\mathcal{M}_\pi$ par lesdites limites (ou ce qui revient au même, qu'il soit stable en lui-même par ces $\varinjlim$\,, et que le foncteur d'inclusion y commute). Alors $\mathcal{M}$ et $\underline{S}$ sont stables par $\varinjlim$ de type $\Phi$, et le foncteur d'inclusion $\underline{S}\longrightarrow\mathcal{M}$ y commute, i.e. $\mathcal{M}$ est stable par ces limites, et $\underline{S}$ stable dans $\mathcal{M}$ par ces mêmes limites.} % [sic] « soient stables par lim de type I_alpha » printed twice -> once, repetition

\msp{page 341}
\textsc{Démonstration}. On applique 4.11.1 (p. 267-268) séparément à $\mathcal{M}\simeq\Ind_\pi\mathcal{M}_\pi$ et à $\underline{S}\simeq\Ind_\pi\underline{S}(\pi)$, pour conclure que $\mathcal{M}$ et $\underline{S}$ sont stables séparément par $\varinjlim$ de type $\Phi$. Il reste à voir que le foncteur $\underline{S}\longrightarrow\mathcal{M}$ y commute. Cela se voit comme dans la démonstration de 4.11.1, on obtient plus généralement ceci :

\textbf{Lemme 4.13.20} (cor. à 4.11.1). \textit{Soient $f_0 : C\longrightarrow C'$ un foncteur entre catégories équivalentes à des petites catégories, $\pi$ un cardinal infini, $f:\mathcal{M}\stackrel{\textup{déf}}{=}\Ind_\pi C\longrightarrow\mathcal{M}'\stackrel{\textup{déf}}{=}\Ind_\pi C'$ l'extension canonique $\Ind_\pi(f_0)$ de $f_0$. Soient $\Phi$ et $\alpha$ comme dans 4.11.1, et supposons que $I_\alpha\in\Phi$ et que $f_0$ commute aux $\varinjlim$ de type $\Phi$. Alors $f$ commute à ces mêmes limites inductives.} (\footnote{Il vaut mieux prendre des foncteurs $f_0 : C\longrightarrow\mathcal{N}$, et $\mathcal{N}$ [?] que $f_0$ commute aux $\varinjlim$ de type $\Phi$, donc il en est de même de son extension $\Ind_\pi(C,\mathcal{N})\longrightarrow\mathcal{N}$. (\textbf{N.B.} On suppose que $\mathcal{N}$ satisfait $\mathrm{L}_\pi$.)}).

\textbf{Corollaire 4.13.21.} \textit{Sous les conditions de 4.13.20, supposons que $\Phi$ contienne les catégories de cardinal $\leqslant\pi$ (resp. les catégories filtrantes de cardinal $\leqslant\pi$, resp. les catégories filtrantes grandes devant $\pi_0$ de cardinal $\leqslant\pi$, resp. les catégories discrètes de cardinal $\leqslant\pi$). Alors $f$ commute aux petites $\varinjlim$ (resp. aux petites $\varinjlim$ filtrantes, resp. aux petites $\varinjlim$ filtrantes grandes devant $\pi_0$, resp. aux petites sommes). (Réciproque claire et “laissée au lecteur”.)}

Se voit comme 4.4.14 et 4.4.14.1 (p. 97, 98).

\msp{page 342}
Voici enfin, il me semble, la forme définitive (\,?) de 4.13.15, 4.13.18.

\textbf{Théorème 4.13.22} (\footnote{cf. variante 4.13.36, page 358.}). \textit{Soient $\mathcal{M}$ une catégorie accessible, $\pi$ un cardinal infini, utile pour $\mathcal{M}$, et $\Phi$ un ensemble de catégories satisfaisant les deux conditions :}
\begin{itemize}[leftmargin=2.5em]
\item[\textit{1°)}] $J\in\Phi$ \textit{implique} $\cardred J\leqslant\pi$.
\item[\textit{2°)}] \textit{Il existe un} ordinal limite $\alpha_0$ \textit{avec} $I_{\alpha_0}\in\Phi$.
\end{itemize}
\textit{On suppose $\mathcal{M}$ stable par $\varinjlim$ de type $\Phi$. Soient $\underline{S}$ une sous-catégorie} strictement pleine \textit{de $\mathcal{M}$, et $\underline{S}(\pi)=\underline{S}\cap\mathcal{M}_\pi$. Désignons par $L(\varphi)$ la condition sur $\underline{S}(\pi)$ : $\underline{S}(\pi)$ est stable dans} % [sic] L(\varphi) -> L(\Phi), the condition is later written L(\Phi)
\printedpage{130}
% head: A. Grothendieck : Dérivateurs, Ch. XVIII, Édition des Universités de Montpellier II et Paris VII
% foot: Version 17 décembre 2025
\textit{$\mathcal{M}$ (où ce qui revient au même, dans $\mathcal{M}_\pi$) par $\varinjlim$ de type $\Phi$. Ceci posé, les conditions suivantes sur $\underline{S}$ sont équivalentes :} % [sic] où ce qui revient -> ou ce qui revient, conjunction "ou"
\begin{itemize}[leftmargin=2.5em]
\item[\cercle{a}] \textit{$\underline{S}$ est $\pi$-accessible dans $\mathcal{M}$, et stable dans $\mathcal{M}$ par $\varinjlim$ de type $\Phi$.}
\item[a$'$)] \textit{$\underline{S}$ est $\pi$-accessible dans $\mathcal{M}$, et $\underline{S}(\pi)$ satisfait à $L(\Phi)$.}
\end{itemize}
\msp{page 343}
\begin{itemize}[leftmargin=2.5em]
\item[b)] \textit{$\underline{S}(\pi)$ satisfait à $L(\Phi)$, et pour tout $x\in\Ob\mathcal{M}$, on a ceci : $x\in\Ob\underline{S}$ si et seulement si la sous-catégorie strictement pleine $\underline{S}(\pi)/x$ de $\mathcal{M}(\pi)/x$ est cofinale dans $\mathcal{M}(\pi)/x$} \texttt{[lire} $\mathcal{M}_\pi/x$\texttt{]}.
\item[\cercle{c}] \textit{$\underline{S}(\pi)$ satisfait à $L(\Phi)$, et pour tout $x\in\Ob\mathcal{M}$ et toute représentation $x\simeq\varinjlim\limits_{i\in I}x_i$ qui est $\Phi$-$\pi$-adaptée (cf. définition plus bas), on a ceci : $x\in\Ob\underline{S}$ si et seulement si $I_{\underline{S}}$ est cofinale dans $I$, où $I_{\underline{S}}$ est la sous-catégorie strictement pleine de $I$ formée des $i\in\Ob I$ tels que $x_i\in\Ob\underline{S}$ (ou encore, $x_i\in\Ob\underline{S}(\pi)$).} \textbf{N.B.} \textit{On dit qu'un système inductif $(x_i)_{i\in I}$ dans $\mathcal{M}$ est $\Phi$-$\pi$-adapté si}
\begin{itemize}[leftmargin=2em]
\item[\textit{1)}] \textit{$I$ équivalente à une petite catégorie, et $I$ grande devant $\pi$,}
\item[\textit{2)}] \textit{les $\varinjlim$ de type $\Phi$ existent dans $I$,}
\item[\textit{3)}] \textit{le foncteur $i\longmapsto x_i$ de $I$ dans $\mathcal{M}$ commute aux $\varinjlim$ de type $\Phi$, et enfin}
\item[\textit{4)}] \textit{les $x_i\in\Ob\mathcal{M}_\pi$.}
\end{itemize}
\item[c$'$)] \textit{(Si $\pi^{\pi_0}=\pi$, où $\pi_0$ est un cardinal $\leqslant\pi$ utile pour $\mathcal{M}$, et si $J\in\Phi$ implique que $J$ est grande devant $\pi_0$.) Comme c), avec $I$} ordonnée.
\end{itemize}
(\footnote{\textbf{N.B.} Quand les conditions restrictives ci-contre \texttt{[parenthèse dans c$'$)]} ne sont pas satisfaites, alors a, a$'$, b, c sont encore équivalentes, et impliquent c$'$, mais ne sont peut-être pas équivalentes à c$'$).},\footnote{\textbf{N.B.} Il faudrait réénoncer les conditions annexes [?] (Si \dots), cf. p. 359-360.}).

\textbf{Définition 4.13.23} (\footnote{\textbf{N.B.} Il y aurait intérêt à remonter cette définition (\textit{sans} b), dès avant 4.13.15.}). Soient $\mathcal{M}$ une catégorie accessible, $\pi$ un cardinal, $\Phi$ un ensemble de catégories, satisfaisant les deux conditions :
\begin{itemize}[leftmargin=2.5em]
\item[a)] $J\in\Phi\Longrightarrow\cardred J\leqslant\pi$.
\item[b)] Il existe un ordinal limite $\alpha_0$, tel que $I_{\alpha_0}\in\Phi$ (\footnote{Si c'est pour poser une définition, on peut se passer de la condition b).}).
\end{itemize}
\msp{page 344}
Soit $S$ une partie de $\Ob\mathcal{M}$, $\underline{S}$ la sous-catégorie pleine qu'elle définit. On dit que $S$ est une \textit{partie $\Phi$-$\pi$-accessible de} $\Ob\mathcal{M}$, ou que $\underline{S}$ est une \textit{sous-catégorie $\Phi$-$\pi$-accessible dans} $\mathcal{M}$, si les conditions suivantes sont satisfaites.
\begin{itemize}[leftmargin=2.5em]
\item[1°)] $\pi$ est utile pour $\mathcal{M}$, $\mathcal{M}$ est stable par $\varinjlim$ de type $\Phi$ (\footnote{Condition préliminaire sur $\mathcal{M}$.}).
\item[2°)] $\underline{S}$ est une sous-catégorie $\pi$-accessible dans $\mathcal{M}$.
\item[3°)] $\underline{S}$ est stable dans $\mathcal{M}$ par $\varinjlim$ de type $\Phi$.
\end{itemize}
(\footnote{\textbf{N.B.} Quand $\Phi=\{I_{\alpha_0}\}$, on dit aussi que $\underline{S}$ est $\alpha_0$-$\pi$-accessible dans $\mathcal{M}$. Si $\Phi$ est formé des catégories de cardinal $\leqslant\pi$ et grandes devant $\pi_0$, on dit que $\underline{S}$ est $\pi_0$-$\pi$-accessible dans $\mathcal{M}$.}).

\textbf{N.B.} Quand la condition préliminaire 1°) \textit{sur} $\mathcal{M}$ est satisfaite, la validité des conditions 2°) 3°) équivaut donc à la condition que $\underline{S}$ soit \textit{strictement} pleine dans $\mathcal{M}$, et que de
\printedpage{131}
% head: A. Grothendieck : Dérivateurs, Ch. XVIII, Édition des Universités de Montpellier II et Paris VII
% foot: Version 17 décembre 2025
plus $\underline{S}$ satisfasse la condition a) du théorème 4.13.22. D'autre part, cette condition est équivalente à chacune des conditions a$'$, b, c de ce même théorème, et aussi à la dernière condition c$'$, sous réserve qu'on ait $\pi^{\pi_0}=\pi$ (pour un cardinal $\pi_0<\pi$ utile pour $\mathcal{M}$) et que $J\in\Phi$ implique que $J$ est grande devant $\pi_0$.

Conceptuellement, c'est la définition donnée dans 4.13.23, i.e. la condition \cercle{a} de 4.13.22, qui paraît la plus claire, ou sa variante immédiate a$'$. Mais du point de vue pratique, du travail avec les parties accessibles, c'est par propriété \cercle{c} qui paraît la plus commode. Le poids principal du théorème 4.13.22 est bien
\msp{page 345}
l'implication a $\Longrightarrow$ c (via a $\Longrightarrow$ a$'$ $\overset{!}{\Longrightarrow}$ c). Donc il faut surtout voir ce théorème comme donnant un critère remarquable d'appartenance à $\Ob\underline{S}$ d'un objet $x$ de $\mathcal{M}$ qui est présenté par un système inductif $\Phi$-$\pi$-adapté :
\[ x\ \in\ \Ob\underline{S}\qquad\quad\Longleftrightarrow\qquad\quad I_{\underline{S}}\ \text{cofinal}\ , \tag*{$(*)$} \]
et dans cette équivalence, c'est l'implication $\Longrightarrow$ qui est la partie non soritale, et remarquable.

\textbf{Proposition 4.13.24.} \textit{Soient $\mathcal{M}$, $\pi$, $\Phi$ comme dans 4.13.23} (\footnote{\textbf{N.B.} La condition $I_{\alpha_0}\in\Phi$ n'est pas utile.})\textit{, on suppose que la condition 1° de 4.13.23 est satisfaite ($\pi$ utile pour $\mathcal{M}$, $\mathcal{M}$ stable par $\varinjlim$ de type $\Phi$). Soit ${\Sigma'}^{\Phi}_{\pi}$ l'ensemble des sous-catégories $\underline{S}$ de $\mathcal{M}$ qui sont $\Phi$-$\pi$-accessibles dans $\mathcal{M}$, et $\Sigma^{\Phi}_{\pi}$ l'ensemble des sous-catégories $\underline{S}(\pi)$ de $\mathcal{M}_\pi$ qui sont strictement pleines, et stables dans $\mathcal{M}_\pi$ par $\varinjlim$ de type $\Phi$. (Ces sous-catégories sont forcément karoubiennes, car stables par $\varinjlim$ de type $I_{\alpha_0}$, donc aussi par facteurs directs.) Alors les bijections (4.13.4.1) (page 288) induisent des bijections inverses l'une à l'autre}
\[
\begin{tikzcd}
\Sigma^{\Phi}_{\pi} \arrow[r, shift left] & {\Sigma'}^{\Phi}_{\pi} \arrow[l, shift left]
\end{tikzcd}\ .
\]

\msp{page 346}
% marge : trait manuscrit (reproduit en image) en regard de la Proposition 4.13.25
\textbf{Proposition 4.13.25.} \textit{Soient $\mathcal{M}$, $\pi$, $\Phi$ comme dans 4.13.23, avec $I_{\alpha_\pi}\in\Phi$ ($\alpha_\pi$ étant le plus petit ordinal de cardinal $\Phi$} \texttt{[plutôt} $\pi$\texttt{]}\textit{). Alors l'ensemble ${\Sigma'}^{\Phi}_{\pi}$ des sous-catégories de $\mathcal{M}$ $\Phi$-$\pi$-accessibles dans $\mathcal{M}$ est stable par intersection de familles de cardinal $\leqslant\pi$} (\footnote{Si $\pi$ n'est pas cardinal limite. Sinon, il faut dire : familles de cardinal $<\pi$. Cf. 4.13.10.3, où je donne un énoncé rectifié légèrement plus général.}).

(\textbf{N.B.} La bijection $\Sigma^{\Phi}_{\pi}\xrightarrow{\ \sim\ }{\Sigma'}^{\Phi}_{\pi}$ de 4.13.24 respecte les relations d'ordre naturelles d'inclusion, donc transforme intersections en intersections.)

Se voit comme 4.13.7 (qui est contenu dans l'énoncé actuel), en utilisant le lemme 4.13.10.

\textbf{4.13.26.} Soient $\mathcal{M}$, $\pi$, $\Phi$ comme dans 4.13.23, et soit $\underline{S}$ une sous-catégorie \textit{strictement pleine} de $\mathcal{M}$, $f:\underline{S}\longrightarrow\mathcal{M}$ le foncteur d'inclusion. La condition que $\underline{S}$ soit $\Phi$-$\pi$-accessible dans $\mathcal{M}$ se décompose en :
\begin{itemize}[leftmargin=2.5em]
\item[1)] $\pi$ est utile pour $\mathcal{M}$ \textit{et} pour $\underline{S}$.
\item[2)] $\mathcal{M}$ et $\underline{S}$ sont stables par $\varinjlim$ de type $\Phi$.
\item[3)] $f$ commute auxdites limites.
\item[4)] $f$ est $\pi$-accessible.
\item[5)] $f(\underline{S}_\pi)\subset\mathcal{M}_\pi$.
\end{itemize}

\printedpage{132}
% head: A. Grothendieck : Dérivateurs, Ch. XVIII, Édition des Universités de Montpellier II et Paris VII
% foot: Version 17 décembre 2025
\msp{page 347}
Soit alors $\pi'>\pi$ un autre cardinal, quand peut-on assurer que si $\underline{S}$ est $\Phi$-$\pi$-accessible dans $\mathcal{M}$, il soit aussi $\Phi$-$\pi'$-accessible\,? Les conditions 2) 3) 4) sont visiblement stables par passage de $\pi$ à $\pi'$. Restent les conditions 1) et 5). On sait que si $\pi'^{\pi}=\pi'$, alors la condition 1) reste valable pour $\pi'$. Plus généralement, si $\pi_0$ est un cardinal $\leqslant\pi$ tel que $\pi_0$ soit utile pour $\mathcal{M}$ et pour $\underline{S}$, alors si $\pi'>\pi_0$ est tel que $\pi'^{\pi_0}=\pi'$, alors $\pi'$ est aussi utile pour $\mathcal{M}$ et pour $\underline{S}$. Si on admet l'HCG (hypothèse du continu généralisée), il en est ainsi pour tout $\pi'>\pi_0$, donc pour tout $\pi'>\pi$.

Reste la condition 5), pour laquelle une réponse est donnée par 4.12.1 (p. 273). Si $\pi_0\leqslant\pi$ est choisi de façon que $f$ soit $\pi_0$-accessible (p.ex. $\pi_0=\pi$), et si
\[
\begin{aligned}
\pi'_1 &= \cardred\underline{S}_{\pi_0} \qquad\quad (\,\leqslant\cardred\underline{S}(\pi)\,)\\
\pi''_1 &= \sup(\pi'_1,\pi)\ ,
\end{aligned}
\]
on voit que si $\pi'\geqslant\pi''_1$ (i.e. $\pi'\geqslant\pi'_1,\pi$) est tel que $\pi'^{\pi_0}=\pi'$, alors on a bien
\[ f(\underline{S}_{\pi'})\ \subset\ \mathcal{M}_{\pi'}\ . \]

\msp{page 348}
Si donc de plus $\pi'>\pi\geqslant\pi_0$ (comme supposé au début), et si on suppose l'HCG, alors la condition $\pi'^{\pi_0}=\pi'$ est automatique, et il reste la seule condition $\pi'\geqslant\pi''_1$, i.e. en définitive $\pi'\geqslant\cardred\underline{S}_{\pi_0}$. Pour résumer :

\textbf{Proposition 4.13.27.} \textit{Sous les conditions de 4.13.23 (sans condition b) sur $\Phi$), soit $\underline{S}$ une sous-catégorie de $\mathcal{M}$ qui est $\Phi$-$\pi$-adaptée. Soit $\pi_0\leqslant\pi$ un cardinal qui est utile pour $\underline{S}$ et pour $\mathcal{M}$, et tel que $f:\underline{S}\longrightarrow\mathcal{M}$ soit $\pi_0$-accessible, i.e. tel que $\underline{S}$ soit stable dans $\mathcal{M}$ par $\varinjlim$ grandes devant $\pi_0$ (on peut p.ex. prendre $\pi_0=\pi$). Alors pour tout cardinal $\pi'>\pi$ tel que}
\begin{itemize}[leftmargin=2.5em]
\item[\textit{a)}] \textit{$\pi'^{\pi_0}=\pi'$ et}
\item[\textit{b)}] \textit{$\pi'\geqslant\cardred\underline{S}_{\pi_0}$,}
\end{itemize}
\textit{$\underline{S}$ est $\Phi$-$\pi'$-accessible dans $\mathcal{M}$. Donc si $\pi'>\pi$ est tel que $\pi'^{\pi}=\pi'$, $\pi'>\cardred\underline{S}_\pi$, alors $\underline{S}$ est $\Phi$-$\pi'$-accessible. (Il suffit p.ex. de prendre $\pi'=2^c$, avec $c\geqslant\sup(\pi,\cardred\underline{S}(\pi))$.) Si on admet l'HCG, il suffit donc qu'on ait $\pi'>\sup(\pi,\cardred\underline{S}(\pi))$.}

\msp{page 349}
\textbf{Proposition 4.13.28} (\footnote{À rejeter plus bas, après 4.13.31, qui est la première application importante de 4.13.22, alors que 4.13.28 se prouve sans avoir à utiliser 4.13.22 ou ses variantes.}). \textit{Soient $\mathcal{M}$, $\pi$, $\Phi$ comme dans 4.13.23. Supposons de plus que $J\in\Phi$ implique que $J$ est} filtrante. \textit{Alors l'ensemble $\tilde{\Sigma}^{\Phi}_{\pi}$ des sous-ensembles $\Phi$-$\pi$-accessibles de $\mathcal{M}$ est stable par réunions finies. Plus généralement, soit $\lambda$ un cardinal $<\pi$ tel que $\forall\ J\in\Phi$, $J$ soit grande devant $\lambda$. Alors l'ensemble $\tilde{\Sigma}^{\Phi}_{\pi}$ est stable par réunions $(S_r)_{r\in A}$, indexées par des ensembles $A$ de cardinal $\leqslant\lambda$.}

\textbf{Corollaire 4.13.29.} \textit{L'ensemble des parties accessibles de} $\Ob\mathcal{M}$ \textit{est stable par “petites” réunions -- i.e. si $(S_r)_{r\in A}$ est une petite famille de parties accessibles de} $\Ob\mathcal{M}$, \textit{alors $\bigcup_{r\in A}S_r$ est aussi accessible.}

La démonstration de 4.13.28 repose sur le lemme suivant (je n'explicite par la démonstration) : % [sic] par la démonstration -> pas la démonstration, sense of the sentence

\printedpage{133}
% head: A. Grothendieck : Dérivateurs, Ch. XVIII, Édition des Universités de Montpellier II et Paris VII
% foot: Version 17 décembre 2025
\msp{page 350}
\textbf{Lemme 4.13.30.} \textit{Soit $I$ une catégorie grande devant $\pi$, et $(J_r)_{r\in A}$ une famille des sous-catégories pleines telles que $\card A\leqslant\pi$, telle que $\bigcup_{r\in A}\Ob J_r$ soit cofinal dans $\Ord(I)$, i.e. pour tout $i\in\Ob I$ $\exists\ j\in\bigcup_{r\in A}\Ob J_r$, et une flèche $i\longrightarrow j$. Alors il existe $r\in A$ tel que $J_r$ soit cofinal dans $I$.}

\textbf{Théorème 4.13.31.} \textit{Soient $\mathcal{M}$, $\mathcal{N}$ des catégories accessibles, $f:\mathcal{M}\longrightarrow\mathcal{N}$ un foncteur accessible, $T\subset\Ob\mathcal{N}$ une partie de $\mathcal{N}$, $\underline{T}\subset\mathcal{N}$ la sous-catégorie strictement pleine correspondante de $\mathcal{N}$, $S=f^{-1}(T)$, $\underline{S}=f^{-1}(\underline{T})$ la sous-catégorie pleine de $\mathcal{M}$ associée. Si $T$ est accessible, alors $S$ est accessible\,; en d'autres termes, si $\underline{T}$ est accessible dans $\mathcal{N}$, $\underline{S}$ est accessible dans $\mathcal{M}$.}

Cela résulte de l'énoncé plus précis :

\textbf{Corollaire 4.13.32.} \textit{Soient $\pi_0$ et $\pi>\pi_0$ des cardinaux ayant les propriétés suivantes :}

\msp{page 351}
\begin{itemize}[leftmargin=2.5em]
\item[\textit{a)}] \textit{$\underline{T}$ est $\pi_0$-$\pi$-accessible dans $\mathcal{N}$ (4.13.23).}
\item[\textit{b)}] \textit{$\pi_0$ est utile pour $\mathcal{M}$ (il l'est pour $\mathcal{N}$, à cause de a)).}
\item[\textit{c)}] \textit{$f(\mathcal{M}_\pi)\subset\mathcal{N}_\pi$.}
\item[\textit{d)}] \textit{$\mathcal{M}$ satisfait à $L_{\pi_0}$, et $f$ est $\pi_0$-accessible.}
\end{itemize}
\textit{(L'existence d'un tel couple $(\pi_0,\pi)$ résulte aussitôt de 4.13.16 et de 4.12.1.) Alors $\underline{S}$ est $\pi_0$-$\pi$-accessible dans $\mathcal{N}$.}

Comme $f$ commute aux $\varinjlim$ grandes devant $\pi_0$ en vertu de d), et que $\underline{T}$ est stable par ces limites par a) (cf. 4.13.23, 3°), il en est donc de même de $\underline{S}$, i.e. le 3° de la définition 4.13.23 est satisfait. D'autre part, la condition 1° ($\pi$ utile pour $\mathcal{M}$, $\mathcal{M}$ stable par $\varinjlim$ grandes devant $\pi_0$ et de cardinal $\leqslant\pi$) est vérifiée par b) et d). Donc il reste à vérifier la condition 2°, i.e. que $\underline{S}$ est $\pi$-accessible. Posons
\[
\begin{aligned}
\underline{T}(\pi) &= \underline{T}\cap\mathcal{N}_\pi\\
\underline{S}(\pi) &= \underline{S}\cap\mathcal{M}_\pi\ \ \underbrace{\ \subset\ }_{\text{grâce à c)}}\ \ f^{-1}(\underline{T}(\pi))\ .
\end{aligned}
\]
Comme $\underline{S}$ est stable dans $\mathcal{M}$ par $\varinjlim$ grandes devant $\pi_0$ (lesquelles existent dans $\mathcal{M}$ par d), et a fortiori par $\varinjlim$ grandes devant $\pi$, on a un foncteur canonique pleinement fidèle
\[ \Ind_\pi\underline{S}(\pi)\ \longrightarrow\ \underline{S}\ . \]

\msp{page 352}
D'ailleurs, $\underline{T}$ étant strictement pleine dans $\mathcal{N}$, $\underline{S}$ est strictement pleine dans $\mathcal{M}$, et il reste à vérifier que le foncteur précédent est essentiellement surjectif. Soit donc $x\in\Ob\underline{S}$. On aura, grâce à b),
\[ x\ =\ \varinjlim_{I}x_i\ ,\qquad\quad x_i\in\mathcal{M}_\pi,\ I=\mathcal{M}_\pi/x\ . \]
D'ailleurs, le système inductif précédent est $\pi_0$-$\pi$-adapté, i.e. (p. 339) :
\begin{itemize}[leftmargin=2.5em]
\item[1)] $I$ équivalente à une petite catégorie, $I$ grande devant $\pi$ (résulte du fait que $\pi$ est utile pour $\mathcal{M}$, cf.\,\cercle{b}\,).
\end{itemize}

% ---- from der-ch18_p136-140.tex
\printedpage{134}
% head: A. Grothendieck : Dérivateurs, Ch. XVIII, Édition des Universités de Montpellier II et Paris VII
% foot: Version 17 décembre 2025
\begin{itemize}[leftmargin=2.5em]
\item[2)] Les limites inductives grandes devant $\pi_0$ est de cardinal $\leqslant \pi$ existent dans $I$ (résulte de $\mathrm{L}_{\pi_0}$ sur $\mathcal{M}$ (cf.\,\cercle{d})). % [sic] est de cardinal -> et de cardinal, coordination ("grandes devant π0 et de cardinal ≤ π")
\item[3)] Le foncteur $\xi : I \longrightarrow \mathcal{M}$ y commute.
\item[4)] Les $x_i$ sont dans $\mathcal{M}_\pi$.
\end{itemize}
Comme $f$ commute aux $\limind$ grandes devant $\pi_0$ (cf.\,\cercle{1}), a fortiori à celles grandes devant $\pi$, on a donc
\[ f(x) \;\simeq\; \varinjlim_{I} f(x_i) \ , \]
je dis que le système inductif $(f(x_i))_{i\in I}$ dans $\mathcal{N}$ est $\pi_0$-$\pi$-adapté. Les conditions 1), 2) sont \textit{ne varietur} (ne concernent que $I$), le foncteur $I \longrightarrow \mathcal{N}$ est le composé $I \xrightarrow{\xi} \mathcal{M} \xrightarrow{f} \mathcal{N}$, donc commute aux $\limind$ grandes devant $\pi_0$ et de cardinal $\leqslant \pi$ (\footnote[155]{cette commutation devrait avoir un nom, condition $\mathrm{L}(\pi_0,\pi)$ sur un foncteur.}), puisqu'il en est ainsi de $\xi$ et de $f$, enfin les $f(x_i)$ sont dans $\mathcal{N}_\pi$ puisque par \cercle{c} $f(\mathcal{M}_\pi) \subset \mathcal{N}_\pi$.

\msp{page 353}

Par ailleurs $f(x) \in \Ob \underline{T}$ (puisque $x \in \Ob \underline{S}$, où $\underline{S} \defeq f^{-1}(\underline{T})$). Donc on peut appliquer 4.13.22 a $\Longrightarrow$ c à la situation $(\mathcal{N}, \pi_0, \pi, \underline{S})$ et au système inductif $(f(x_i))_{i\in I}$, pour conclure que la sous-catégorie $I'$ de $I$ formée des $i$ tels que $f(x_i) \in \Ob \underline{S}$ \texttt{[plutôt Ob}~$\underline{T}$\texttt{]}, i.e. tels que $x_i \in \underline{S}$, donc $x_i \in \underline{S}(\pi)$, est \textit{cofinale} dans $I$. Donc on a
\[ x \;\simeq\; \varinjlim_{i\in I'} x_i \ , \]
avec cette fois les $x_i$ dans $\underline{S}(\pi)$, et $I'$ grande devant $\pi$ (car cofinale et pleine dans $I$, qui est grande devant $\pi$), ce qui prouve que $x$ est dans l'image essentielle de $\Ind_\pi(\underline{S}(\pi))$, q.e.d.

\textbf{Question 4.13.33.} On vient de voir que si $f : \mathcal{M} \longrightarrow \mathcal{N}$ est un foncteur accessible entre catégories accessibles, alors l'image inverse d'une partie accessible de $\Ob \mathcal{N}$ est une partie accessible de $\Ob \mathcal{M}$. Je conjecture que l'énoncé similaire, non pas pour les images directes, mais pour les \textit{images essentielles}, est
\msp{page 354}
vraie, i.e. que pour toute sous-catégorie $\underline{S}$ de $\mathcal{M}$ accessible dans $\mathcal{M}$, son image essentielle dans $\mathcal{N}$ est accessible dans $\mathcal{N}$ (\footnote[156]{Faux, même si on prend la $\natural$-clôture de l'image essentielle (“image $\natural$-essentielle”).}). Dans cette question, on peut supposer que $\underline{S} = \mathcal{M}$, et la question revient à celle-ci : tout foncteur accessible $f : \mathcal{M} \longrightarrow \mathcal{N}$ entre catégories accessibles se factorise-t-il en un foncteur (accessible) \textit{essentiellement surjectif}, et un foncteur accessible pleinement fidèle $\mathcal{M} \longrightarrow \underline{T} \longrightarrow \mathcal{N}$ (où $\underline{T}$ est elle-même accessible) (\footnote[157]{Il faut se borner tout au moins à essentiellement $\natural$-surjectif – mais même cela est trop optimiste.})~? Si oui, cela me paraîtrait peut-être le résultat le plus profond (peut-être le seul véritablement profond) en théorie des catégories. % [sic] vraie -> vrai, agreement with « l'énoncé similaire »

Voici une question qui me paraît très voisine. Soient $\mathcal{M}$ une catégorie accessible, $A$ une sous-catégories pleine petite. Supposons que $\mathcal{M}$ satisfasse à $\mathrm{L}_{\pi_0}$ (p.ex. $\pi_0$ utile pour $\mathcal{M}$), et soit $\underline{S}$ la plus petite sous-catégorie strictement pleine de $\mathcal{M}$ contenant $A$, et stable par limites grandes devant $\pi_0$. Je conjecture que $\underline{S}$ est accessible dans $\mathcal{M}$. % [sic] sous-catégories -> sous-catégorie, singular after « une »

\msp{page 355}

De même, si $\Phi$ est un petit ensemble de petites catégories, tel que $\mathcal{M}$ soit stable par $\limind$ de type $\Phi$, on peut prendre l'enveloppe $\underline{S}$ de $A$ pour les $\limind$ grandes devant $\pi_0$ \textit{et} les
\printedpage{135}
% head: A. Grothendieck : Dérivateurs, Ch. XVIII, Édition des Universités de Montpellier II et Paris VII
% foot: Version 17 décembre 2025
$\limind$ de type $\Phi$, et je conjecture encore que $\underline{S}$ est accessible. Dans cette direction, on a du moins le résultat très partiel suivant.

\textbf{Proposition 4.13.34.} \textit{Soit $\mathcal{M}$ une catégorie accessible stable par petites $\limind$ filtrantes, et soit $\underline{S}$ une sous-catégorie strictement pleine de $\mathcal{M}$. Conditions équivalentes :}
\begin{itemize}[leftmargin=2.5em]
\item[\textit{a)}] \textit{$\underline{S}$ est accessible dans $\mathcal{M}$, et stable dans $\mathcal{M}$ par petites $\limind$ filtrantes.}
\item[\textit{b)}] \textit{Il existe une petite sous-catégorie pleine $A$ de $\mathcal{M}$, telle que $\underline{S}$ soit la plus petite sous-catégorie strictement pleine de $\mathcal{M}$ contenant $A$ et stable par petites $\limind$ filtrantes.}
\end{itemize}

\msp{page 356}

On a a) $\Longrightarrow$ b), car si $\underline{S}$ est $\pi$-accessible dans $\mathcal{M}$, il est immédiat (sans autre condition de stabilité sur $\underline{S}$) que $\underline{S}$ est la plus petite sous-catégorie strictement pleine de $\mathcal{M}$ contenant $\underline{S}(\pi) \defeq \underline{S}\cap\mathcal{M}_\pi$ et stable par $\limind$ grandes devant $\pi$. Si donc on sait déjà que $\underline{S}$ est même stable par \textit{toutes} $\limind$ filtrantes petites, il s'ensuit que $\underline{S}$ est la plus petite sous-catégorie de $\mathcal{M}$ ayant cette propriété de stabilité plus forte, et contenant $A$.

Prouvons b $\Longrightarrow$ a. Soit $\pi$ un cardinal utile pour $\mathcal{M}$, et tel que $A \subset \mathcal{M}_\pi$. Soit $B$ la plus petite sous-catégorie strictement pleine de $\mathcal{M}$ contenant $A$ et stable par $\limind$ filtrantes de cardinal $\leqslant \pi$. On a donc $B \hookrightarrow \mathcal{M}_\pi$ et $B \subset \underline{S}$, et $\underline{S}$ est aussi la L-enveloppe de $B$. D'ailleurs $B$ est stable par
\msp{page 357}
facteurs directs. Considérons alors l'image essentielle $\underline{S}'$ de $\Ind_\pi(B)$ dans $\mathcal{M}$. On a $\underline{S}' \subset \underline{S}$, et tout revient à prouver que $\underline{S}' = \underline{S}$, d'où résultera en effet que $\underline{S}$ est $\pi$-accessible dans $\mathcal{M}$, a fortiori accessible dans $\mathcal{M}$, d'où a). Pour prouver que $\underline{S}' = \underline{S}$, par définition de $\underline{S}$ comme la L-enveloppe de $A$, il suffit de voir que $\underline{S}'$ est stable par petites $\limind$ filtrantes. Or cela résulte du fait que $B \subset \mathcal{M}_\pi$ est stable dans $\mathcal{M}_\pi$ par $\limind$ filtrantes de cardinal $\leqslant \pi$, et de 4.13.19, qui implique tout au moins, comme $B = \underline{S}' \cap \mathcal{M}_\pi$ (en vertu de 4.13.4) que $\underline{S}'$ est stable par $\limind$ filtrantes \textit{de cardinal} $\leqslant \pi$. Qu'il soit même stable par petites $\limind$ filtrante quelconques, résulte de ceci, et du fait qu'il l'est par $\limind$ grandes devant $\pi$. (Cf. 4.13.21 \dots) % [sic] filtrante quelconques -> filtrantes quelconques, plural agreement

\msp{page 358}

\textbf{Définition 4.13.35.} Soient $\mathcal{M}$ une catégorie accessible, stable par petites $\limind$ filtrantes, $S$ une partie de $\Ob \mathcal{M}$, $\underline{S}$ la sous-catégorie pleine correspondante dans $\mathcal{M}$. On dit que $S$ est \textit{L-accessible}, ou que $\underline{S}$ est L-accessible dans $\mathcal{M}$, si $\underline{S}$ est accessible dans $\mathcal{M}$, et de plus stable par petites $\limind$ filtrantes (\footnote[158]{\textbf{N.B.} L'ensemble des parties L-accessibles de $\Ob \mathcal{M}$ est stable par petites intersections, et par réunions \textit{finies}. L'image inverse d'une partie L-accessible par un foncteur (\footnotemark[159]) $f : \mathcal{M} \longrightarrow \mathcal{N}$ de catégories accessibles est L-accessible.}).\footnotetext[159]{un \textit{$L_0$-foncteur} : qui commute aux petites $\limind$ filtrantes.}

\textbf{Proposition 4.13.36.} \textit{Soient $\mathcal{M}$, $S$, $\underline{S}$ comme dans 4.13.35, et soit $\pi$ un cardinal utile pour $\mathcal{M}$. Supposons $\underline{S}$ stable dans $\mathcal{M}$ par petites $\limind$ filtrantes, ou seulement $\underline{S}(\pi) \defeq \underline{S} \cap \mathcal{M}_\pi$ stable par petites $\limind$ filtrantes de cardinal $\leqslant \pi$. Conditions équivalentes :}
\begin{itemize}[leftmargin=2.5em]
\item[\textit{a)}] \textit{$\underline{S}$ est $\pi$-accessible. (Il sera automatiquement stable par petites $\limind$ filtrantes, donc L-accessible.)}
\item[\textit{b)}] \textit{Pour tout $x \in \Ob \mathcal{M}$, on a $x \in \Ob \underline{S}$ si et seulement si $\underline{S}(\pi)/x$ est cofinale dans $\mathcal{M}_\pi/x$.}
\printedpage{136}
% head: A. Grothendieck : Dérivateurs, Ch. XVIII, Édition des Universités de Montpellier II et Paris VII
% foot: Version 17 décembre 2025
\item[\textit{c)}] \textit{Pour tout $x \in \Ob \mathcal{M}$, et toute représentation $x \simeq \varinjlim\limits_{i\in I} x_i$}
\msp{page 359}
\textit{qui est L-$\pi$-adaptée (cf. page 343, c), on a ceci : $x \in \Ob S$ si et seulement si $I_{\underline{S}}$ est cofinal dans $I$.}
\item[\textit{c}$'$\textit{)}] \textit{Comme c), avec $I$ ordonnée.}
\item[\textit{c}$''$\textit{)}] \textit{(Si on a choisi, pour tout $x \in \Ob \mathcal{M}$,} une \textit{représentation $\pi$-adaptée $x \simeq \varinjlim\limits_{I} x_i$) : $x \in \Ob \underline{S} \Longleftrightarrow$ l'ensemble $I_{\underline{S}}$ est cofinal dans $I$.}
\end{itemize}
\textsc{Démonstration}. C'est essentiellement l'énoncé de 4.13.22 dans le cas particulier où $\Phi =$ ensemble des catégories \textit{filtrantes} de cardinal $\leqslant \pi$, mais où, pour simplifier, on a supposé d'avance que $\underline{S}(\pi)$ satisfait à $\mathrm{L}(\Phi)$. L'équivalence de a) b) c) est contenue dans 4.13.22, ainsi que c $\Longrightarrow$ c$'$. Mais pour c$'$ $\Longrightarrow$ a, l'hypothèse accessoire $\pi^{\pi_0} = \pi$ et que les $J \in \Phi$ seront grandes devant $\pi_0$, n'est pas satisfaite ici, et il faut un démonstration indépendante, calquée sur celle des pages 333-336. Il suffit pour cela, à l'endroit \texttt{[de]} $\square$, page 333, d'invoquer le % [sic] un démonstration -> une démonstration, gender agreement

\textbf{Lemme 4.13.37.} \textit{Soient $\mathcal{M}$ une catégorie}
\msp{page 360}
\textit{accessible stable par $\limind$ filtrantes, $\pi$ un cardinal utile pour $\mathcal{M}$. Alors pour tout $x \in \Ob \mathcal{M}$, on a une représentation}
\[ x \;=\; \varinjlim_{i\in I} x_i \]
\textit{L-$\pi$-adaptée (i.e. avec}
\begin{itemize}[leftmargin=2.5em]
\item[\textit{1)}] \textit{$I$ équivalente à une petite catégorie, et $I$ grande devant $\pi$,}
\item[\textit{2)}] \textit{$I$ stable par $\limind$ filtrantes de cardinal $\leqslant \pi$,}
\item[\textit{3)}] \textit{$I \longrightarrow \mathcal{M}$ commute à ces $\limind$ ,}
\item[\textit{d)}] \textit{$x_i \in \Ob \mathcal{M}_\pi \ \forall\, i \in I$),} % [sic] d) -> 4), the items are numbered 1), 2), 3)
\end{itemize}
\textit{avec de plus $I$} ordonné.

En effet, on part de
\[ x \;=\; \varinjlim_{j\in J} x_j \ , \qquad\quad J \;\simeq\; \underbrace{(\mathcal{M}_\pi/x)_0}_{\substack{\text{sous-catégorie réduite}\\ \text{équivalente à } \mathcal{M}_\pi/x}} \qquad . \]
On définit $I$ comme l'ensemble des sous-catégories \textit{filtrantes} $J'$ de $J$, de cardinal \texttt{[“cardinal” biffé]} $\leqslant \pi$, $I$ ordonné par inclusion, et on pose, pour $J' \in I$,
\[ x_{J'} \;=\; \varinjlim_{j\in J'} x_j \ . \]
Il est immédiat qu'on obtient ainsi un système inductif dans $\mathcal{M}$, indexé par $I$.

\msp{page 361}

D'autre part, il est clair que $I$ est petite, je dis que $I$ est stable par $\limind$ filtrantes de cardinal $\leqslant \pi$, et que le foncteur $J' \longmapsto x_{J'} : I \longrightarrow \mathcal{M}$ y commute, enfin que les $x_{J'}$ sont dans $\mathcal{M}_\pi$. Il reste à voir que $I$ est grande devant $\pi$, i.e. que si on a $(J'_r)_{r\in A}$, avec $\card A \leqslant \pi$, les $J'_r$ dans $I$, alors $\exists\, J' \in I$ qui majore les $J'_r$. Soit $K$ la sous-catégorie de $J$ engendrée par les $J'_r$, on voit facilement que $\card K \leqslant \pi$. Tout \texttt{[“tout” biffé]} revient à prouver :

\printedpage{137}
% head: A. Grothendieck : Dérivateurs, Ch. XVIII, Édition des Universités de Montpellier II et Paris VII
% foot: Version 17 décembre 2025
\textbf{Lemme 4.13.38} (\footnote[160]{\textbf{N.B.} Pour la démonstration de 4.13.37, il suffit de prouver 4.13.38 dans le cas où $J$ est ordonnée, auquel cas les contorsions de la fin sont inutiles.}). \textit{Soient $J$ une catégorie grande devant $\pi$, $K$ une sous-catégorie de $J$ de cardinal $\leqslant \pi$. Alors il existe une sous-catégorie $K'$} filtrante \textit{de $J$, contenant $K$, de cardinal $\leqslant \pi$. Si on suppose que l'on n'est pas, dans le cas où $J$ contient un objet maximal, et que $K$ contienne tous les objets maximaux de $J$,}
\msp{page 362}
\textit{on peut même prendre $K'$ telle que $\Ob K' = \Ob K \cup \{e\}$, où $e$ est un} objet final \textit{de $K'$, et $K$ est une sous-catégorie pleine de $K'$. (Donc si $K' \neq K$, $K'$ est isomorphe au cône à droite construit sur $K$, “déduit de $K$ en lui ajoutant un objet final”.)}

\textsc{Démonstration}. Si $K$ admet un objet final, on prend $K' = K$. Sinon, comme $J$ est grande devant $\card \Ob K$, il existe un $j \in \Ob J$ qui majore tous les objets de $K$. Choisissons, pour tout $k \in \Ob K$, une flèche
\[ \alpha_k \;:\; k \longrightarrow j \ . \]
Pour toute flèche
\[ u \;:\; k' \longrightarrow k \]
dans $K$, considérons la double-flèche
\[ (\alpha_k \circ u, \alpha_{k'}) \;:\; k' \rightrightarrows j \ . \]
L'ensemble de ces double-flèches, de but $j$, est de cardinal $\leqslant \card \Fl K \leqslant \pi$, donc $J$ étant grande devant $\pi$, il existe une flèche $\rho : j \longrightarrow e$ dans $J$ qui égalise toutes ces double-flèches. Soit $K' \subset J$, ayant comme objets $\Ob K \cup \{e\}$, et ayant comme flèches, en plus de
\msp{page 363}
celles de $K$, les flèches $\rho\alpha_k : k \longrightarrow e$ (pour $k \in \Ob K$), et $\id_e$. \textit{Si} $e \notin \Ob K$, on voit que c'est bien une sous-catégorie, admettant $e$ comme objet final et $K$ comme sous-catégorie pleine. Si on avait $e \in \Ob K$, et s'il existe un objet $e'$ de $J$ qui majore $e$ et qui n'est pas dans $K$, on se ramène aussitôt au cas précédent, en remplaçant $e$ par $e'$.

$\Biggl($ (\footnote[161]{La construction qui suit est contenue dans l'énoncé plus général et plus précis 4.13.39.}). Donc il reste le cas où il existe un objet $e$ dans $K$ qui est un plus grand élément dans $K$, et tel que tout objet de $J$ majorant $e$ soit dans $K$. Cela implique, $J$ étant filtrante, que $e$ est un plus grand objet de $J$. (Car si $x \in J$, $\exists\, x'$ majorant commun de $x$ et $e$, donc $x'$ est dans $K$, donc majoré par $e$, qui majore donc $x$.) On voit comme tantôt qu'il existe une flèche $p : e \longrightarrow j$ dans $J$, qui égalise toutes les double-flèches $k \rightrightarrows e$ de but $e$ \textit{dans}~$K$. Quitte à composer avec $j \longrightarrow e$, on peut supposer $j = e$. On a donc un endomorphisme $p$ de $e$ dans $J$, tel que $pu = pv$ pour toute double-flèche $k \overset{u}{\underset{v}{\rightrightarrows}} e$ dans $K$, de but $e$.

\msp{page 364}

Pour $k \in \Ob K$, soit
\[ f_k \;:\; k \longrightarrow e \]
la valeur commune des $pu$, pour $u \in \Hom_K(k,e)$. On voit tout de suite que si on a $k' \xrightarrow{v} k$ dans $K$, alors
\[ f_{k'} \;=\; f_k \circ v \ . \]
\printedpage{138}
% head: A. Grothendieck : Dérivateurs, Ch. XVIII, Édition des Universités de Montpellier II et Paris VII
% foot: Version 17 décembre 2025
En particulier, prenant $k' = e$, on trouve
\[ f_k \circ v \;=\; f_e \,(=\ p\ ) \qquad\quad \text{si } v : e \longrightarrow k \ . \]
Soit donc $K'$ la sous-catégorie de $J$ ayant mêmes objets que $K$, et comme flèches, en plus de celles de $K$, les flèches composées $k \xrightarrow{u} e \xrightarrow{p^r} e \xrightarrow{v} k'$, avec $u$, $v$ dans $K$, et $r \geqslant 1$, i.e. les flèches composées $k \xrightarrow{f_k} e \xrightarrow{p^{r'}} e \xrightarrow{v} k'$, avec $k, k' \in \Ob K$, $v \in \Fl K$, $r' \geqslant 0$. Il faut voir que c'est stable par composition (pour avoir bien une sous-catégorie). C'est clair pour $vu$ si $u$ ou $v$ sont dans $\Fl K$. Si \texttt{[ni]} $u$ ni $v$ ne le sont, on doit prendre un composé
\[
\begin{tikzcd}[column sep=large]
k \arrow[r, "u_1"] & e \arrow[r, "p^r", "r\geqslant 1"'] & e \arrow[r, "u_2"] \arrow[rr, bend right=20] & k' \arrow[r, "v_1"] & e \arrow[r, "p^s", "s\geqslant 1"'] & e \arrow[r, "v_2"] & k'' \ ,
\end{tikzcd}
\]
avec $u_1$, $u_2$, $v_1$, $v_2$ dans $\Fl K$, $r, s \geqslant 1$. Or on a
\[ p(v_1u_2) \;=\; p \ , \qquad\quad \text{d'où} \qquad\quad p^s(v_1u_2) \;=\; p^{s-1}(pv_1u_2) \;=\; p^{s-1}p \;=\; p^s \ , \]
donc le composé est
\[ k \xrightarrow{\;u_1\;} e \xrightarrow{\;p^{r+s}\;} e \xrightarrow{\;v_2\;} k'' \ , \]
qui
\msp{page 365}
est donc bien dans $\Fl K'$.

Ainsi, $K'$ est une sous-catégorie de $J$, mais est-elle filtrante ? La condition PS 1$'$ est évidente, puisque $e$ est un plus grand élément de $J$, reste la condition PS 2 de la double-flèche. On voit aisément que cette condition est satisfaite dans $K'$, si et seulement si elle est satisfaite pour toute double-flèche
\[ e \overset{p^r}{\underset{p^s}{\rightrightarrows}} e \qquad\quad \text{avec } r, s \geqslant 0 \ . \]
(En effet, on se ramène aussitôt aux double-flèches de but $e$, et on note que toute flèche de $K'$ de but $e$ s'écrit comme un composé $k \xrightarrow{u} e \xrightarrow{p^r} e \xrightarrow{v} e$, avec $r \geqslant 0$ et $u$, $v$ des flèches dans $K'$. En composant avec $e \xrightarrow{p} e$, on trouve alors, comme $pv = pu$, $k \xrightarrow{u} e \xrightarrow{p^{r+1}} e$, soit $k \xrightarrow{f_k} e \xrightarrow{p^r} e$.) Mais on voit aussitôt que la possibilité d'égaliser $(p^r, p^s)$ dans $K'$ équivaut à l'existence de $t \geqslant 0$ tel que $p^{r+t} = p^{s+t}$. Donc $K'$ est filtrante si et seulement si $r \longmapsto p^r$ est une suite stationnaire (faire $r = 0$, $s = 1$, d'où $p^r = p^{r+1}$ pour $r$ grand). Ce n'est
\msp{page 366}
pas forcément le cas.

Posons donc $K = K_0$, $K' = K_1$, et construisons par induction une suite croissante de sous-catégories $K_i$ ($i \in \mathbf{N}$), avec $\Ob K_i = \Ob J$, et (pour $i \geqslant 1$) $K_i$ étant déduites de $K_{i-1}$ en lui rajoutant une flèche $p_i : e \longrightarrow e$, comme ci-dessus. Ainsi toute double-flèche dans un $K_i$ peut s'égaliser dans $K_{i+1}$. Prenons $K_\infty = \limind K_i$, $K_\infty$ est une sous-catégorie \textit{filtrante} de $J$, contenant $K$ et de cardinal $\leqslant \pi$, q.e.d. $\Biggr)$

On peut préciser cette construction ainsi :

% ---- from der-ch18_p141-145.tex
\printedpage{139}
% head: A. Grothendieck : Dérivateurs, Ch. XVIII, Édition des Universités de Montpellier II et Paris VII
% foot: Version 17 décembre 2025

\noindent\textbf{Corollaire 4.13.39.} {\itshape Sous les conditions de 4.13.38, supposons que $K$ contienne un objet $e$ maximal dans $J$. Soit $\Sigma$ l'ensemble des endomorphismes $p$ de $e$ dans $J$, tels que pour toute double-flèche $k \overset{u}{\underset{v}{\rightrightarrows}} e$ \textup{dans} $K$, on ait $pu=pv$. Alors $\Sigma$ est stable par composition, et il existe une partie non vide $\Sigma_\infty$ \edbr{[}$=\Sigma_0$\edbr{]} de $\Sigma$, telle que pour $p,p'\in\Sigma_0$, existe
\msp{page 367}
$q\in\Sigma_0$ tel que $qp=qp'$ (i.e. telle que la sous-catégorie $J_0$ de $J$ dont l'ensemble des objets est $\{e\}$, et l'ensemble des flèches $\{\id_e\}\cup\Sigma_\infty$, soit filtrante). Si $\Sigma_0$ est choisie ainsi, soit $\mathcal{F}$ l'ensemble réunion de $\Fl K$ et de l'ensemble des flèches composées
\[
k \xrightarrow{\ u\ } e \xrightarrow{\ p\ } e \xrightarrow{\ v\ } k'\ , \tag*{$(*)$}
\]
avec $u,v\in\Fl J$, $p\in\Sigma_\infty$. Alors $\mathcal{F}$ est l'ensemble des flèches d'une sous-catégorie $K'$ de $J$, ayant mêmes objets que $J$, et admettant $e$ pour plus grand élément. De plus, $K'$ contient $K$ et est filtrante.}

\noindent\textsc{Démonstration}. Que $\Sigma$ soit stable par composition est clair. Pour construire $\Sigma_\infty$ ayant les propriétés énoncées, on note que, $J$ étant grande devant $\pi$, $\Sigma$ est non vide (\footnote[162]{cf. argument plus bas pour construction de $p_i$.}). On construit par récurrence une suite d'éléments
\[
p_0\ ,\ p_1\ ,\ \ldots\ ,\ p_n\ ,\ \ldots
\]
de $\Sigma$, en prenant pour $p_0$ un élément quelconque de $\Sigma$, et pour $i\geqslant 1$, prenant $p_i\in\Sigma$ tel que pour toute double-flèche $(p,p')$ du sous-monoïde $M_{i-1}$ de $\End_J(e)$ engendré par $(p_j)_{0\leqslant j\leqslant i-1}$, on ait $p_ip=p_ip'$. Un tel $p_i$ existe, car prenant l'ensemble des double-flèches de la forme précédente $(p,p')$, avec $p,p'\in M_{i-1}$, et de celles de la forme $(u,v) : k \rightrightarrows e$, avec $u,v\in\Fl J$, l'ensemble de
\msp{page 368}
toutes ces double-flèches est de cardinal $\leqslant\pi$ ($M_{i-1}$ étant dénombrable, donc de cardinal $\leqslant\pi$), donc il existe une flèche $e\longrightarrow j$ dans $J$ qui les égalise toutes ($J$ étant grande devant $\pi$). Quitte à composer avec $j\longrightarrow e$ ($e$ étant objet maximal de $J$), on trouve le $p_i : e\longrightarrow e$ cherché.

Soit alors $\Sigma_0$ le sous-ensemble multiplicatif de $\Sigma$ engendré par les $p_i$ ($i\in\mathbf{N}$), limite inductive des ensembles similaires $\Sigma_i$ engendrés par les $p_j$ pour $j\in[0,i]$. Si $p$, $p'$ sont dans $\Sigma_\infty$, il sont dans un même $\Sigma_i$, donc on aura $p_{i+1}p=p_{i+1}p'$, ce qui prouve que $\Sigma_\infty$ a la propriété filtrante exigée. % [sic] il sont -> ils sont, agreement of the verb with the plural subject

Prouvons maintenant que $\mathcal{F}$ est l'ensemble de flèches d'une sous-catégorie $K'$ de $J$, ayant $\Ob K$ comme ensemble d'objets. Il reste à établir la stabilité par composition $vu$. Mais si $u$ et $v$ sont dans $\Fl K$, ou si l'un des deux est dans $\Fl K$, l'autre non (donc de la forme $(*)$), cette stabilité est évidente. Reste à voir le cas où ni \edbr{[}$u$\edbr{]} ni $v$ sont dans $\Fl K$, donc sont tous deux de la forme $(*)$ :
\[
k \xrightarrow{\ u\ } e \xrightarrow{\ p\ } e \xrightarrow{\ v\ } k' \xrightarrow{\ u'\ } e \xrightarrow{\ p'\ } e \xrightarrow{\ v'\ } k''\ ,
\]
avec $u,v,u',v'\in\Fl K$, $p,p'\in\Sigma_\infty$. Comme $(u'v,\id_e)$ est une double-flèche de
\printedpage{140}
% head: A. Grothendieck : Dérivateurs, Ch. XVIII, Édition des Universités de Montpellier II et Paris VII
% foot: Version 17 décembre 2025
\msp{page 369}
$K$ de but $e$, on a
\[
p'u'v \;=\; p'\id_e \;=\; p'\ ,
\]
donc le composé en question s'écrit aussi
\[
k \xrightarrow{\ u\ } e \xrightarrow{\ p'p\ } e \xrightarrow{\ v'\ } k''\ ,
\]
où $p'p\in\Sigma_\infty$ puisque $\Sigma_\infty$ est stable par composition, donc il est dans $\mathcal{F}$, q.e.d.

Il est clair que $K'$ contient $K$, et ayant mêmes objets que $K$ qui admet $e$ comme plus grand élément, que $e$ est un plus grand élément de $K'$. Reste à prouver que $K'$ est filtrante. Comme $K'$ a un plus grand élément, il reste à voir que toute double-flèche $k\rightrightarrows l$ dans $K'$ peut être égalisée dans $K'$. Quitte à composer avec $l\longrightarrow e$, on peut supposer $l=e$. Or on voit que toute flèche $k\longrightarrow e$ dans $K'$ est de la forme
\[
k \xrightarrow{\ u\ } e \xrightarrow{\ p\ } e \xrightarrow{\ v\ } e\ ,
\]
avec $u,v\in\Fl K$ et $p\in\Sigma_\infty\cup\{\id_e\}$. Quitte à composer avec $p_0\in\Sigma_\infty$ fixé, cette flèche devient (puisque $p_0v=p_0\id_e=p_0$)
\[
k \xrightarrow{\ u\ } e \xrightarrow{\ p_0p\ } e\ , \qquad\quad \text{i.e. } p_0pu\ .
\]
Pour une flèche $v'p'u' : k\longrightarrow e$, on trouve $p_0p'u'$. Soit alors $q\in\Sigma_0$ tel que $q(p_0p)=q(p_0p')$, soit $r$ ($q$ existe puisque $p_0p,p_0p'\in\Sigma_\infty$). Donc $(q(p_0pu),q(p_0p'u'))=(ru,ru')$, et on a $ru=ru'$ puisque $r\in\Sigma_\infty$. Au total, $qp_0$ égalise les deux flèches $vpu$ et $v'p'u'$, q.e.d.

\msp{page 370}

\section*{4.14\quad Critères d'accessibilité pour grosses catégories}

\noindent\textbf{Théorème 4.14.1.} {\itshape Soient $\mathcal{C}$ une catégorie équivalente à une petite catégorie, $\Phi$, $\Psi$ deux petits ensembles de catégories, équivalentes à des petites catégories. Pour toute catégorie $I$, on désigne par $I_{\bullet}$ (resp. ${}_{\bullet}I$) le cône droit (resp. gauche) sur $I$, déduit de $I$ en ajoutant un objet final (resp. initial) $e=e_I$. Soit $\Lambda_1$ (resp. $\Lambda_2$) un ensemble de couples $(I,\varphi)$ (resp. $(I,\psi)$), où $I\in\Phi$ (resp. $I\in\Psi$), et où
\[
\varphi\ :\ I \longrightarrow \mathcal{C} \qquad\quad \text{(resp. }\psi\ :\ I \longrightarrow \mathcal{C}\ \ \text{)} \tag*{\textup{(4.14.1.1)}}
\]
est un foncteur. Pour toute catégorie $\mathcal{N}$ stable (\footnote[163]{Cette propriété de stabilité n'est pas nécessaire pour la définition de $\uHom^{\Lambda_2}_{\Lambda_1}(\mathcal{C},\mathcal{N})$, et il est presque sûr que le théorème est valable sans faire cette hypothèse sur $\mathcal{N}$. Cf. 4.14.16, où on voit que c'est en tous cas inutile pour $\Psi$.}), par $\varinjlim$ de type $\Phi$ et par $\varprojlim$ de type $\Psi$ on désigne par
\[
\uHom^{\Lambda_2}_{\Lambda_1}(\mathcal{C},\mathcal{N})\ \subset\ \uHom(\mathcal{C},\mathcal{N}) \tag*{\textup{(4.14.1.2)}}
\]
la sous-catégorie strictement pleine de $\uHom(\mathcal{C},\mathcal{N})$ formée des foncteurs $f:\mathcal{C}\longrightarrow\mathcal{N}$ ayant la propriété suivante :
\begin{enumerate}[leftmargin=2.5em]
\item[\textit{a)}] $\forall\ (I,\varphi)\in\Lambda_1$, le foncteur composé
\[
I_{\bullet} \xrightarrow{\ \varphi\ } \mathcal{C} \xrightarrow{\ f\ } \mathcal{N}
\]
est $I$-exact à gauche, i.e. fait de $(f\varphi)(e_I)$ une $\varinjlim\limits_{I}$ de $f\varphi|I$.
\printedpage{141}
% head: A. Grothendieck : Dérivateurs, Ch. XVIII, Édition des Universités de Montpellier II et Paris VII
% foot: Version 17 décembre 2025
\msp{page 371}
\item[\textit{b)}] Pour tout $(I,\psi)\in\Lambda_2$, le foncteur composé
\[
{}_{\bullet}I \xrightarrow{\ \psi\ } \mathcal{C} \xrightarrow{\ f\ } \mathcal{N}
\]
est $I$-exact à droite, i.e. fait de $(f\psi)(e_I)$ une limite projective de $(f\psi)|I$.
\end{enumerate}
\noindent Avec ces notations : si $\mathcal{N}$ est une catégorie accessible, il en est de même de $\uHom^{\Lambda_2}_{\Lambda_1}(\mathcal{C},\mathcal{N})$.}

\noindent\textbf{Commentaire.} On utilisera ce théorème surtout quand on a $\Lambda_2=\emptyset$ ou $\Lambda_1=\emptyset$, auquel cas on note $\uHom_{\Lambda_1}(\mathcal{C},\mathcal{N})$, resp. $\uHom^{\Lambda_2}(\mathcal{C},\mathcal{N})$, la catégorie $\uHom^{\Lambda_2}_{\Lambda_1}(\mathcal{C},\mathcal{N})$. On appliquera le théorème surtout quand $\mathcal{C}$ elle-même est stable par $\varinjlim$ de type $\Phi$ et par $\varprojlim$ de type $\Psi$, et quand $\Lambda_1$ (resp. $\Lambda_2$) est déjà formé de foncteurs $\varphi$ (resp. $\psi$) qui sont $I$-exacts à gauche (resp. à droite). Donc $\uHom^{\Lambda_2}_{\Lambda_1}(\mathcal{C},\mathcal{N})$ est alors formée des $f:\mathcal{C}\longrightarrow\mathcal{N}$ qui commutent à certaines $\varinjlim$ de type $\Phi$, et à certaines $\varprojlim$ de type $\Psi$. En particulier, on trouve :

\msp{page 372}

\noindent\textbf{Corollaire 4.14.2.} {\itshape Soient $\mathcal{C}$ une catégorie équivalente à une petite catégorie, $\mathcal{N}$ une $\mathfrak{U}$-catégorie, $\Phi$ et $\Psi$ deux petits ensembles de catégories $I$, équivalentes à des petites catégories. On suppose $\mathcal{C}$ et $\mathcal{N}$ stables par $\varinjlim$ de type $\Phi$, par $\varprojlim$ de type $\Psi$. Soit $\uHom^{\Psi}_{\Phi}(\mathcal{C},\mathcal{N})$ la sous-catégorie strictement pleine de $\uHom(\mathcal{C},\mathcal{N})$ formée des $f:\mathcal{C}\longrightarrow\mathcal{N}$ qui commutent aux $\varinjlim$ de type $\Phi$ et aux $\varprojlim$ de type $\Psi$. Alors, si $\mathcal{N}$ est accessible, il en est de même de $\uHom^{\Psi}_{\Phi}(\mathcal{C},\mathcal{N})$.}

Voici des cas particuliers remarquables de ce corollaire :

\noindent\textbf{Corollaire 4.14.3.} {\itshape Soient $\mathcal{M}$, $\mathcal{N}$ deux catégories accessibles, et $\pi_0$ un cardinal infini, tel que $\mathcal{M}$ satisfasse à $L_{\pi_0}$. Soit $\uHom_{\pi_0}(\mathcal{M},\mathcal{N})$ la sous-catégorie strictement pleine de $\uHom(\mathcal{M},\mathcal{N})$ formée des foncteurs $\pi_0$-accessibles. Alors $\uHom_{\pi_0}(\mathcal{M},\mathcal{N})$ est accessible.}

\msp{page 373}

(Le cas où $\pi_0$ est utile pour $\mathcal{M}$ a déjà été traité.)

\noindent\textsc{Démonstration}. Dans le cas général, soit $\pi>\pi_0$ un cardinal utile pour $\mathcal{M}$, de sorte qu'on a
\[
\mathcal{M} \xleftarrow{\ \sim\ } \Ind_\pi(\mathcal{C}) \qquad\quad \text{avec } \mathcal{C}=\mathcal{M}_\pi\ .
\]
On peut supposer $\mathcal{M}=\Ind_\pi(\mathcal{C})$. Le fait que $\mathcal{M}$ satisfasse à $\mathrm{L}_{\pi_0}$ implique que $\mathcal{C}$ est stable par $\varinjlim$ grandes devant $\pi_0$, et de cardinal $\leqslant\pi$. Soit $\Phi$ l'ensemble des catégories de cette nature. $\Phi$ n'est pas petit, mais il existe $\Phi_0\subset\Phi$ petit, tel que toute $I\in\Phi$ soit isomorphe à une $I_0\in\Phi_0$, ce qui suffira pour la validité de 4.14.1 et 4.14.2 (où on fera ici $\Psi=\emptyset$). On peut considérer $\uHom_{!\,\pi_0}(\mathcal{M},\mathcal{N})$ comme une sous-catégorie strictement pleine de $\uHom_{!\,\pi}(\mathcal{M},\mathcal{N})\xrightarrow{\ \sim\ }\uHom(\mathcal{C},\mathcal{N})$ (cf. 4.12.3, p. 275). Pour un foncteur $f:\mathcal{M}\longrightarrow\mathcal{N}$ qui est $\pi$-accessible, dire qu'il est même $\pi_0$-accessible revient à dire que sa restriction à $\mathcal{C}$, i.e. son image dans $\uHom(\mathcal{C},\mathcal{N})$, commute aux $\varinjlim$ de type $\Phi$. (Cf. 4.13.20, page 341 (\footnote[164]{La référence n'est pas tout à fait adéquate. Il faudrait réunir dans un paragraphe à part tous les énoncés de ce type, dont certains utilisent le théorème de représentation indicielle.}), où on prend $\alpha=\alpha_\pi$, notant que $J_{\alpha_\pi}$ grande devant $\pi_0$ puisque $\pi>\pi_0$).

\msp{page 374}

Donc on trouve
\[
\uHom_{!\,\pi_0}(\mathcal{M},\mathcal{N}) \xrightarrow{\ \sim\ } \uHom_{\Phi}(\mathcal{C},\mathcal{N})\ ,
\]
\printedpage{142}
% head: A. Grothendieck : Dérivateurs, Ch. XVIII, Édition des Universités de Montpellier II et Paris VII
% foot: Version 17 décembre 2025
et le corollaire 4.14.2 donne la conclusion voulue.

\noindent\textbf{Corollaire 4.14.4.} {\itshape Soient $\mathcal{M}$, $\mathcal{N}$ deux catégories accessibles, $\Phi$ et $\Psi$ deux ensembles des catégories comme dans 4.14.2, on suppose que $\mathcal{M}$ et $\mathcal{N}$ satisfassent $L_{\pi_0}$, et \edbr{[soient]} stables par $\varinjlim$ de type $\Phi$ et par $\varprojlim$ de type $\Psi$. Soit $\uHom^{\Psi}_{\Phi,\pi_0}(\mathcal{M},\mathcal{N})$ la sous-catégorie strictement pleine de $\uHom(\mathcal{M},\mathcal{N})$ formée des foncteurs $\pi_0$-accessibles qui commutent aux $\varinjlim$ de type $\Phi$, et aux $\varprojlim$ de type $\Psi$. Alors $\uHom^{\Psi}_{\Phi,\pi_0}(\mathcal{M},\mathcal{N})$ est accessible.}

\noindent\textsc{Démonstration}. Cette fois, on choisit $\pi>\pi_0$ tel qu'on ait $\cardred I\leqslant\pi$ pour tout $I$ dans $\Phi$ ou dans $\Psi$. La catégorie envisagée s'interprète alors (posant encore $\mathcal{C}=\mathcal{M}_\pi$) comme
\msp{page 375}
\[
\uHom^{\Psi}_{\Phi,\pi_0}(\mathcal{M},\mathcal{N}) \xrightarrow{\ \sim\ } \uHom^{\Psi}_{\Phi'}(\mathcal{C},\mathcal{N})\ ,
\]
où $\Phi'$ est formée des catégories qui sont soit dans $\Phi$, ou qui sont grandes devant $\pi_0$ et de cardinal $\leqslant\pi$. Il suffit donc d'appliquer encore 4.14.2.

\noindent\textbf{Corollaire 4.14.5.} {\itshape Considérons les trois ensembles suivants :
\begin{enumerate}[leftmargin=2.5em]
\item[\textit{a)}] $\Phi_1$ est formée de toutes les catégories équivalentes à des petites catégories.
\item[\textit{b)}] $\Phi_2$ comme a), mais en restreignant aux catégories filtrantes.
\item[\textit{c)}] $\Phi_3$ formé des catégories discrètes, et de celles qui sont grandes devant $\pi_0$ ($\pi_0$ un cardinal infini donné).
\end{enumerate}
\noindent Soit $\Phi$ l'un des ensembles de catégories $\Phi_1$, $\Phi_2$, $\Phi_3$, $\Phi_2\cup\Phi_3$ ou la réunion d'un de ceux-ci et d'un petit ensemble $\Phi_0$ de catégories essentiellement petites (\footnote[165]{On pourrait faire ici un énoncé qui englobe 4.14.4 en rajoutant un
\begin{enumerate}[leftmargin=2.5em]
\item[d)] $\Phi_4$ formé des catégories grandes devant $\pi_0$,
\end{enumerate}
\noindent et prendre $\Phi$ de la forme de $\Phi_1$, $\Phi_2$, $\Phi_4$, $\Phi_3$, $\Phi_3\cup\Phi_2$, $\Phi_3\cup\Phi_4$, ou la réunion de ceux-ci et d'un $\Phi_0$.}). Soit $\Psi$ un petit ensemble de catégories. Soient $\mathcal{M}$, $\mathcal{N}$ des catégories accessibles, stables par $\varinjlim$ de type $\Phi$, et par $\varprojlim$ de type $\Psi$. Alors $\uHom^{\Psi}_{\Phi}(\mathcal{M},\mathcal{N})$ est accessible.}

\msp{page 376}

\noindent\textsc{Démonstration}. Soit $\pi'_0\geqslant\pi_0$ un cardinal utile pour $\mathcal{M}$, $\pi$ un cardinal ayant les propriétés suivantes.
\begin{enumerate}[leftmargin=2.5em]
\item[a)] $\pi\geqslant\sup(\pi'_0,\cardred\mathcal{M}_{\pi'_0})$, $\pi>\pi_0$.
\item[b)] $\pi^{\pi'_0}=\pi$.
\item[c)] $I\in\Phi_0\cup\Psi\Longrightarrow\cardred I\leqslant\pi$.
\end{enumerate}
\noindent Soit $\mathcal{C}=\mathcal{M}_\pi$, d'où $\mathcal{M}\simeq\Ind_\pi(\mathcal{C})$. Dans tous les cas envisagés, $\Phi$ contient l'ensemble des catégories essentiellement petites qui sont grandes devant $\pi_0$, a fortiori celles qui sont grandes devant $\pi$, donc la catégorie $\uHom^{\Psi}_{\Phi}$ s'interprète comme une sous-catégorie strictement pleine de
\[
\uHom(\mathcal{C},\mathcal{N})\ .
\]
Si d'autre part $\Phi_\pi$ désigne l'ensemble des $I\in\Phi$ telles que $\card I\leqslant\pi$, on trouve
\[
\uHom^{\Psi}_{\Phi}(\mathcal{M},\mathcal{N})\ \simeq\ \uHom^{\Psi}_{\Phi_\pi}(\mathcal{C},\mathcal{N})\ ,
\]
et on applique 4.14.2.

\printedpage{143}
% head: A. Grothendieck : Dérivateurs, Ch. XVIII, Édition des Universités de Montpellier II et Paris VII
% foot: Version 17 décembre 2025
\msp{page 377}

En particulier, on trouve :

\noindent\textbf{Corollaire 4.14.6.} {\itshape Soient $\mathcal{M}$, $\mathcal{N}$ deux catégories accessibles, supposons $\mathcal{M}$ stable par petites $\varinjlim$ (i.e. $\mathcal{M}$ est une catégorie paratopos). Alors la catégorie $\uHom_{!}(\mathcal{M},\mathcal{N})$ des foncteurs $f:\mathcal{M}\longrightarrow\mathcal{N}$ qui commutent aux petites $\varinjlim$ est accessible. Si $\mathcal{N}$ est elle aussi un paratopos, $\uHom_{!}(\mathcal{M},\mathcal{N})$ est un paratopos.}

\noindent\textsc{Démonstration}. La première assertion est le cas particulier de 4.14.5 relatif au cas où $\Phi=\Phi_0$, $\Psi=\emptyset$. La deuxième assertion en résulte, puisque $\uHom_{!}(\mathcal{M},\mathcal{N})$ est visiblement stable par petites $\varinjlim$ , donc c'est un paratopos étant accessible (4.8.9, page 207).

\noindent\textbf{Corollaire 4.14.7.} {\itshape Soient $\Phi$, $\Psi$, $\mathcal{M}$ comme dans 4.14.5. Alors pour toute catégorie accessible $\mathcal{N}$, la catégorie
\msp{page 378}
\[
\uHom^{\Phi}_{\Psi}(\mathcal{M}^{\circ},\mathcal{N})\ \simeq\ (\uHom^{\Psi}_{\Phi}(\mathcal{M},\mathcal{N}^{\circ}))^{\circ}
\]
est accessible.}

\noindent\textsc{Démonstration}. On procède comme pour 4.14.5, pour expliciter
\[
\uHom^{\Psi}_{\Phi}(\mathcal{M},\mathcal{N}^{\circ}) \xrightarrow{\ \sim\ } \uHom^{\Psi}_{\Phi_\pi}(\mathcal{C},\mathcal{N}^{\circ})\ ,
\]
d'où
\[
\uHom^{\Phi}_{\Psi}(\mathcal{M}^{\circ},\mathcal{N}) \xrightarrow{\ \sim\ } \uHom^{\Phi_\pi}_{\Psi}(\mathcal{C}^{\circ},\mathcal{N})\ ,
\]
et on applique 4.14.2 à $\mathcal{C}^{\circ}$, $\Psi$, $\Phi_\pi$.

\noindent\textbf{Corollaire 4.14.8.} {\itshape Soient $\mathcal{M}$ et $\mathcal{N}$ deux catégories paratopos. Alors
\[
\mathcal{M}\boxtimes\mathcal{N}\ \simeq\ \uHom^{!}(\mathcal{M}^{\circ},\mathcal{N}) \qquad\quad \text{(cf. 4.2 et 4.3)}
\]
est un paratopos.}

En effet, en vertu de 4.14.7 (avec $\Phi=\Phi_\pi$ [?], $\Psi=\emptyset$), c'est une catégorie accessible.

\msp{page 379}

D'autre part, il est immédiat qu'elle est stable par petites $\varprojlim$ , $\mathcal{N}$ l'étant. Donc on conclut encore par 4.8.9 (page 207) que c'est une catégorie paratopos.

\noindent\textbf{Corollaire 4.14.9.} {\itshape Soient $\mathcal{M}$, $\mathcal{N}$ deux catégories paratopos, et $\uHom^{!}(\mathcal{M},\mathcal{N})$ la catégorie des foncteurs $\mathcal{M}\longrightarrow\mathcal{N}$ qui commutent aux petites $\varprojlim$ . On a alors une équivalence de catégories
\[
\begin{aligned}
\uHom^{!}(\mathcal{M},\mathcal{N}) &\ \xrightarrow{\ \sim\ }\ \uHom_{!}(\mathcal{N},\mathcal{M})^{\circ}\\
f &\ \longmapsto\ \text{adjoint à gauche de } f\ ,
\end{aligned}
\]
\noindent\textup{\textbf{!}}\ \textup{\textbf{!}} et par suite la catégorie \textup{\textbf{opposée}} à $\uHom^{!}(\mathcal{M},\mathcal{N})$ est accessible.} % Verified against PDF leaf 25: a nonverbal marginal curved mark; see the correction register.

Ceci résulte du

\noindent\textbf{Lemme 4.14.10} (pour mémoire). {\itshape Soient $\mathcal{M}$ \edbr{[une]} catégorie pseudo-topos, $f:\mathcal{M}\longrightarrow\mathcal{N}$ un foncteur. Pour que $f$ admette un adjoint à gauche (resp. à droite), il faut et il suffit que $f$ commute aux petites $\varprojlim$ (resp. aux petites $\varinjlim$ ).}

% ---- from der-ch18_p146-150.tex
\printedpage{144}
% head: A. Grothendieck : Dérivateurs, Ch. XVIII, Édition des Universités de Montpellier II et Paris VII
% foot: Version 17 décembre 2025

\msp{page 380}
Ceci équivaut au cas particulier :

\textbf{Corollaire 4.14.11.} \textit{Soit $\mathcal{M}$ un paratopos, et $f$ un foncteur covariant (resp. contravariant) de $\mathcal{M}$ dans \textup{Ens}. Pour que $f$ soit représentable, il faut et il suffit que $f$ commute aux $\varprojlim$ .}

Ça a été vu -- le cas des contrafoncteurs étant le plus évident (cf. 4.1 -- il suffit que $\mathcal{M}$ soit un paratopos, et même que $\mathcal{M}$ admette une petite sous-catégorie pleine génératrice, cf. SGA 4 I 8.12.7 (pour l'énoncé dual). Dans les cas d'un foncteur covariant, la difficulté provient du fait que $\mathcal{M}$ n'admet pas forcément de petite sous-catégorie cogénératrice. Il faut utiliser l'existence d'une filtration cardinale (généralisée), cf. loc. cit. 8.3.12.) % [sic] Dans les cas -> Dans le cas, singular "un foncteur covariant"

Après tous ces beaux corollaires, on est mûr pour s'attaquer à la démonstration de 4.14.1 !

\msp{page 381}
\textsc{Démonstration de} 4.14.1. On sait que $\mathcal{M} = \uHom(\mathcal{C},\mathcal{N})$ est accessible (4.9.5, p. 220), et $\uHom^{\Lambda_2}_{\Lambda_1}(\mathcal{C},\mathcal{N})$ est une sous-catégorie strictement pleine de cette catégorie. Il suffit de prouver que c'est une sous-catégorie accessible dans $\mathcal{M}$ (4.13.1, p. 281). Pour tout $\xi = (I,\varphi) \in \Lambda_1$, on a un foncteur
\[
L_\xi \ : \ \mathcal{M} \ = \ \uHom(\mathcal{C},\mathcal{N}) \longrightarrow \uFl(\mathcal{N}) \ ,
\]
donné par
\[
L_\xi(f) \ = \ \text{flèche canonique } \varinjlim_{i\in I} f\varphi(i) \longrightarrow f\varphi(e_I) \ .
\]
Je dis que le foncteur $L_\xi$ est \textit{accessible.} En effet, soit $\pi_0$ un cardinal tel que $\mathcal{N}$ satisfasse à $\mathrm{L}_{\pi_0}$, il en est donc de même de $\mathcal{M} = \uHom(\mathcal{C},\mathcal{N})$ et de $\uFl(\mathcal{N}) \defeq \uHom(\boldsymbol{\Delta}_1,\mathcal{N})$, les limites en question se calculent argument par argument. Je dis que $L_\xi$ est $\pi_0$-accessible, i.e. commute aux $\varinjlim$ grandes devant $\pi_0$. Il suffit de voir qu'il en est ainsi des foncteurs composés avec les foncteurs
\[
\begin{tikzcd}
\uFl(\mathcal{N}) \arrow[r, shift left, "\sigma"] \arrow[r, shift right, "\rho"'] & \mathcal{N}
\end{tikzcd}
\]
``source'' et ``but'', i.e. des foncteurs
\[
f \longmapsto \varinjlim_{i\in I} f\varphi(i) \qquad \text{et} \qquad f \longmapsto f\varphi(e_I) \ .
\]

\msp{page 382}
Or le premier foncteur, $f \longmapsto \varinjlim\limits_{i\in I} f\varphi(i)$, est le composé de
\[
f \longmapsto f\circ\varphi_I \ : \ \uHom(\mathcal{C},\mathcal{N}) \longrightarrow \uHom(I,\mathcal{N})
\]
(lequel commute aux $\varinjlim$ qui sont représentables dans $\mathcal{N}$, et en particulier aux $\varinjlim$ grandes devant $\pi_0$) et du foncteur
\[
\varinjlim_{I} \ : \ \uHom(I,N) \longrightarrow \mathcal{N} \ , % [sic] \uHom(I,N) -> \uHom(I,\mathcal{N}), italic N printed instead of the script N used everywhere else
\]
qui lui aussi commute aux limites inductives représentables dans $\mathcal{N}$, et en particulier à celles grandes devant $\pi_0$. Quant au foncteur $f \longmapsto (f\circ\varphi)(e_I) = f(\varphi(e_I))$, il est clair qu'il commute également aux $\varinjlim$ représentables dans $\mathcal{N}$, donc il est $\pi_0$-accessible.

\printedpage{145}
% head: A. Grothendieck : Dérivateurs, Ch. XVIII, Édition des Universités de Montpellier II et Paris VII
% foot: Version 17 décembre 2025

Ceci posé, notons que
\[
\uHom_{\Lambda_1}(\mathcal{C},\mathcal{N}) \ = \ \bigcap_{\xi\in\Lambda_1} \uHom_{\{\xi\}}(\mathcal{C},\mathcal{N}) \ ,
\]
et que $\uHom_\xi(\mathcal{C},\mathcal{N})$ est l'image inverse, par le foncteur accessible $L_\xi$ précédent, de la sous-catégorie strictement pleine $\uIs(\mathcal{N})$ de $\uFl(\mathcal{N})$ formée des \textit{isomorphismes} dans $\mathcal{N}$. Si on prouve que cette dernière est accessible dans $\uFl(\mathcal{N})$, il s'ensuit que
\msp{page 383}
$\uHom_\xi(\mathcal{C},\mathcal{N})$ est accessible dans $\mathcal{M}$ (4.3.31, page 350). Donc, $\Phi$ étant petite, il s'ensuit que $\uHom_{\Lambda_1}(\mathcal{C},\mathcal{N})$ est accessible dans $\mathcal{M}$ (4.13.7, page 291).

On prouve de la même façon que $\uHom^{\Lambda_2}(\mathcal{C},\mathcal{N})$ est accessible dans $\mathcal{M} = \mathrm{Hom}(\mathcal{C},\mathcal{N})$. Ici, on doit considérer les foncteurs % [sic] \mathrm{Hom}(\mathcal{C},\mathcal{N}) -> \uHom(\mathcal{C},\mathcal{N}), printed without underline, elsewhere underlined
\[
L'_\xi \ : \ \mathcal{M} = \uHom(\mathcal{C},\mathcal{N}) \longrightarrow \uFl(\mathcal{N})
\]
définis par
\[
L'_\xi \ = \ \text{flèche canonique } f\varphi(e'_I) \longrightarrow \varprojlim_{i\in I} f\varphi(i) \ .
\]
Pour prouver que ce foncteur est accessible, on procède comme tantôt, en prenant cependant $\pi_0$ assez grand pour qu'on ait
\[
\left\{
\begin{array}{l}
I \in \Psi \ \Longrightarrow \ \cardred I \leqslant \pi_0. \\[1ex]
\mathcal{N}_{\pi_0} \text{ est génératrice dans } \mathcal{N}.
\end{array}
\right.
\]
Je dis qu'alors $L'_\xi$ est $\pi_0$-accessible : en effet, en procédant comme tantôt pour $L_\xi$, on est ramené à voir que
\[
\varinjlim_{I} \uHom(I,\mathcal{N}) \longrightarrow \mathcal{N}
\]
est $\pi$-accessible. Or cela est contenu dans SGA 4 I 9.8.

\msp{page 384}
Ainsi, on a prouvé que les deux sous-catégories strictement pleines
\[
\uHom_{\Lambda_1}(\mathcal{C},\mathcal{N}) \ , \qquad \uHom^{\Lambda_2}(\mathcal{C},\mathcal{N})
\]
de $\mathcal{M} \defeq \uHom(\mathcal{C},\mathcal{N})$ sont accessibles dans $\mathcal{M}$. Il en est donc de même de leur intersection (4.13.7)
\[
\uHom^{\Lambda_2}_{\Lambda_1}(\mathcal{C},\mathcal{N}) \ = \ \uHom_{\Lambda_1}(\mathcal{C},\mathcal{N}) \cap \uHom^{\Lambda_2}(\mathcal{C},\mathcal{N}) \ ,
\]
q.e.d.

Il reste à prouver le

\textbf{Lemme 4.14.12.} \textit{Soit $\mathcal{N}$ une catégorie accessible. Alors la catégorie $\uFl(\mathcal{N}) \defeq \uHom(\boldsymbol{\Delta}_1,\mathcal{N})$ est accessible (4.9.5), et la sous-catégorie strictement pleine $\uIs(\mathcal{N})$ formée des flèches inversibles dans $\mathcal{N}$ est accessible dans $\uFl(\mathcal{N})$. (Il en est de même de la sous-catégorie strictement pleine $\uMon(\mathcal{N})$ formée des monomorphismes dans $\mathcal{N}$.)}

\textsc{Démonstration}. Soit $\mathcal{C}$ une petite sous-catégorie strictement génératrice de $\mathcal{N}$, et considérons la petite famille des foncteurs
\[
h_x \ = \ (y \longmapsto \mathrm{Hom}(x,y)) \ : \ \mathcal{N} \longrightarrow \Ens \ ,
\]
\printedpage{146}
% head: A. Grothendieck : Dérivateurs, Ch. XVIII, Édition des Universités de Montpellier II et Paris VII
% foot: Version 17 décembre 2025
\msp{page 385}
d'où des foncteurs
\[
\uFl(h_\xi) \ : \ \uFl(\mathcal{N}) \longrightarrow \uFl(\Ens) \ . % [sic] h_\xi -> h_x, the functors defined just before are h_x
\]
Pour qu'une flèche $u : y \longrightarrow z$ dans $\mathcal{N}$ soit isomorphe (resp. monomorphe), il faut et il suffit que pour tout $x \in \Ob\mathcal{C}$, $h_x(u)$ le soit. Ainsi, $\uIs(\mathcal{N})$ (resp. $\uMon(\mathcal{N})$) apparaît comme l'intersection des images inverses, par les foncteurs $\uFl(h_\xi)$ ci-dessus, de la sous-catégorie $\uIs(\Ens)$ (resp. $\uMon(\Ens)$) de $\uFl(\Ens)$. Cela nous ramène au cas où $\mathcal{N} = (\Ens)$. Mais si $\mathcal{C}$ désigne la sous-catégorie strictement pleine de Ens formée des ensembles \textit{finis}, ceux-ci forment une famille strictement génératrice formée d'objets de \textit{présentation finie}, et comme Ens est stable par petites $\varinjlim$ , donc par $\varinjlim$ filtrantes, il s'ensuit qu'on a % [sic] h_\xi -> h_x, the functors defined just before are h_x
\[
\Ind(\mathcal{C}) \longrightarrow (\Ens) \ , \tag*{$(*)$}
\]
d'où aussitôt (SGA 4 I 8.8.5, p. 104)

\msp{page 386}
\[
\Ind(\uFl(\mathcal{C})) \xrightarrow{\ \sim\ } \uFl(\Ens) \ . \tag*{$(*)$}
\]
Donc le résultat d'accessibilité voulu résulte des deux lemmes suivants :

\textbf{Lemme 4.14.13.} \textit{L'équivalence $(*)$ précédente induit des équivalences}
\[
\begin{aligned}
\Ind(\uIs(\mathcal{C})) \quad &\xrightarrow{\ \sim\ } \quad \uIs(\Ens) \\
\Ind(\uMon(\mathcal{C})) \quad &\xrightarrow{\ \sim\ } \quad \uMon(\Ens)
\end{aligned}
\]
\textit{(où $\mathcal{C} = (\mathrm{Ensf})$ est la sous-catégorie pleine de \textup{(Ens)} formée des ensembles \textup{finis}.)}

\textbf{Lemme 4.14.14} (\footnote[166]{\textbf{N.B.} À équivalence près, on trouve les catégories accessibles $\mathcal{M}$ stables par petites $\varinjlim$ filtrantes, et telles que $\mathcal{M}_{\mathrm{pf}}$ (sous-catégorie des objets de présentation \textit{finie}) soit strictement génératrice.}). \textit{Soit $\mathcal{C}$ une catégorie équivalente à une petite catégorie. Alors :}
\begin{itemize}[nosep]
\item[\textit{a)}] \textit{$\mathcal{M} \defeq \Ind(\mathcal{C})$ est une catégorie accessible.}
\item[\textit{b)}] \textit{Pour toute sous-catégorie pleine $\mathcal{D}$ dans $\mathcal{C}$, l'image essentielle de $\Ind\mathcal{D}$ dans $\mathcal{M} = \Ind(\mathcal{C})$ est une sous-catégorie L-accessible dans $\mathcal{M}$.}
\end{itemize}
\textsc{Démonstration} des deux lemmes immédiate et ``laissée au lecteur''. Le b) dans 4.14.14 provient du fait que le foncteur $\Ind(\mathcal{D}) \longrightarrow \Ind(\mathcal{C})$ provenant d'un foncteur $\mathcal{D} \longrightarrow \mathcal{C}$ (ici un
\msp{page 387}
foncteur d'inclusion) est un L-foncteur (i.e. commute aux petites $\varinjlim$ filtrantes), et a fortiori est accessible (et même $\aleph_0$-accessible).

\textbf{4.14.15.} On ne peut s'empêcher de se demander si pour $\mathcal{N}$ accessible, la sous-catégorie strictement pleine $\uEp(\mathcal{N})$ dans $\uFl(\mathcal{N})$, formée des épimorphismes dans $\mathcal{N}$, est accessible dans $\mathcal{N}$. Il est immédiat que si $\mathcal{N}$ satisfait à $\mathrm{L}_{\pi_0}$, alors $\uEp(\mathcal{N})$ est stable dans $\uFl(\mathcal{N})$ par les $\varinjlim$ de type $\mathrm{L}_{\pi_0}$, a fortiori par facteurs directs. Elle aurait donc dû mal à ne pas être accessible dans $\uFl(\mathcal{N})$. Notons quand même : % [sic] dû mal -> du mal, "avoir du mal à"

\textbf{Corollaire 4.14.15.1} (de 4.14.13). \textit{Soit $\mathcal{N}$ une catégorie accessible stable par sommes amalgamées $y \amalg_x z$} (\footnote[167]{Je suis presque sûr que cette hypothèse est superflue. Comparer page 390.}). \textit{Alors $\uEp(\mathcal{N})$ est une sous-catégorie accessible dans $\uFl(\mathcal{N})$.}

\printedpage{147}
% head: A. Grothendieck : Dérivateurs, Ch. XVIII, Édition des Universités de Montpellier II et Paris VII
% foot: Version 17 décembre 2025

\msp{page 388}
Considérons en effet le foncteur
\[
\begin{aligned}
\gamma \ : \ \uFl(\mathcal{N}) \quad &\longrightarrow \quad \uFl(\mathcal{N}) \\
(x \longrightarrow y) \quad &\longmapsto \quad (y \amalg_x y \xrightarrow{\ \ \text{codiag}\ \ } y) \ .
\end{aligned}
\]
Ce foncteur commute aux types de $\varinjlim$ qui sont représentables dans $\mathcal{N}$, en particulier, il est accessible. D'autre part, $\uEp(\mathcal{N})$ est l'image inverse par $\gamma$ de la sous-catégorie strictement pleine $\uIs(\mathcal{N})$ de $\uFl(\mathcal{N})$. Donc le corollaire résulte du fait que par 4.14.13 cette dernière est accessible dans $\uFl(\mathcal{N})$.

\textbf{4.14.16.} Je présume que 4.14.1 reste valable sans supposer $\mathcal{N}$ stable par les $\varinjlim$ de type $\Phi$ et par les $\varprojlim$ de type $\Psi$. Si on a une essentiellement petite catégorie $I$, la question est de savoir si la sous-catégorie strictement pleine $\uHom_{\square}(I_\bullet,\mathcal{N})$ de $\uHom(I_\bullet,\mathcal{N})$, formée des foncteurs satisfaisant
\msp{page 389}
la condition a) de 4.14.1, est accessible dans $\uHom(I_\bullet,\mathcal{N})$, et de même pour $\uHom^{\square}({}_\bullet I,\mathcal{N})$. Pour ce deuxième cas, ça a l'air OK, car si $\mathcal{C}$ est une petite sous-catégorie strictement génératrice de $\mathcal{N}$, et si pour tout $x \in \Ob$ \texttt{[}$\mathcal{C}$\texttt{]} on considère encore le foncteur
\[
h_x \ : \ \mathcal{N} \longrightarrow \Ens \ ,
\]
d'où
\[
\bar{h}_x \ : \ \uHom({}_\bullet I,\mathcal{N}) \longrightarrow \uHom({}_\bullet I,\Ens) \ ,
\]
on voit que $\uHom^{\square}({}_\bullet I,\mathcal{N})$ est l'intersection des images inverses des sous-catégories $\uHom^{\square}({}_\bullet I,\Ens)$ de $\uHom({}_\bullet I,\Ens)$. Donc ces dernières étant accessibles dans $\uHom({}_\bullet I,\Ens)$, par 4.14.1 ((Ens) étant stable par petites $\varprojlim$ ), il s'ensuit bien que $\uHom^{\square}({}_\bullet I,\mathcal{N})$ l'est dans $\uHom({}_\bullet I,\mathcal{N})$, q.e.d.

Le cas de $\uHom_{\square}(I_\bullet,\mathcal{N})$ dans $\uHom(I_\bullet,\mathcal{N})$ semble plus délicat, dans le cas où
\msp{page 390}
\[
I = 
\begin{tikzcd}[sep=small]
 & {} \\
\bullet \arrow[ur] \arrow[dr] & \\
 & {}
\end{tikzcd}
\quad , \qquad \text{donc} \qquad
I = 
\begin{tikzcd}[sep=small]
 & {} \arrow[dr] & \\
\bullet \arrow[ur] \arrow[dr] & & {} \\
 & {} \arrow[ur] &
\end{tikzcd}
\quad ,
\]
si l'énoncé envisagé est vrai, on en conclura que $\uEp(\mathcal{N})$ est accessible dans $\uFl(\mathcal{N})$, pour toute catégorie accessible, car c'est l'image inverse de $\uHom_{\square}(I_\bullet,\mathcal{N})$ par le foncteur canonique (accessible)
\[
\begin{aligned}
\uFl(\mathcal{N}) \quad &\longrightarrow \quad \uHom(I_\bullet,\mathcal{N}) \\
(x \xrightarrow{\ u\ } y) \quad &\longmapsto \quad
\left(
\begin{tikzcd}[sep=small, ampersand replacement=\&]
 \& y \arrow[dr, "\mathrm{id}_y"] \& \\
x \arrow[ur, "u"] \arrow[dr, "u"'] \& \& y \\
 \& y \arrow[ur, "\mathrm{id}_y"'] \&
\end{tikzcd}
\right) \ .
\end{aligned}
\]

\printedpage{148}
% head: A. Grothendieck : Dérivateurs, Ch. XVIII, Édition des Universités de Montpellier II et Paris VII
% foot: Version 17 décembre 2025

C'est dire que la généralisation envisagée de 4.14.1 n'est pas oiseuse, mais est sans doute utile.

\msp{page 391}
\textbf{Théorème 4.14.17} (\footnote[168]{Cf. généralisation dans \texttt{[le]} théorème page 422, th. 4.14.26.}). \textit{Soit $p : \mathcal{E} \longrightarrow \mathcal{C}$ un foncteur cofibrant ou fibrant, avec $\mathcal{C}$ essentiellement petite (i.e. équivalente à une petite catégorie), et tel que les catégories fibres soient accessibles, et les foncteurs cochangement de base (cas cofibrant) ou de changement de base (cas fibrant) associés aux flèches $u$ de $\mathcal{C}$ soient accessibles. Alors la catégorie $\uGamma(\mathcal{E}/\mathcal{C})$ des sections de $\mathcal{E}$ sur $\mathcal{C}$ est accessible, et pour tout $s \in \Ob\mathcal{C}$, le foncteur ``évaluation en $s$'' $\varepsilon_s : \uGamma(\mathcal{E}/\mathcal{C}) \longrightarrow \mathcal{E}_s$ est accessible.}

\textbf{Corollaire 4.14.18.} \textit{Soit $\pi_0$ un cardinal infini tel que les catégories fibres satisfassent la condition $L_{\pi_0}$, et que les foncteurs cochangement de base (resp. changement de base) entre icelles satisfassent également $L_{\pi_0}$. (Un tel $\pi_0$ existe.)}
\begin{itemize}[nosep]
\item[\circled{a}] \textit{Alors $\mathcal{M} = \uGamma(\mathcal{E}/\mathcal{C})$ satisfait la condition $\underline{\mathrm{L}}_{\pi_0}$, et les foncteurs $\varepsilon_s$ aussi, i.e. les $\varinjlim$ grandes devant $\pi_0$ dans $\uGamma(\mathcal{E}/\mathcal{C})$ se calculent argument par argument.}
\item[\circled{b}] \textit{Soit un cardinal $\pi \geqslant \sup(\pi_0,\cardred\mathcal{C})$. Pour qu'un objet $f$ de $\mathcal{M}$ appartienne à $\mathcal{M}_\pi$, il \textup{suffit} que $\forall\ s \in \Ob C$, $f(s) \defeq \varepsilon_s(x)$ \textup{\texttt{[plutôt} $\varepsilon_s(f)$\texttt{]}} soit dans $\mathcal{E}_{s\,\pi}$. (Comparer 4.5.8, page 119.)} % [sic] \Ob C -> \Ob\mathcal{C}, italic C printed instead of script C
\item[\circled{c}] \textit{Supposons $pi_0$ choisi de façon} % [sic] pi_0 -> \pi_0, the letters "pi" printed instead of the Greek letter
\msp{page 392}

\textit{à satisfaire les conditions supplémentaires suivantes : $\pi_0$ est utile pour $\mathcal{M} = \uGamma(\mathcal{E}/\mathcal{C})$, et pour chacune des catégories fibres $\mathcal{E}_s$ ($s \in \Ob\mathcal{C}$) (un tel $\pi_0$ existe). Soit}
\[
\pi_1 \ = \ \sup(\pi_0, \cardred\mathcal{C}, \cardred\mathcal{M}_{\pi_0}) \ .
\]
\textit{Alors pour tout cardinal $\pi \geqslant \pi_1$ et tel que $\pi^{\pi_0} = \pi$, on a}
\[
\mathcal{M}_\pi \ = \ \left\{
\begin{array}{l}
\textit{sous-catégorie strictement pleine de $\mathcal{M}$ formée des $f \in$} \\
\textit{$\Ob\mathcal{M}$ telles que $f(s) \in \Ob\mathcal{E}_{s\,\pi}$ $\forall\ s \in S$ \textup{\texttt{[plutôt}} $s \in$} \\
\textit{$\Ob C$\textup{\texttt{]}}.} % [sic] \Ob C -> \Ob\mathcal{C}, italic C printed instead of script C
\end{array}
\right.
\]
\textit{Si de plus $\pi$ est choisi tel qu'on ait}
\[
\pi \ \geqslant \ \sup_{s\in\Ob\mathcal{C}} \cardred\mathcal{E}_{s\,\pi_0} \ ,
\]
\textit{alors la sous-catégorie strictement pleine $\mathcal{E}(\pi)$ de $\mathcal{E}$ telle que $\Ob\mathcal{E}(\pi) = \bigcup_{s\in\Ob\mathcal{C}} \Ob\mathcal{E}_{s\,\pi}$ est une sous-catégorie $\mathcal{C}$-cofibrée (resp. $\mathcal{C}$-fibrée) de $\mathcal{E}$, et on a}
\[
\uGamma(\mathcal{E}/\mathcal{C})_\pi \xleftarrow{\ \sim\ } \uGamma(\mathcal{E}(\pi)/C), % [sic] C -> \mathcal{C}, italic C printed instead of script C
\]
\textit{donc ($\pi$ étant utile pour $\mathcal{E}$) on a une équivalence de catégories}
\[
\uGamma(\mathcal{E}/\mathcal{C}) \xleftarrow{\ \sim\ } \Ind_\pi(\uGamma(\mathcal{E}(\pi)/\mathcal{C})) \ .
\]
\end{itemize}

% ---- from der-ch18_p151-155.tex
\printedpage{149}
% head: A. Grothendieck : Dérivateurs, Ch. XVIII, Édition des Universités de Montpellier II et Paris VII

\msp{page 393}
Toutes ces assertions se prouvent facilement dans l’ordre où elles sont données, en utilisant 4.14.17 et les résultats généraux déjà prouvés. Bien sûr, on peut commencer par remplacer $\mathcal{C}$ par une sous-catégorie réduite $\mathcal{C}_0$ équivalente et $\mathcal{E}$ par $\mathcal{E}_0 = \mathcal{E}|\mathcal{C}$ \texttt{[plutôt} $\mathcal{E}|\mathcal{C}_0$\texttt{]}, de façon à se réduire (dans b) et c)) au cas où $\operatorname{card}\mathcal{C} \leqslant \pi$.

\textbf{Corollaire 4.14.19.} \textit{Sous les conditions de 4.14.17, la sous-catégorie $\uGamma_{\cocart}(\mathcal{E}/\mathcal{C})$ de $\mathcal{M} = \uGamma(\mathcal{E}/\mathcal{C})$ (cas $p$ cofibrant), resp. la sous-catégorie $\uGamma_{\cart}(\mathcal{E}/\mathcal{C})$ (cas $p$ fibrant) est accessible dans $\mathcal{M}$, donc accessible et le foncteur d’inclusion} \texttt{[est]} \textit{accessible.}

(\textbf{N.B.} Ces deux catégories sont strictement pleines dans $\uGamma(\mathcal{E}/\mathcal{C})$, donc l’assertion qu’elles soient accessibles dans celle-ci \textit{équivaut} à la dernière assertion.)

\textsc{Démonstration de} 4.14.19. On peut supposer $\mathcal{C}$ réduite, donc petite. Considérons le cas où $p$ est cofibrant. Pour toute flèche $\lambda : s \longrightarrow t$ dans $\mathcal{C}$, considérons le foncteur
\[
\varphi_\lambda \ : \ \mathcal{M} \longrightarrow \uFl(\mathcal{E}_t)
\]
défini par
\msp{page 394}
\[
\varphi_\lambda(f) \ = \ (\lambda_*(f(s)) \xrightarrow{f(\lambda)^\natural} f(t)) \ .
\]
C’est un foncteur accessible, car son composé avec les foncteurs source et but sont les foncteurs $\lambda_*\varepsilon_s$ et $\varepsilon_t$, donc accessibles, $\varepsilon_s$, $\varepsilon_t$ et $\lambda_*$ étant accessibles. Considérons la sous-catégorie strictement pleine
\[
\uIs(\mathcal{E}_t) \ \subset \ \uFl(\mathcal{E}_t) \ ,
\]
on sait qu’elle est accessible dans $\uFl(\mathcal{E}_t)$ (4.14.12). Donc son image inverse par $\varphi_\lambda$ est accessible dans $\mathcal{M}$, et par suite aussi
\[
\uGamma_{\cart}(\mathcal{E}/\mathcal{C}) \ \overset{\text{déf}}{=} \ \bigcap_{\lambda\in\Ob\mathcal{C}} \varphi_\lambda^{-1}(\uIs(\mathcal{E}_t))
\]
% [sic] \uGamma_{\cart} -> \uGamma_{\cocart}, the case treated is p cofibrant
% [sic] \lambda\in\Ob\mathcal{C} -> \lambda\in\Fl\mathcal{C}, \lambda runs over arrows of \mathcal{C}
comme intersection d’une \textit{petite} famille de sous-catégories accessibles dans $\mathcal{M}$, q.e.d. Le cas où $p$ est fibrant se traite de façon similaire, en introduisant les foncteurs
\[
\psi_\lambda \ : \ \mathcal{M} = \uGamma(\mathcal{E}/\mathcal{C}) \longrightarrow \uFl(\mathcal{E}_s)
\]
donnés par
\[
\psi_\lambda(f) \ = \ (f(s) \xrightarrow{f(\lambda)^\flat} \lambda^* f(t)) \ .
\]

\msp{page 395}
\textbf{Corollaire 4.14.20.} \textit{Soit $\pi_0$ un cardinal infini choisi comme dans 4.14.18 a).}
\begin{itemize}[leftmargin=2.5em]
\item[\circled{a}] \textit{Alors $\mathcal{M}_{\cc} = \uGamma_{\cocart}(\mathcal{E}/\mathcal{C})$ (resp. $\mathcal{M}_{\cc} = \uGamma_{\cart}$) est $L_{\pi_0}$-stable dans $\mathcal{M} = \uGamma(\mathcal{E}/\mathcal{C})$, ou ce qui revient au même (comme elle est strictement pleine), elle satisfait $L_{\pi_0}$, et le foncteur d’inclusion satisfait à $L_{\pi_0}$. Les $\varinjlim$ grandes devant $\pi_0$ dans cette catégorie se calculent argument par argument.}
\item[\circled{b}] \textit{Pour tout cardinal $\pi \geqslant \pi_0$, tout objet de $\mathcal{M}_{\cc}$ qui est $\pi$-accessible dans $\mathcal{M}$ l’est (tautologiquement) dans $\mathcal{M}_{\cc}$. Si $\pi \geqslant \sup(\pi_0, \cardred\mathcal{C})$, alors (par 4.14.18 b) tout objet $f$ de $\mathcal{M}_{\cc}$ tel que $f(s) \in \Ob\mathcal{E}_{s\,\pi}$ $\forall\ s \in S$} \texttt{[plutôt} $s \in \Ob C$\texttt{]} \textit{est dans $\mathcal{M}_\pi$} \textup{et} \textit{dans $\mathcal{M}_{\cc\,\pi}$.}
% foot: Version 17 décembre 2025

\printedpage{150}
% head: A. Grothendieck : Dérivateurs, Ch. XVIII, Édition des Universités de Montpellier II et Paris VII

\item[\circled{c}] \textit{Supposons $\pi_0$ choisi comme dans 4.14.18 c, et que de plus $\mathcal{M}_{\cc}$ soit $\pi_0$-accessible dans $\mathcal{M}$. (Un tel $\pi_0$ existe.) On a donc, pour $\pi \geqslant \pi_1 = \sup(\pi_0, \cardred\mathcal{C}, \cardred\mathcal{M}_\pi)$ tel que $\pi^{\pi_0} = \pi$,}
\[
(\mathcal{M}_{\cc})_\pi \ = \ \mathcal{M}_\pi{\cap}\mathcal{M}_{\cc} \ = \quad
\begin{array}{l}
\textit{Sous-catégorie strictement pleine de $\mathcal{M}_{\cc}$ formée des $f$}\\
\textit{telles que $f(s) \in \Ob\mathcal{E}_{s\,\pi}$ $\forall\ s \in \Ob\mathcal{C}$.}
\end{array}
\]

\msp{page 396}
\textit{Si de plus $\pi$ est choisi tel qu’on ait $\pi \geqslant \sup_{s\in\Ob\mathcal{C}} \cardred\mathcal{E}_{s\,\pi_0}$, alors}
\[
\left\{
\begin{array}{rrcl}
 & \mathcal{M}_{\cc\,\pi} & \xleftarrow{\ \sim\ } & \uGamma_{\cocart}(\mathcal{E}(\pi)/\mathcal{C})\\[1ex]
\textit{resp.} & \mathcal{M}_{\cc\,\pi} & \xleftarrow{\ \sim\ } & \uGamma_{\cart}(\mathcal{E}(\pi)/\mathcal{C}) \ ,
\end{array}
\right.
\]
\textit{donc ($\pi$ étant utile pour $\mathcal{M}_{\cc}$) une équivalence de catégories}
\[
\begin{array}{rrcl}
 & \uGamma_{\cocart}(\mathcal{E}/\mathcal{C}) & \xleftarrow{\ \sim\ } & \mathrm{Ind}_\pi(\uGamma_{\cocart}(\mathcal{E}(\pi)/\mathcal{C}))\\
\textit{resp.}\quad & \uGamma_{\cart}(\mathcal{E}/\mathcal{C}) & \xleftarrow{\ \sim\ } & \mathrm{Ind}_\pi(\uGamma_{\cart}(\mathcal{E}(\pi)/\mathcal{C})).
\end{array}
\]
\end{itemize}

\textsc{Démonstration}. Tout est clair dans l’ordre où c’est donné en utilisant les deux corollaires précédents.

\textbf{Corollaire 4.14.21.} \textit{Considérons un diagramme de foncteurs}
\[
\begin{tikzcd}[column sep=small, row sep=small]
\mathcal{E}_1 \arrow[dr, "p_1"] & & \mathcal{E}_2 \arrow[dl, "p_2"'] \\
& \mathcal{E}_0 \ , &
\end{tikzcd}
\tag*{$(*)$}
\]
\textit{les catégories et les foncteurs étant accessibles.}

\msp{page 397}
\begin{itemize}[leftmargin=2.5em]
\item[\textit{a)}] \textit{Le produit 2-fibré}
\[
\mathcal{M}_0 \ = \ \mathcal{E}_1 \overset{2}{\times}_{\mathcal{E}_0} \mathcal{E}_2
\]
\textit{est accessible, et les foncteurs canoniques $\mathcal{M} \longrightarrow \mathcal{E}_1$, $\mathcal{M} \longrightarrow \mathcal{E}_2$ sont accessibles.} % [sic] \mathcal{M} \longrightarrow \mathcal{E}_1, \mathcal{M} \longrightarrow \mathcal{E}_2 -> \mathcal{M}_0 \longrightarrow \mathcal{E}_1, \mathcal{M}_0 \longrightarrow \mathcal{E}_2, assertion a) concerns \mathcal{M}_0
\item[\textit{b)}] (\footnote[169]{Ceci est contenu dans 4.14.22.1 ci-dessus (p. 406), appliqué à $\Pi = \mathcal{E}_1{\times}\mathcal{E}_2$, $\Pi' = \mathcal{E}_0$, et $F, G : \Pi \rightrightarrows \Pi'$ donnés par $F = p_1\circ\pr_1$ et $G = p_2\circ\pr_2$. D’autre part, a) et c) résultent facilement de b).}). \textit{La catégorie $\mathcal{M}$ des triples $(x_1,x_2,u)$ avec $x_1 \in \Ob\mathcal{E}_1$, $x_2 \in \Ob\mathcal{E}_2$, et}
\[
u \ : \ p_1(x_1) \ \longrightarrow \ p_2(x_2) \ ,
\]
\textit{est accessible, et les foncteurs canoniques $\mathcal{M} \longrightarrow \mathcal{E}_1$, $\mathcal{M} \longrightarrow \mathcal{E}_2$ sont accessibles.}
\item[\textit{c)}] \textit{La sous-catégorie strictement pleine $\mathcal{M}_0$ de $\mathcal{M}$ est accessible dans $\mathcal{M}$.}
\end{itemize}

\noindent\textsc{Démonstration}. Soit $\mathcal{C}$ la catégorie
\[
\begin{tikzcd}[column sep=small]
1 \arrow[dr] & & 2 \arrow[dl] \\
& 0 \ , &
\end{tikzcd}
\]
et $\mathcal{E}$ la catégorie cofibrée sur $\mathcal{C}$ définie par le diagramme $(*)$. On est sous les conditions de 4.14.17. La catégorie $\Gamma = \uGamma(\mathcal{E}/\mathcal{C})$ s’interprète
% foot: Version 17 décembre 2025
\printedpage{151}
% head: A. Grothendieck : Dérivateurs, Ch. XVIII, Édition des Universités de Montpellier II et Paris VII
\msp{page 398}
comme la catégorie des quintuplets $(x_0,x_1,x_2,u_1,u_2)$, avec $x_i \in \Ob\mathcal{E}$ pour $i \in \{0,1,2\}$, et $u_1, u_2 \in \Fl\mathcal{E}_0$,
\[
\begin{tikzcd}[column sep=small, row sep=small]
p_1(x_1) \arrow[dr, "u_1"] & & p_2(x_2) \arrow[dl, "u_2"'] \\
& x_0 \ . &
\end{tikzcd}
\tag*{$(**)$}
\]
En vertu de 4.14.18, cette catégorie $\uGamma$ est accessible. La catégorie $\mathcal{M}$ est équivalente à la sous-catégorie strictement pleine $\Gamma_0$ de $\Gamma$, formée des quintuplets pour lesquels on a $u_2 \in \uIs\mathcal{E}_0$. On a en effet des foncteurs quasi-inverses l’un de l’autre % [sic] \Gamma_0 -> \Gamma_1, the subcategory is called \Gamma_1 in what follows
\[
\begin{tikzcd}
\mathcal{M} \arrow[r, shift left, "\alpha"] & \Gamma_1 \arrow[l, shift left, "\beta"]
\end{tikzcd}
\]
donnés par
\[
\begin{aligned}
\alpha(x_1,x_2,u) \ &= \ (p_2(x_2),x_1,x_2,u,\id_{p_2(x_2)})\\
\beta(x_0,x_1,x_2,u_1,u_2) \ &= \ (x_1,x_2,u_2^{-1}u_1) \ ,
\end{aligned}
\]
on a alors
\[
\begin{aligned}
\beta\alpha \quad &= \quad \id_{\mathcal{M}}\\
\alpha\beta \quad &\xrightarrow{\ \sim\ } \quad \id_{\Gamma_1} \ .
\end{aligned}
\]
D’autre part, $\Gamma_1$ est évidemment accessible dans $\mathcal{M}$, comme image inverse de $\uIs(\mathcal{E}_0)$ (qui % [sic] accessible dans \mathcal{M} -> accessible dans \Gamma, \Gamma_1 is a subcategory of \Gamma
\msp{page 399}
est accessible dans $\uFl(\mathcal{E}_0)$ par le foncteur accessible
\[
\uGamma \ \longrightarrow \ \uFl(\mathcal{E}_0)
\]
donné par $(x_1, x_1, x_2, u_1, u_2) \longmapsto u_2$. Cela prouve que $\mathcal{M}$, qui est équivalente à $\Gamma_1$, est également accessible, et l’accessibilité du foncteur canonique $\mathcal{M} \longrightarrow \mathcal{E}_i$ ($i \in \{1,2\}$) résulte de celle de $\Gamma \longrightarrow \mathcal{E}_i$, impliquant celle du composé $\Gamma_1 \hookrightarrow \Gamma \longrightarrow \mathcal{E}_i$. Cela prouve donc b). Pour prouver a), il suffit de voir que $\mathcal{M}_0$ est accessible dans $\mathcal{M}$, or c’est encore l’image inverse de $\uIs(\mathcal{E}_0) \subset \uFl(\mathcal{E}_0)$ par le foncteur accessible canonique % [sic] (x_1, x_1, x_2, u_1, u_2) -> (x_0, x_1, x_2, u_1, u_2), the quintuplets begin with x_0
\[
\mathcal{M} \ \longrightarrow \ \uFl(\mathcal{E}_0) \ ,
\]
donc $\mathcal{M}$ est bien accessible dans $\mathcal{M}$, comme image inverse d’une sous-catégorie accessible dans $\uFl(\mathcal{E})$. On peut d’ailleurs aussi interpréter $\mathcal{M}$, à équivalence de catégories près, comme % [sic] donc \mathcal{M} est bien accessible -> donc \mathcal{M}_0 est bien accessible, \mathcal{M}_0 is the subcategory in question
% [sic] \uFl(\mathcal{E}) -> \uFl(\mathcal{E}_0), as in the preceding display
% [sic] interpréter \mathcal{M} -> interpréter \mathcal{M}_0, as in the following display
\msp{page 400}
\[
\mathcal{M}_0 \ \simeq \ \uGamma_{\cocart}(\mathcal{E}/\mathcal{C}) \ ,
\]
les foncteurs canoniques $\mathcal{M}_0 \longrightarrow \mathcal{E}_i$ ($i \in \{0,1\}$) s’interprétant comme des foncteurs évaluation, ce qui donne a) directement, comme cas particulier de 4.14.19. % [sic] i \in \{0,1\} -> i \in \{1,2\}, the canonical functors of a) go to \mathcal{E}_1, \mathcal{E}_2
% foot: Version 17 décembre 2025

\printedpage{152}
% head: A. Grothendieck : Dérivateurs, Ch. XVIII, Édition des Universités de Montpellier II et Paris VII

\textbf{4.14.22.} \textsc{Démonstration de} 4.14.17. Le cas où $p$ est cofibrant a déjà été démontré (th. 4.10.8, p. 263). Nous allons voir que le cas où $p$ est fibrant se ramène au théorème précédent (dans un cas très particulier), en utilisant à fond les résultats généraux sur les parties accessibles d’une catégorie accessible. Notons que les assertions a) et b) du corollaire 4.14.18 se prouvent de façon AQT \texttt{[âne qui trotte]} sans avoir à faire appel à 4.14.17. Le a) implique déjà que les foncteurs évaluation $\varepsilon_s : \mathcal{M} \longrightarrow \mathcal{E}_s$ sont accessibles. Il reste donc à prouver
\msp{page 401}
que $\mathcal{M} = \uGamma(\mathcal{E}/\mathcal{C})$ est accessible, dans le cas où $p$ est \textit{fibrant}. On supposera de plus que $\mathcal{C}$ est petite (quitte à remplacer $\mathcal{C}$ par une petite sous-catégorie pleine équivalente).

Soient
\[
\begin{aligned}
\Pi \ &= \ \prod\nolimits_{s\in\Ob\mathcal{C}} \mathcal{E}_s\\
\Pi' \ &= \ \prod\nolimits_{\lambda\in\Fl\mathcal{C}} \mathcal{E}_{s(\lambda)} \ ,
\end{aligned}
\]
où $s(\lambda)$ désigne la \textit{source} $s$ de la flèche $\lambda : s \longrightarrow t$ dans $\mathcal{C}$. $\Pi$, $\Pi'$ sont des catégories accessibles, comme (petits) produits de catégories accessibles. (Cas particulier trivial de 4.10.8, pour une catégorie d’indices discrète.) On va définir deux foncteurs
\[
\Pi \ \overset{F}{\underset{G}{\rightrightarrows}} \ \Pi'
\]
en définissant pour tout $\lambda : s_0 \longrightarrow t_0$ dans $\Fl\mathcal{C}$ deux foncteurs
\msp{page 402}
\[
\left\{
\begin{array}{lcll}
\Pi & \overset{F_\lambda}{\underset{G_\lambda}{\rightrightarrows}} & \mathcal{E}_{s(\lambda)=s_0} & \\
F_\lambda & = & \pr_{s_0} & \text{projection de } \Pi = \prod\nolimits_s \mathcal{E}_s \text{ sur le facteur } \mathcal{E}_{s_0}\\
G_\lambda & = & \lambda^*\pr_{s_1} & \text{composé } \Pi \xrightarrow{\pr_{s_1}} \mathcal{E}_{s_1} \xrightarrow{\lambda^*} \mathcal{E}_{s_0}\\
 & & & \text{\texttt{[plutôt} } \lambda^*\pr_{t_0} \text{ \texttt{resp.} } \Pi \xrightarrow{\pr_{t_0}} \mathcal{E}_{t_0} \xrightarrow{\lambda^*} \mathcal{E}_{s_0}\text{\texttt{]}}
\end{array}
\right.
\qquad \bigl(\mbox{\footnotemark[170]}, \mbox{\footnotemark[171]}\bigr) \ .
\]
\footnotetext[170]{Il est immédiat que les $F_\lambda$, $G_\lambda$, donc $F$ et $G$, sont des foncteurs \textit{accessibles}.}
\footnotetext[171]{\textbf{N.B.} On suppose choisi pour tout $\lambda \in \Fl\mathcal{C}$ un foncteur image inverse $\lambda^* : \mathcal{E}_t \longrightarrow \mathcal{E}_s$. On suppose que si $\lambda$ est une flèche identique, $\id_s$, \texttt{[on]} choisit $\lambda^* = \id_{\mathcal{E}_s}$.}

Soit $f \in \mathcal{M} = \uGamma(\mathcal{E}/\mathcal{C})$, alors pour tout $\lambda \in \Fl\mathcal{C}$, $f$ définit une flèche
\[
\alpha_\lambda(f) \ : \quad \underbrace{f(s)}_{= F_\lambda(\varepsilon(f))} \quad \xrightarrow{f(\lambda)^\flat} \quad \underbrace{\lambda^*(f(t))}_{= G_\lambda(\varepsilon(f))} \ ,
\]
\textit{fonctorielle} en $f$. On a un foncteur
\[
\varepsilon \ = \ (\varepsilon_s)_{s\in\Ob\mathcal{C}} \ : \ \mathcal{M} \ \longrightarrow \ \Pi \ ,
\]
et pour tout $\lambda \in \Fl\mathcal{C}$ un homomorphisme fonctoriel
\[
\alpha_\lambda \ : \ F_\lambda \circ \varepsilon \ \longrightarrow \ G_\lambda \circ \varepsilon \ ,
\]
d’où un homomorphisme fonctoriel
\[
\alpha \ = \ (\alpha_\lambda)_{\lambda\in\Fl\mathcal{C}} \ : \ F \circ \varepsilon \ \longrightarrow \ G \circ \varepsilon \ .
\]
Soit alors $\mathcal{M}'$ la catégorie suivante :
% foot: Version 17 décembre 2025
\printedpage{153}
% head: A. Grothendieck : Dérivateurs, Ch. XVIII, Édition des Universités de Montpellier II et Paris VII
\msp{page 403}
\[
\mathcal{M}' \quad = \quad \parbox[t]{0.5\textwidth}{catégorie des couples $(X,u)$, avec $X$ dans $\Ob\Pi$, et $u : F(X) \longrightarrow G(X)$ une flèche dans $\Pi'$.}
\]
L’homomorphisme fonctoriel $\alpha : F \circ \varepsilon \longrightarrow G \circ \varepsilon$ définit un foncteur
\[
\mathcal{M} \ \longrightarrow \ \mathcal{M}' \ ,
\]
et il est immédiat que ce foncteur est pleinement fidèle. En fait, la catégorie $\mathcal{M}'$ est isomorphe à la catégorie formée des couples $(f_0,f_1)$ de deux applications
\[
\begin{array}{lcccccc}
f_0 : & \Ob\mathcal{C} & \longrightarrow & \Ob\mathcal{E} & s & \longmapsto & f_0(s)\\[1ex]
f_1 : & \Fl\mathcal{C} & \longrightarrow & \Fl\mathcal{E} & \lambda & \longmapsto & f_1(\lambda)
\end{array}
\]
de telle façon que $f_1$ et $f_0$ soient compatibles avec les applications source et but, i.e. qu’ils définissent un homomorphisme de \textit{diagrammes} $f : \diag\mathcal{C} \longrightarrow \diag\mathcal{E}$, et que cet homomorphisme soit inverse à droite de $\mathrm{diag}(p) : \diag\mathcal{E} \longrightarrow \diag\mathcal{C}$.

\msp{page 404}
Le foncteur précédent $\mathcal{M} \hookrightarrow \mathcal{M}'$ s’interprète comme le foncteur qui associe à toute sections $f$ de $\mathcal{E}$ sur $\mathcal{C}$, le morphisme de diagrammes correspondant. C’est donc l’inclusion d’une sous-catégorie \textit{strictement pleine} de $\mathcal{M}'$. Donc pour prouver que $\mathcal{M}$ est accessible, il suffit de prouver : % [sic] toute sections -> toute section, agreement of toute
\begin{itemize}[leftmargin=2.5em]
\item[a)] $\mathcal{M}'$ est accessible.
\item[b)] $\mathcal{M}$ est accessible dans $\mathcal{M}'$.
\end{itemize}
Prouvons d’abord b) (en admettant a). La sous-catégorie $\mathcal{M}$ de $\mathcal{M}'$ est définie par deux types des conditions, sur un objet $f = (f_0,f_1)$ de $\mathcal{M}'$ :
\begin{itemize}[leftmargin=2.5em]
\item[1°)] $\forall\ s \in \Ob\mathcal{C}$, $f_1(\id_s) = \id_{f_0(s)}$.
\item[2°)] Si $s \xrightarrow{\lambda} t \xrightarrow{\mu} u$ sont deux flèches composables dans $\mathcal{C}$, alors
\[
f_1(\mu\lambda) \ = \ f_1(\mu)f_1(\lambda) \ .
\]
\end{itemize}
Ces conditions s’explicitent l’une et \texttt{[l’autre par]} l’égalité d’une double-flèche dans $\mathcal{E}_s$,
\msp{page 405}
savoir
\[
f_0(s) \overset{f_1(\id_s)}{\underset{\id_{f_0(s)}}{\rightrightarrows}} f_0(s)
\]
dans le cas 1°), et
\[
f_0(s) \overset{f_1(\mu\lambda)^\flat}{\underset{(f_1(\mu)f_1(\lambda))^\flat}{\rightrightarrows}} (\mu\lambda)^*(f_0(u))
\]
dans le cas 2°). Dans l’un et l’autre cas, on a un foncteur
\[
\mathcal{M}' \ \xrightarrow{\gamma_s \text{ resp. } \gamma_{\lambda,\mu}} \ \uHom(I,\mathcal{E}_s)
\]
% foot: Version 17 décembre 2025

% ---- from der-ch18_p156-160.tex
\printedpage{154}
% head: A. Grothendieck : Dérivateurs, Ch. XVIII, Édition des Universités de Montpellier II et Paris VII
dans la catégorie des double-flèches de $\mathcal{E}_s$ ($I$ désignant la catégorie double-flèche type $0\rightrightarrows 1$), foncteur qui est (trivialement) accessible. Soit alors $I\xrightarrow{\;i\;}\bDelta_1$ le foncteur canonique, et
\[
\underbrace{\uHom(\bDelta_1,\mathcal{E}_s)}_{\simeq\,\uFl(\mathcal{E}_s)}\ \hookrightarrow\ \uHom(I,\mathcal{E}_s)
\]
le foncteur $u\longmapsto u\circ i$, qui est un foncteur strictement pleine fidèle, ayant comme image la sous-catégorie des double-flèches $(u,v)$ telles que $u=v$. Comme le foncteur d'inclusion est accessible et $\uHom(\bDelta_1,\mathcal{E}_s)$ et $\uHom(I,\mathcal{E}_s)$ accessibles
\msp{page 406}
(cas particulier de 4.14.1, ou aussi de 4.10.8!), on voit que la sous-catégorie strictement pleine de $\uHom(I,\mathcal{E}_s)$ des double-flèches égalisées dans $\mathcal{E}_s$, est accessible dans $\uHom(I,\mathcal{E}_s)$. Donc, son image inverse dans $\mathcal{M}'$ est accessible. Donc l'intersection de ces images inverses pour $s\in\Ob\mathcal{C}$ et $(\lambda,\mu)\in\Fl_2(\mathcal{C})$ variables, l'est également ($\mathcal{C}$ étant petite).

Donc tout revient à prouver que $\mathcal{M}'$ est accessible, ce qui est contenu dans le

% vertical bar printed in the left margin alongside Lemme 4.14.22.1
\textbf{Lemme 4.14.22.1.} \textit{Soient $\Pi$, $\Pi'$ deux catégories accessibles, $F,G:\Pi\rightrightarrows\Pi'$ deux foncteurs accessibles, et $\mathcal{M}'$ la catégorie des couples $(X,u)$, avec $X\in\Ob\Pi$ et $u:F(x)\longrightarrow G(x)$ une flèche de $\Pi'$. Alors $\mathcal{M}'$ est accessible, et le foncteur canonique $\mathcal{M}'\longrightarrow\Pi$ est accessible.} % [sic] (X,u), X \in Ob Pi ... F(x) -> the same letter (X or x) throughout, capital X and lowercase x used for the same object

\textsc{Démonstration du Lemme.} La donnée de $\Pi$, $\Pi'$, $F$, $G$ définit une catégorie cofibrée $\mathcal{F}$ sur la catégorie $I=(0\rightrightarrows 1)$
\msp{page 407}
précédente. La catégorie $\uGamma(\mathcal{F}/I)$ s'interprète, à isomorphisme près, comme la catégorie des triples $(x,u,v)$ avec $x\in\Ob\Pi$, $x'\in\Ob\Pi'$, et $u$ et $v$ des flèches
\[
\begin{tikzcd}[column sep=small]
F(x) \arrow[dr, "u"'] & & G(x) \arrow[dl, "v"] \\
& x' &
\end{tikzcd}
\]
dans $\Pi'$. La catégorie $\mathcal{M}'$ est équivalente à la sous-catégorie strictement pleine de $\uGamma(\mathcal{F}/I)$, formée des triples $(x,u,v)$ tels que $v$ soit un isomorphisme, i.e. à la sous-catégorie de $\uGamma(\mathcal{F}/I)$, image inverse de $\uIs(\Pi')$ par le foncteur $(x,u,v)\longmapsto v$,
\[
\uGamma(\mathcal{F}/I)\ \longrightarrow\ \uFl(\Pi')\ .
\]
Or par 4.10.8, $\uGamma(\mathcal{F}/I)$ est accessible, et le foncteur précédent est bien sûr accessible. Enfin, par 4.14.12,
\msp{page 408}
$\uIs(\Pi')$ est accessible dans la catégorie accessible $\uFl(\Pi')$. Donc son image inverse $\mathcal{M}'$ dans $\uGamma(\mathcal{F}/I)$ est accessible dans $\uGamma(\mathcal{F}/I)$, donc accessible. Donc il en est de même de la catégorie équivalente, q.e.d.

Ainsi, grâce aux résultats généraux sur les parties accessibles des catégories accessibles, le théorème 4.14.17, dans le cas où $p$ est fibrant, est ramené aux cas cofibrant, dans le cas très particulier où la catégorie d'indices est $I=(0\rightrightarrows 1)$. On pourrait d'ailleurs faire une réduction similaire dans le cas où $p$ est cofibrant, par une démonstration pratiquement % [sic] aux cas cofibrant -> au cas cofibrant, singular noun
% foot: Version 17 décembre 2025
\printedpage{155}
% head: A. Grothendieck : Dérivateurs, Ch. XVIII, Édition des Universités de Montpellier II et Paris VII
identique. En d'autres termes, le théorème 4.10.8 se ramène au cas où la catégorie d'indices est $I=(0\rightrightarrows 1)$. Mais je suspecte que la démonstration du théorème dans ce cas particulier, n'est pas essentiellement plus simple que dans le cas général.

\msp{page 409}

\textbf{4.14.23.} On peut se demander de trouver un énoncé assurant l'accessibilité d'une catégorie $\uGamma(\mathcal{E}/\mathcal{C})$, qui coiffe les deux cas envisagés dans 4.14.17. Voici ce que je peux proposer. Je vais faire l'hypothèse suivante sur les foncteur $p:\mathcal{E}\longrightarrow\mathcal{C}$ (en plus de celle que $\mathcal{C}$ soit essentiellement petite) : % [sic] les foncteur -> le foncteur, singular functor p

\begin{itemize}[leftmargin=*]
\item[H)] Pour tout $\lambda\in\Fl\mathcal{C}$, la catégorie $\mathcal{E}_\lambda=\mathcal{E}\times_{\mathcal{C}}(\bDelta_1,i_\lambda)$, où $i_\lambda:\bDelta_1\longrightarrow\mathcal{C}$ est le foncteur défini par $\lambda$, est accessible, et le foncteur canonique $\mathcal{E}_\lambda\xrightarrow{\;p_\lambda\;}\bDelta_1$ est accessible.
\end{itemize}

Notons que si $\lambda:s\longrightarrow t$ dans $\mathcal{C}$ est tel que $\mathcal{E}_\lambda$ et $p_\lambda$ soient accessibles, alors $\mathcal{E}_s$ et $\mathcal{E}_t$ sont accessible, car isomorphes aux catégories image inverse par $p_\lambda:\mathcal{E}_\lambda\longrightarrow\bDelta_1$ des sous-catégories $\{0\}$ resp. $\{1\}$ de $\Delta_1$, lesquelles sont accessibles dans la catégorie accessible $\bDelta_1$, donc % [sic] sont accessible -> sont accessibles, plural agreement
\msp{page 410}
$\mathcal{E}_{\lambda_0}$ et $\mathcal{E}_{\lambda_1}$ sont accessibles dans $\mathcal{E}_\lambda$, donc sont accessibles, donc aussi $\mathcal{E}_s$ et $\mathcal{E}_t$. Ainsi, l'hypothèse H) implique que les catégories fibres sont accessibles. Notons de plus :

\textbf{Lemme 4.14.23.1.} \textit{Supposons que $p:\mathcal{E}\longrightarrow\mathcal{C}$ soit cofibrant ou fibrant, et soit $(\lambda:s\longrightarrow t)\in\Fl(\mathcal{C})$. Pour que $\mathcal{E}_\lambda$ soit accessible et $\mathcal{E}_\lambda\xrightarrow{\;p_\lambda\;}\bDelta_1$ accessible, il faut et il suffit que $\mathcal{E}_s$, $\mathcal{E}_t$ soient accessibles et $\lambda_*:\mathcal{E}_s\longrightarrow\mathcal{E}_t$ resp. $\lambda^*:\mathcal{E}_t\longrightarrow\mathcal{E}_s$ soit un foncteur accessible.}

\textsc{Démonstration.} On peut supposer que $\mathcal{C}=\bDelta_1$, et que $\mathcal{E}$ est la catégorie cofibrée (resp. fibrée) sur $\bDelta_1$ définie par un foncteur $\varphi:\mathcal{E}_0\longrightarrow\mathcal{E}_1$ (resp. $\psi:\mathcal{E}_1\longrightarrow\mathcal{E}_0$), et à montrer que $\mathcal{E}$ et $p:\mathcal{E}\longrightarrow\bDelta_1$ sont accessibles si et seulement si $\mathcal{E}_0$ et $\mathcal{E}_1$ et $\varphi$ (resp. $\psi$) le sont.

\msp{page 411}

Prouvons le “il faut”. On a déjà vu que si $\mathcal{E}$, $p$ sont accessibles, $\mathcal{E}_0$ et $\mathcal{E}_1$ sont accessibles dans $\mathcal{E}_s$, donc accessibles, et les foncteurs d'inclusion % [sic] dans E_s -> dans E, here C = Delta_1 and the subcategories lie in E
\[
\mathcal{E}_0\ \hookrightarrow\ \mathcal{E}\ \hookleftarrow\ \mathcal{E}_1
\]
accessibles. Dire que $p$ est cofibrant (resp. fibrant) signifie que le foncteur d'inclusion $\mathcal{E}_1\hookrightarrow\mathcal{E}$ admet un adjoint à droite $\bar{\varphi}$ (resp. que le foncteur $\mathcal{E}_0\hookrightarrow\mathcal{E}$ admet un adjoint à gauche $\tilde{\psi}$). Alors $\varphi$ (resp. $\psi$) est le composé $\mathcal{E}_0\longrightarrow\mathcal{E}\xrightarrow{\;\bar{\varphi}\;}\mathcal{E}_1$ (resp. le composé $\mathcal{E}_1\hookrightarrow\mathcal{E}\xrightarrow{\;\tilde{\psi}\;}\mathcal{E}_1$). Donc pour prouver que $\varphi$ (resp. $\psi$) est accessible, il suffit de prouver que $\bar{\varphi}$ (resp. $\tilde{\psi}$) l'est. Cela est un cas particulier du % [sic] E -> E_1 (target of psi-tilde composite) -> E_0, psi goes from E_1 to E_0

\textbf{Lemme 4.14.23.2.} \textit{Soient $\mathcal{E}$, $\mathcal{E}'$ deux catégories accessibles} (\footnote[172]{\textbf{N.B.} Il suffit que $\mathcal{E}$ admette une petite sous-catégorie génératrice, et que les objets de $\mathcal{E}'$ soient accessibles.}), \textit{et $\mathcal{E}\underset{G}{\overset{F}{\rightleftarrows}}\mathcal{E}'$ un couple de foncteurs adjoints entre $\mathcal{E}$ et $\mathcal{E}'$. Alors $F$ et $G$ sont accessibles.}

\msp{page 412}

En effet, $F$ est accessible, car il commute aux petites $\varinjlim$ quelconques, a fortiori il satisfait $\mathrm{L}_{\pi_0}$ (où $\pi_0$ est tel que $\mathcal{E}$, $\mathcal{E}'$ satisfassent $\mathrm{L}_{\pi_0}$). Prouvons que $G$ est accessible. Si $S$ est un
% foot: Version 17 décembre 2025
\printedpage{156}
% head: A. Grothendieck : Dérivateurs, Ch. XVIII, Édition des Universités de Montpellier II et Paris VII
petit ensemble d'éléments générateur dans $\mathcal{E}$, de sorte que la petite famille des foncteurs
\[
h_\xi\ =\ (x\longmapsto\Hom_{\mathcal{E}}(\xi,x))\ :\ \mathcal{E}\ \longrightarrow\ \Ens
\]
est conservative, pour prouver que $G$ est accessible, il suffit de prouver que les $h_\xi\circ G$ ($\xi\in S\subset\Ob\mathcal{E}$) le sont. Or
\[
L_\xi\circ G(x')\ =\ \Hom(\xi,G(x'))\ \simeq\ \Hom_{\mathcal{E}'}(F(\xi),x')\ , % [sic] L_xi -> h_xi, the functors were named h_xi
\]
et comme $F(\xi)$ est un objet accessible de $\mathcal{E}'$, le foncteur
\[
x'\ \longmapsto\ (L_\xi\circ G)(x')\ \simeq\ \Hom_{\mathcal{E}'}(F(\xi),x') % [sic] L_xi -> h_xi, the functors were named h_xi
\]
est bien accessible, q.e.d.

Cela prouve le “il faut” dans 4.14.23.1. Prouvons le “il suffit”, i.e. supposons $\mathcal{E}_s$, $\mathcal{E}_t$ accessibles, et $\varphi$ (resp. $\psi$) accessible.

\msp{page 413}

Dans les deux cas, $\mathcal{E}_0$ est un ouvert de $\mathcal{E}$, $\mathcal{E}_1$ est le fermé complémentaire, et $\mathcal{E}$ est la catégorie “recollée” à l'aide du bifoncteur
\[
H\ :\ \mathcal{E}_0^{\circ}\times\mathcal{E}_1\ \longrightarrow\ \Ens
\]
donné par
\[
\left\{
\begin{aligned}
H(x_0,x_1)\ &=\ \Hom(\varphi(x_0),x_1)\\
&\hspace{-2em}\text{resp.}\\
H(x_0,x_1)\ &=\ \Hom(x_0,\psi(x_1))\ ,
\end{aligned}
\right.
\]
suivant qu'on est dans le cas $p$ cofibrant ou $p$ fibrant. On va utiliser la

\textbf{Proposition 4.14.24} (\footnote[173]{cf. énoncés plus généraux et démonstration sympathique, 4.15.1 et 4.15.4 (p. 425, 432).}). \textit{Soit $\mathcal{E}$ une catégorie, $\mathcal{E}_0$ une sous-catégorie ouverte, $\mathcal{E}_1$ la sous-catégorie fermée complémentaire. Conditions équivalentes :}
\begin{itemize}[leftmargin=*]
\item[\textit{a)}] \textit{$\mathcal{E}$ est accessible, et $\mathcal{E}\longrightarrow\bDelta_1$ (de fibres $\mathcal{E}_0$, $\mathcal{E}_1$) est accessible.}
\item[\textit{b)}] \textit{$\mathcal{E}$ est accessible, et $\mathcal{E}_0$, $\mathcal{E}_1$ sont accessibles dans $\mathcal{E}$.}
\item[\textit{c)}] \textit{$\mathcal{E}_0$, $\mathcal{E}_1$ sont accessibles, et pour tout $x_0\in\Ob\mathcal{E}_0$, le foncteur}
\[
x_1\ \longmapsto\ \uHom_{\mathcal{E}}(x_0,x_1)\ :\ \mathcal{E}_1\ \longrightarrow\ \Ens
\]
\msp{page 414}
\textit{est accessible, et il existe $\pi$ tel que pour tout $x_1\in\Ob\mathcal{E}_1$, le foncteur}
\[
x_0\ \longmapsto\ \uHom_{\mathcal{E}}(x_0,x_1)^{\circ}\ :\ \mathcal{E}_0\ \longrightarrow\ (\Ens)^{\circ}
\]
\textit{soit $\pi$-accessible.}
\end{itemize}
% foot: Version 17 décembre 2025

\printedpage{157}
% head: A. Grothendieck : Dérivateurs, Ch. XVIII, Édition des Universités de Montpellier II et Paris VII
Admettons ce lemme \texttt{[plutôt cette proposition]}, achevons la démonstration de 4.14.23.1. On suppose $\mathcal{E}_0$, $\mathcal{E}_1$ et $\varphi$ (resp. $\psi$) accessibles, on veut prouver la condition a) de 4.14.23.3. Pour ceci, il suffit de prouver la condition c). Or on a, dans le cas $p$ cofibrant, % [sic] 4.14.23.3 -> 4.14.24, the conditions a), c) are those of Proposition 4.14.24
\[
\Hom_{\mathcal{E}}(x_0,x_1)\ \simeq\ \Hom_{\mathcal{E}_1}(\varphi(x_0),x_1)\ ,
\]
c'est bien accessible en $x_1$, puisque $\mathcal{E}_1$ est accessible, et $\pi$-accessible en $x_0$ si $\varphi$ l'est, comme composé de $\varphi$ avec $\mathcal{E}_1\longrightarrow\Ens$, $\Hom(-,x_1)$, qui l'est. Dans le cas $p$ fibrant, on a
\[
\Hom_{\mathcal{E}}(x_0,x_1)\ \simeq\ \Hom_{\mathcal{E}_0}(x_0,\psi(x_1))\ ,
\]
c'est accessible en $x_1$, comme le composé du foncteur accessible $\psi:\mathcal{E}_1\longrightarrow\mathcal{E}_0$, et du foncteur également accessible $\Hom(x_0,-):\mathcal{E}_0\longrightarrow\Ens$, qui est $\pi$-accessible en $x_0$ tautologiquement.

Pour achever la démonstration de 4.14.23.1, il reste donc à prouver la proposition 4.14.24. On a déjà noté qu'on a a $\Longrightarrow$ b, et b $\Longrightarrow$ c clair, puisque le foncteur envisagé $\mathcal{E}\longrightarrow\Ens$ est le composé $\mathcal{E}_1\overset{\text{incl.}}{\hookrightarrow}\mathcal{E}\xrightarrow{\;\Hom_{\mathcal{E}}(x_0,-)\;}\Ens$, est accessible comme composé de deux
\msp{page 415}
foncteurs accessibles. Reste donc à prouver c $\Longrightarrow$ a. Soit donc $\pi_0$ un cardinal utile pour $\mathcal{E}_0$, et tel que les foncteurs
\[
x_0\ \longmapsto\ H(x_0,x_1)\ :\ \mathcal{E}_0\longrightarrow\Ens^{\circ}
\]
soient $\pi_0$-accessibles, soit $\pi$ un cardinal $\geqslant\sup(\pi_0,\mathrm{card}\,\text{réd}\,\mathcal{E}_{0\,\pi_0})$ tel que $\pi^{\pi_0}=\pi$ (donc utile pour $\mathcal{E}_0$), utile pour $\mathcal{E}_1$, et tel que les foncteurs
\[
x_1\ \longmapsto\ H(x_0,x_1)\ :\ \mathcal{E}_1\longrightarrow\Ens
\]
soient $\pi$-accessibles, pour $x_0\in\Ob\mathcal{E}_{0\,\pi_0}$. Soit $\mathcal{E}(\pi)$ la sous-catégorie strictement pleine de $\mathcal{E}$, telle que
\[
\Ob\mathcal{E}(\pi)\ =\ \Ob\mathcal{E}_{0\,\pi}\cup\Ob\mathcal{E}_{1\,\pi}\ .
\]
Je dis que
\begin{itemize}[leftmargin=*]
\item[1°)] $\mathcal{E}$ satisfait à $\mathrm{L}_\pi$, et $\mathcal{E}_0$ et $\mathcal{E}_1$ sont stables dans $\mathcal{E}$ sous $\mathrm{L}_\pi$, les foncteurs $\mathcal{E}_0\longrightarrow\mathcal{E}$, $\mathcal{E}_1\longrightarrow\mathcal{E}$, $\mathcal{E}\longrightarrow\bDelta_1$ satisfont $\mathrm{L}_\pi$ (\footnote[174]{il suffit que $\mathcal{E}_0$ et $\mathcal{E}_1$ satisfassent $\mathrm{L}_\pi$, et que
\[
x_0\ \longmapsto\ H(x_0,x_1)\ :\ \mathcal{E}_0\longrightarrow\Ens^{\circ}
\]
satisfasse $\mathrm{L}_\pi$ $\forall\ x\in\Ob\mathcal{E}$.}).
\item[2°)] $\mathcal{E}(\pi)=\mathcal{E}_\pi$ (d'où un foncteur pleinement fidèle
\[
\mathrm{Ind}_\pi(\mathcal{E}(\pi))\ \longrightarrow\ \mathcal{E}\ \ ).
\]
\item[3°)] Ledit foncteur est essentiellement surjectif,
\msp{page 416}
donc une équivalence de catégories.
\end{itemize}
Notons que 2°) et 3°) impliquent que $\mathcal{E}$ est accessible, et 1°) que $\mathcal{E}\longrightarrow\bDelta_1$ est accessible, d'où la condition a) du lemme qu'il s'agit d'établir. Donc tout revient à prouver les 1°, 2°, 3° ci-dessus.
% foot: Version 17 décembre 2025

\printedpage{158}
% head: A. Grothendieck : Dérivateurs, Ch. XVIII, Édition des Universités de Montpellier II et Paris VII
1°) Comme $\mathcal{E}_1$ est fermée dans $\mathcal{E}$, il est clair que $\mathcal{E}$ est stable par toutes $\varinjlim$ représentables dans $\mathcal{E}$, et plus précisément, que pour un foncteur $I\longrightarrow\mathcal{E}_1$, sa limite inductive dans $\mathcal{E}_1$ existe si et seulement si elle existe dans $\mathcal{E}$, et les deux sont isomorphes. Supposons $I$ grande devant $\pi$, et soit $\varphi:I\longrightarrow\mathcal{E}$, $i\longmapsto x_i$. Distinguons deux cas.

a) $\varphi$ prend ses valeurs dans $\mathcal{E}_0$. Soit $x_0$ sa limite inductive \textit{dans} $\mathcal{E}_0$. Je dis que c'est aussi une limite inductive dans $\mathcal{E}$. Il faut prouver que pour $x_1\in\Ob\mathcal{E}_1$, la flèche
\[
\Hom_{\mathcal{E}}(x_0,x_1)\ \longrightarrow\ \varprojlim_{i\in I}\Hom_{\mathcal{E}}(x_i,x_1)
\]
est isomorphe. Or c'est le cas
\msp{page 417}
par hypothèse sur $\pi$,
\[
\xi\ \longmapsto H(\xi,x_1)\ :\ \mathcal{E}_0\ \longrightarrow\ (\Ens)^{\circ}
\]
satisfait à $\mathrm{L}_\pi$ pour tout $x_1$ dans $\Ob\mathcal{E}_1$.

b) $\varphi$ ne se factorise pas par $\mathcal{E}_0$, i.e. $\exists\ i_0\in\Ob I$ avec $\varphi(i_0)\in\Ob\mathcal{E}_1$. Alors on a $\varphi(i)\in\Ob\mathcal{E}_1$ pour tout $i\in\Ob I[i_0]$, où $I[i_0]$ désigne la sous-catégorie strictement pleine de $I$ formée des $i$ tels que $i\geqslant i_0$. Cette catégorie est cofinale dans $I$, donc grande devant $\pi$ comme $I$, et pour que $\varinjlim\limits_{I}x_i$ existe, il faut et il suffit que $\varinjlim\limits_{I[i_0]}x_i$ existe, ce qui est le cas d'après les cas a) déjà traité. % [sic] les cas a) déjà traité -> le cas a) déjà traité, singular agreement

Cela prouve donc en même temps que $\mathcal{E}$ satisfait à $\mathrm{L}_\pi$, et que $\mathcal{E}_0$ et $\mathcal{E}_1$ sont $\mathrm{L}_\pi$-stables. La démonstration montre en même temps que le foncteur $p:\mathcal{E}\longrightarrow\bDelta_1$ satisfait à $\mathrm{L}_\pi$.

2°) De ceci résulte aussitôt que $\mathcal{E}_\pi\subset\mathcal{E}(\pi)$. Prouvons qu'on a aussi
\[
\mathcal{E}(\pi)\ \subset\ \mathcal{E}_\pi\ ,
\]
et distinguons encore deux cas, pour
\msp{page 418}
$x\in\Ob E(\pi)$, suivant que $x$ est dans $\mathcal{E}_{0\,\pi}$ ou $\mathcal{E}_{1\,\pi}$. S'il est dans $\mathcal{E}_{1\,\pi}$, pour prouver que % [sic] E(pi) -> \mathcal{E}(pi), script letter elsewhere
\[
\Hom_{\mathcal{E}}(x,\underbrace{\varinjlim y_i}_{y})\ \xleftarrow{\ \sim\ }\ \varinjlim_{i}\Hom_{\mathcal{E}}(x,y_i)
\tag*{$(*)$}
\]
(pour un système inductif $(y_i)_{i\in I}$ dans $\mathcal{E}$ grand devant $\pi$), si les $y_i$ sont dans $\mathcal{E}_0$, donc $y\in\mathcal{E}_0$, alors les deux membres sont vides, et c'est OK. Sinon, on peut supposer (quitte à remplacer $I$ par un $I[i_0]$) que les $x_i$ sont dans $\mathcal{E}_1$, et alors c'est OK aussi, puisque $y$ est aussi la $\varinjlim$ dans $\mathcal{E}_1$, d'après 1°). % [sic] les x_i -> les y_i, the inductive system is (y_i)

Si $x$ est dans $\mathcal{E}_{0\,\pi}$, alors $(*)$ est clair si les $y_i$ sont dans $\mathcal{E}_0$. Sinon, on peut supposer encore que les $y_i$ sont dans $\mathcal{E}_0$, et $(*)$ s'écrit % [sic] y_i sont dans E_0 (second occurrence) -> E_1, the case "sinon" and the following "les y_i dans E_1"
\[
H(x,\varinjlim y_i)\ \xleftarrow{\ \sim\ }\ \varinjlim H(x,y_i)\ ,
\]
pour $x\in\Ob\mathcal{E}_{0\,\pi}$, les $y_i$ dans $\mathcal{E}_1$, $I$ grande devant $\pi$. En d'autres termes, il faut prouver que le foncteur
\[
x_1\ \longmapsto\ H(x_0,x_1)\ :\ \mathcal{E}_1\ \longrightarrow\ \Ens
\]
% foot: Version 17 décembre 2025

% ---- from der-ch18_p161-165.tex
\setcounter{footnote}{174}

%=====================================================================
\printedpage{159}
% head: A. Grothendieck : Dérivateurs, Ch. XVIII, Édition des Universités de Montpellier II et Paris VII

\msp{page 419}
est $\pi$-accessible, quand $x_0$ est dans $\cE_{0\,\pi}$. Or par hypothèse sur $\pi$, il est $\pi$-accessible si $x_0$ est dans $\cE_{0\,\pi_0}$. Mais si $x$ est dans $\cE_{0\,\pi}$, on aura % [sic] x -> x_0, the formula that follows concerns x_0
\[
x_0 \;=\; \varinjlim_{i\in I} x_{0\,i} \qquad
\left\{
\begin{array}{l}
I\ \ \text{grande}\ \ \text{devant}\ \ \pi_0,\\
\text{les}\ \ x_{0\,i}\ \ \text{dans}\ \ \cE_{0\,\pi_0},\\
\card I \leqslant \pi
\end{array}
\right.
\]
(car $\pi_0$ est utile pour $\cE_0$ par hypothèse, et $\pi$ satisfait $\pi^{\pi_0} = \pi \geqslant \cardred \cE_{\pi_0}$). Comme $x_0 \longrightarrow H(x_0,x_1) : \cE_0 \longrightarrow \Ens$ satisfait à $\mathrm{L}_{\pi_0}$, pour tout $x_1$ dans $\cE_1$, on aura donc % [sic] x_0 \longrightarrow H(x_0,x_1) -> x_0 \longmapsto H(x_0,x_1), elsewhere such functors are written with \longmapsto
\[
H(x_0) \;\simeq\; \varprojlim_{i\in I} H(x_{0\,i})
\]
(où pour $x_0$ dans $\cE_0$, $H(x)$ désigne le foncteur $x_1 \longmapsto H(x_0,x_1) : \cE_1 \longrightarrow (\Ens)$). Comme % [sic] H(x) -> H(x_0), the functor depends on x_0
les foncteurs $H(x_{0\,i})$ sont $\pi$-accessibles, et $\card I \leqslant \pi$, il s'ensuit que leur $\varprojlim\limits_{I}$ $H(x_0)$ est $\pi$-accessible (SGA 4 I 9.6 (ii)). Cela achève de prouver donc 2°).

3°). Pour prouver que $\Ind_\pi(\cE(\pi)) \longrightarrow \cE$ est essentiellement surjectif, il suffit de voir que tout $x \in \Ob\cE_0$ et tout $x \in \Ob\cE_1$
\msp{page 420}
est dans l'image essentielle. Or cela résulte de l'hypothèse que $\pi$ est utile pour $\cE_0$ et $\cE_1$, et de 1°) qui implique que $\cE_0 \longrightarrow \cE$ et $\cE_1 \longrightarrow \cE$ satisfont $\mathrm{L}_\pi$.

Cela achève de prouver la proposition 4.14.24, et de même le résultat suivant, plus précis que l'implication c) $\Longrightarrow$ a).

\textbf{Corollaire 4.14.25.} \textit{Soient $\cE$ une catégorie, $\cE_0$ une sous-catégorie ouverte, $\cE_1$ la sous-catégorie fermée complémentaire,}
\[
H \;=\; ((x_0,x_1) \longmapsto \Hom_{\cE}(x_0,x_1)) \;:\; \cE_0^\circ \times \cE_1 \;\longrightarrow\; \Ens
\]
\textit{le foncteur recollement.}
\begin{enumerate}
\item[\textit{1°)}] \textit{Soit $\pi$ un cardinal infini. Pour ce que $\cE$ et $\cE_0$ satisfassent à $L_\pi$, et le foncteur d'inclusion $\cE_0 \hookrightarrow \cE$ itou, il faut et il suffit que $\cE_0$ satisfasse à $L_\pi$, et que pour tout $x_1$ dans $\cE_1$, le foncteur $x_0 \longmapsto H(x_0,x_1) : \cE_0 \longrightarrow \Ens^\circ$ satisfasse à $L_\pi$. Alors $\cE_1$ et $\cE_1 \hookrightarrow \cE$ satisfont aussi à $L_\pi$,}
\msp{page 421}
\textit{et on a}
\[
\cE_\pi \;\subset\; \cE(\pi) \;\defeq\;
\begin{array}{l}
\text{\itshape sous-catégorie\ \ strictement\ \ pleine\ \ de\ \ $\cE$\ \ telle\ \ que}\\
\Ob\cE(\pi) = \Ob\cE_{0\,\pi} \cup \Ob\cE_{1\,\pi}.
\end{array}
\]
\item[\textit{2°)}] \textit{Supposons que la condition c) de 4.14.24 soit satisfaite, et soit $\pi_0$ un cardinal infini tel que $\pi_0$ soit utile pour $\cE_0$, et que les foncteurs $x_0 \longmapsto H(x_0,x_1) : \cE_0 \longrightarrow \Ens^\circ$ définis par les $x_1 \in \Ob\cE_1$ soient $\pi_0$-accessibles. Soit $\pi$ un cardinal tel que $\pi^{\pi_0} = \pi \geqslant \sup(\pi_0, \cardred\cE_{0\,\pi_0})$, et tel que pour tout $x_0 \in \Ob\cE_{\pi_0}$, le foncteur $x_1 \longmapsto H(x_0,x_1) : \cE_1 \longrightarrow \Ens$ soit $\pi$-accessible. Alors on a}
\[
\cE_\pi \;=\; \cE(\pi)\ .
\]
% foot: Version 17 décembre 2025

%=====================================================================
\printedpage{160}
% head: A. Grothendieck : Dérivateurs, Ch. XVIII, Édition des Universités de Montpellier II et Paris VII

\item[\textit{3°)}] \textit{Si $\pi_0$ et $\pi$ sont comme dans 2°), et si de plus $\pi$ est utile pour $\cE_1$, alors $\pi$ est utile pour $\cE$, donc (en vertu de 2°)}
\[
\cE \;\xleftarrow{\ \sim\ }\; \Ind_\pi(\cE(\pi))\ .
\]
\end{enumerate}
Nous pouvons maintenant énoncer la généralisation cherchée de 4.14.17 :

\msp{page 422}
\textbf{Théorème 4.14.26} (généralisation de 4.14.13). \textit{Soit $p : \cE \longrightarrow \cC$ un foncteur, avec $\cC$ petite. On suppose que $p$ satisfait à l'hypothèse H) de la page 409, qui équivaut (moyennant 4.14.24) à la suivante :}
\begin{enumerate}
\item[\textit{H$'$)}] \textit{$\forall\ s \in \Ob\cC$, la catégorie fibre $\cE_s$ est accessible, et pour tout $\lambda : s \longrightarrow t$ dans $\Fl\cC$, le foncteur}
\[
\left\{
\begin{array}{lcl}
H_\lambda : \cE_s^\circ \times \cE_t & \longrightarrow & \Ens\\[1ex]
\qquad (x_0,x_1) & \longmapsto & \Hom_\lambda(x_0,x_1)
\end{array}
\right.
\]
\textit{satisfait les deux conditions suivantes :}
\begin{enumerate}
\item[\textit{a)}] \textit{$\forall\ x_0 \in \Ob\cE_s^\circ$, $x_1 \longmapsto H_\lambda(x_0,x_1) : \cE_t \longrightarrow \Ens$ est accessible.}
\item[\textit{b)}] \textit{Il existe un cardinal infini $\pi_\lambda$, tel que pour tout $x_1 \in \Ob\cE_t$, le foncteur $x_0 \longmapsto H_\lambda(x_0,x_1) : \cE_s \longrightarrow (\Ens)^\circ$ soit $\pi_\lambda$-accessible.}
\end{enumerate}
\end{enumerate}
\textit{Ceci posé, la catégorie $\underline{\Gamma}(\cE/\cC)$ des sections de $\cE$ sur $\cC$ est accessible, et les foncteurs $\varepsilon_s : \underline{\Gamma}(\cE/\cC) \longrightarrow \cE_s$ sont accessibles. On peut préciser de façon similaire à 4.14.18 (énoncé en forme laissé au lecteur).}

\msp{page 423}
\textsc{Démonstration} de 4.14.26. Elle se fait exactement comme celle de 4.14.17 (p. 400 ff.), mais est conceptuellement plus simple. On suppose $\cC$ petite, et on introduit les catégories
\[
\Pi \;=\; \prod_{s\,\in\,\Ob\cC} \cE_s\ , \qquad \Pi' \;=\; \prod_{\lambda\,\in\,\Fl\cC} \cE_\lambda\ ,
\]
où
\[
\cE_\lambda \;=\; \cE \times_{\cC} (\boldsymbol{\Delta}_1, i_\lambda)\ ,
\]
et les deux foncteurs canoniques
\[
\Pi \;\underset{G}{\overset{F}{\rightrightarrows}}\; \Pi'
\]
de composantes
\[
\begin{aligned}
F_\lambda \;&=\; i_\lambda \circ \pr_{s(\lambda)}\\
G_\lambda \;&=\; j_\lambda \circ \pr_{b(\lambda)}\ ,
\end{aligned}
\]
où pour $\lambda : s \longrightarrow t$, on pose
\[
s(\lambda) \;=\; s\ , \qquad b(\lambda) \;=\; t\ ,
\]
et $i_\lambda$, $j_\lambda$ sont les inclusions canoniques
% In the print, the vertical arrow p_lambda and Delta_1 are enclosed together in a pair of large parentheses.
\[
\begin{tikzcd}
\cE_s \arrow[r,"i_\lambda"] \arrow[d,dash,"\wr"] & \cE_\lambda \arrow[dd,"p_\lambda"] & \cE_t\ . \arrow[l,"j_\lambda"'] \arrow[d,dash,"\wr"] \\
(\cE_\lambda)_0 & & (\cE_\lambda)_1 \\
& \boldsymbol{\Delta}_1 &
\end{tikzcd}
\]
% foot: Version 17 décembre 2025

%=====================================================================
\printedpage{161}
% head: A. Grothendieck : Dérivateurs, Ch. XVIII, Édition des Universités de Montpellier II et Paris VII

\msp{page 424}
On désigne par $\cM'$ la catégorie
\[
\cM' \;=\;
\begin{array}{l}
\text{catégorie des couples } (X,u)\text{, avec}\\
X \in \Ob\Pi,\ u : F(x) \longrightarrow G(x),
\end{array}
\] % [sic] F(x), G(x) -> F(X), G(X), the object is named X
qui est constructible par 4.14.22.1 est accessible (car $\Pi$ et $\pi'$ et les foncteurs $F$ et $G$ le sont, par l'hypothèse H) faite sur $p$). On peut l'interpréter comme la catégorie des sections % [sic] \pi' -> \Pi', the category introduced is \Pi'
de $\mathrm{Diag}(\cE)$ sur $\mathrm{Diag}(\cC)$. On a un foncteur strictement pleinement fidèle
\[
\cM \;\defeq\; \underline{\Gamma}(\cE/\cC) \;\hookrightarrow\; \cM'\ ,
\]
qui permet d'interpréter $\cM$ comme une sous-catégorie strictement pleine de $\cM'$. On voit alors comme dans le cas de 4.14.17 que cette sous-catégorie est accessible dans $\cM'$ -- c'est essentiellement sorital. Il s'ensuit que $\cM$ est accessible et que les foncteurs $\cM \longrightarrow \cE_s$ le sont aussi, comme composés de
\[
\cM \;\overset{\text{\upshape acc.}}{\hookrightarrow}\; \cM' \;\overset{\text{\upshape acc.}}{\hookrightarrow}\; \Pi \;\xrightarrow{\ \ \pr_s\ \ \text{\upshape acc.}\ \ }\; \cE_s\ ,
\]

\hfill q.e.d.

\msp{page 425}

\subsection*{4.15\quad Catégorie accessible sur une autre}

\textbf{Proposition 4.15.1.} \textit{Soit $f : \cE \longrightarrow \cC$ un foncteur avec $\cC$ petite. On suppose de plus satisfaite la condition H (p. 409), qu'on explicite ici sous la forme équivalente H$'$ :}
\begin{enumerate}
\item[\textit{H 1)}] \textit{Pour tout $s \in \Ob\cC$, $\cE_s$ est accessible.}
\item[\textit{H 2)}] \textit{Pour toute $\lambda \in \Fl\cC$, $\lambda : s \longrightarrow t$, existe un cardinal infini $\pi_\lambda$ tel que pour tout $x_1 \in \Ob\cE_t$, le foncteur}
\[
x_0 \;\longmapsto\; \Hom_\lambda(x_0,x_1) \;:\; \cE_s \;\longrightarrow\; (\Ens)^\circ
\]
\textit{est $\pi_\lambda$-accessible ($\cE_s$ satisfaisant de plus à $L_{\pi_\lambda}$).}
\item[\textit{H 3)}] \textit{Pour toute $\lambda \in \Fl\cC$, $\lambda : s \longrightarrow t$, et pour tout $x_0 \in \Ob\cE_s$, le foncteur}
\[
x_1 \;\longmapsto\; \Hom_\lambda(x_0,x_1) \;:\; \cE_t \;\longrightarrow\; (\Ens)
\]
\textit{est accessible.}
\end{enumerate}
\textit{Alors les foncteurs d'inclusion $i_s : \cE_s \longrightarrow \cE$ sont accessibles. De plus, si $\cC$ ``est sans projecteurs'' (i.e. si tout endomorphisme $p$ d'un objet $s$ de $\cC$ tel que $p^2 = p$, est l'identité), alors $\cE$ est accessible, et $f : \cE \longrightarrow \cC$ est accessible.}

Cela résulte du résultat plus précis :
% foot: Version 17 décembre 2025

%=====================================================================
\printedpage{162}
% head: A. Grothendieck : Dérivateurs, Ch. XVIII, Édition des Universités de Montpellier II et Paris VII

\msp{page 426}
\textbf{Corollaire 4.15.2} (\footnote{Pour a) et b) on n'utilise que H 1 et H 2, à l'exclusion de H 3.}). \textit{Les hypothèses sur $f : \cE \longrightarrow \cC$ sont celles de 4.15.1 ($\cC$ petite, et $f$ satisfait $H' =$ H 1) 2) 3)).}
\begin{enumerate}
\item[\textit{a)}] \textit{Soit $s \in \Ob\cC$ et soit $\pi_1$ un cardinal infini tel que pour toute $\lambda : s \longrightarrow t$ (flèche de $\cC$ de source $t$), on ait $\pi_1 \geqslant \pi_\lambda$ (où $\pi_\lambda$ est défini dans l'énoncé de H 2). Alors le foncteur $i_s : \cE_s \longrightarrow \cE$ est $\pi_1$-accessible.} % [sic] source t -> source s, lambda : s -> t has source s (cf. p. 163 "de source s")
\item[\textit{b)}] \textit{On suppose $\cC$ sans projecteurs. Soit $\pi_1$ un cardinal $\geqslant \sup(\pi_\lambda, \card\cC)$. Alors $\cE$ satisfait $L_{\pi_1}$, et $f : \cE \longrightarrow \cC$ est $\pi$-accessible.}
\item[\textit{c)}] \textit{Supposons encore que $\cC$ soit sans projecteurs (p.ex. une catégorie ``rigide'', p.ex. une catégorie préordonnée). Soient $\pi_0$, $\pi_1$ des cardinaux tels que :}
\begin{enumerate}
\item[\textcircled{\footnotesize $1^o$}] $\forall\ s \in \Ob\cC$, $\pi_0$ est utile pour $\cE_s$ \textit{(un tel $\pi_0$ existe grâce à H 1 et au fait que $\cC$, donc $\Ob\cC$ petits).}
\item[\textcircled{\footnotesize $2^o$}] $\pi_1 \geqslant \sup(\pi_0, \card\cC, \sup_{\lambda\in\Fl\cC}\pi_\lambda, \sup_{s\in\Ob\cC} \cardred \cE_{s\,\pi_0})$.
\item[\textcircled{\footnotesize $3^o$}] \textit{$\forall\lambda \in \Fl\cC$, et $x_0 \in \Ob\cE_{s,\pi_0}$, le foncteur $H_\lambda(x_0)$ de H 3, $x_1 \longmapsto \Hom_\lambda(x_0,x_1) : \cE_t \longrightarrow \Ens$ est $\pi_1$-accessible.}
\end{enumerate}
\textit{Soit $\pi$ un cardinal tel que $\pi^{\pi_0} = \pi \geqslant \pi_1$. Alors par b), $\cE$ satisfait $L_{\pi_1}$, donc $L_\pi$, et $f : \cE \longrightarrow \cC$ est $\pi_1$-accessible. De plus,}
\[
i_s(\cE_{s\,\pi}) \;\subset\; \cE_\pi\ ,
\]
\textit{ou encore}
\[
\cE(\pi) \;\subset\; \cE_\pi\ , \tag*{$(*)$}
\]
\textit{où $\cE(\pi)$ est la sous-catégorie pleine de $\cE$ telle que}
\[
\Ob\cE(\pi) \;=\; \bigcup_{s\in\Ob\cC} \Ob\,(\cE_{s\,\pi})\ .
\]
\textit{Comme $\cE_\pi$ satisfait $L_{\pi_1}$ par b), donc a fortiori $L_\pi$, on a donc par $(*)$ un foncteur} pleinement fidèle
\[
\Ind_\pi(\cE_{s\,\pi}) \;\xrightarrow{\ \sim\ }\; \cE\ ,
\]
\textit{et ce foncteur est une équivalence de catégories (donc $\cE$ est accessible).}
\item[\textit{d)}] \textit{En fait, sous les hypothèses de c) sur $\cC$, $\pi_0$, $\pi_1$, $\pi$, on a $\cE(\pi) = \cE_\pi$ (non seulement $\cE(\pi) \subset \cE_\pi$).}
\end{enumerate}

\msp{page 427}
\textsc{Démonstration de 4.15.1 et 4.15.2.} Il est clair que 4.15.1 est contenu dans 4.15.2 a) et c). Nous allons donc prouver 4.15.2.

a) Soit $(x_i)_{i\in I}$ un système inductif dans $\cE_s$, avec $I$ grand devant $\pi_0$, et $x$ sa limite inductive dans $\cE_s$ (qui existe, car $\cE_s$ satisfait à $\mathrm{L}_{\pi_\lambda}$, donc à $\mathrm{L}_{\pi_0}$). Il faut prouver que c'est une limite dans $\cE$, donc que pour tout $y \in \Ob\cE$, la flèche canonique
\[
\Hom_{\cE}(x,y) \;\longrightarrow\; \varprojlim_{I} \Hom_{\cE}(x_i,y) \tag*{$(*)$}
\]
% foot: Version 17 décembre 2025

%=====================================================================
\printedpage{163}
% head: A. Grothendieck : Dérivateurs, Ch. XVIII, Édition des Universités de Montpellier II et Paris VII
est un isomorphisme. Or $\Hom_{\cE}(x,y)$ et les $\Hom_{\cE}(x_i,y)$ sont des objets de $\Ens/\Hom_{\cC}(s,t)$, où $s = f(x)$, $t = f(y)$, et les flèches de transition sont des flèches dans $\Ens/\Hom_{\cC}(s,t)$. Cela montre que la bijectivité de $(*)$ revient à celle de
\[
\Hom_\lambda(x,y) \;\longrightarrow\; \varprojlim_{I} \Hom_\lambda(x_i,y) \tag*{$(*)$}
\]
pour tout $\lambda \in \Hom_{\cC}(x,y)$ \texttt{[plutôt} $\Hom_{\cC}(f(x),f(y))$\texttt{]}. Pour $\lambda$ fixé, et $y$ fixé dans $\Ob\cE_t$, la bijectivité dans $(*)$, pour tout système inductif grand devant $\pi_1$ dans $\cE_s$, de limite $x$, signifie que le foncteur
\[
x \;\longmapsto\; \Hom_\lambda(x,y) \;:\; \cE_s \;\longrightarrow\; (\Ens)^\circ
\]
est $\pi_0$-accessible. Or c'est bien le cas, par hypothèse sur $\pi_1$, quel que soit la flèche $\lambda : s \longrightarrow t$ de source $s$, et quel que soit $y \in \Ob\cE_t$.

\msp{page 428}
Soit $(x_i)_{i\in I}$ un système inductif grand devant $\pi_1$, et $(s_i)_{i\in I}$ l'image de ce système dans $\cC$. Nous prouverons plus bas le

% In the left margin, beside the statement of Lemme 4.15.3 and beside the following paragraph "Mais attention ...", the edition prints a wavy marginal sign (one for each).
\textbf{Lemme 4.15.3} (\footnote{Cf. démonstration page 452.},\footnote{sous réserve de la validité de 4.8.13 (Mais c'est à présent établi, sous une forme suffisante, dans 4.15.10).}). \textit{Soit $\cC$ une catégorie, $\pi$ un cardinal, tels que} $\card\cC \leqslant \pi$\textit{, et que $\cC$ soit sans projecteurs. Alors pour tout système inductif $(s_i)_{i\in I}$ grand devant $\pi$ dans $\cC$, existe une partie cofinale $J$ dans $I$ telle que le système induit $(s_i)_{i\in J}$ soit} constant, \textit{i.e. tel que les $s_i$ soient égaux à un $s \in \cC$, et que les morphismes de transition soient tous égaux à} $\id_s$.

Mais attention, on s'appuie pour prouver 4.15.3 sur le théorème 4.8.13 (page 212-216), que je n'ai toujours pas prouvé. Dans le cas où $\cC$ est rigide (i.e. $\End_{\cC}(s) = \{\id_s\}\ \forall\ s \in \Ob\cC$), la démonstration est cependant immédiate, cf. plus bas.

\msp{page 429}
Pour prouver que $x = \varinjlim\limits_{I} x_i$ existe dans $\cE$, et commute à $f$, on peut remplacer $I$ par une sous-catégorie cofinale. Cela nous ramène, grâce au lemme, au cas où $(s_i)_{i\in I}$ est constant, donc au cas où le système inductif provient d'un système inductif dans une catégorie fibre $\cE_s$. Mais par a), on sait que sa limite $x$ dans $\cE_s$ est une limite dans $\cE$, d'où l'existence de cette dernière, et la commutation à $f$, puisque $\varinjlim s_i = s$ ( $\varinjlim$ d'un système inductif filtrant constant).

c) Il faut prouver que si $x \in \Ob\cE_s$ est $\pi$-accessible dans $\cE_s$, il l'est dans $\cE$, i.e. que pour tout système inductif $(y_i)_{i\in I}$ grande devant $\pi$, la flèche canonique
\[
\varinjlim_{i} \Hom(x,y_i) \;\longrightarrow\; \Hom(x,y) \tag*{$(**)$}
\]
est bijective. Or comme par 2°) $\pi \geqslant \pi_1 \geqslant \sup(\card S, \sup_\lambda \pi_\lambda)$, il résulte encore du lemme % [sic] card S -> card \mathcal{C}, no S is defined; the condition 2° bounds card \mathcal{C}
que, quitte à remplacer $I$ par une sous-catégorie cofinale, on peut supposer que le système inductif $(y_i)_{i\in I}$ provient d'un système dans une catégorie fibre $\cE_t$, et que $y$ est sa $\varinjlim$ dans $\cE_t$. Mais alors les ensembles $\Hom(x,y)$ et $\Hom(x,y_i)$, et les morphismes de transition entre celles-ci, sont encore dans $\Ens/\Hom(s,t)$, et on est ramené
% foot: Version 17 décembre 2025
\end{document}
